% concatenated from INAR_rough_Heston.tex
\documentclass[letterpaper]{article}
\usepackage[margin=1in]{geometry}
\usepackage{float}
\usepackage{multirow} 
\usepackage{bm}
\usepackage{amsmath,amsthm,amsfonts,amssymb}
\usepackage{enumerate,xcolor}
\usepackage{graphicx}
\usepackage{subcaption}
\usepackage[authoryear,round]{natbib}
\usepackage[ruled,vlined]{algorithm2e}
\usepackage{booktabs,tabularx}

\PassOptionsToPackage{hyphens}{url}
\usepackage{url} 
\Urlmuskip=0mu plus 1mu 

\usepackage{hyperref}
\urlstyle{same}
\hypersetup{
  pdftitle  = {Rough Heston model as the scaling limit of bivariate cumulative heavy-tailed INAR processes: Weak-error bounds and option pricing},
  pdfauthor = {Yingli Wang and Zhenyu Cui and Lingjiong Zhu},
  colorlinks = true,
  linkcolor  = [rgb]{0.10,0.20,0.55},
  citecolor  = [rgb]{0.10,0.20,0.55},
  urlcolor   = [rgb]{0.10,0.20,0.55}
}

\usepackage{booktabs}
\usepackage{array}

\allowdisplaybreaks

\setlength{\parskip}{5pt}



\numberwithin{equation}{section}
\newtheorem{theorem}{Theorem}[section]
\newtheorem{lemma}[theorem]{Lemma}
\newtheorem{proposition}[theorem]{Proposition}
\newtheorem{corollary}[theorem]{Corollary}
\newtheorem{conjecture}[theorem]{Conjecture}
\newtheorem{assumption}[theorem]{Assumption}

\theoremstyle{definition}
\newtheorem{definition}[theorem]{Definition}
\newtheorem{example}[theorem]{Example}
\newtheorem{notation}[theorem]{Notation}
\newtheorem{question}[theorem]{Question}
\newtheorem{claim}{Claim}

\theoremstyle{remark}
\newtheorem{remark}[theorem]{Remark}
\newcommand{\red}[1]{\textcolor{red}{#1}}
\newcommand{\add}[1]{\textcolor{blue}{#1}}
\newcommand{\Real}{\mathbb{R}}
\newcommand{\set}[1]{\{#1\}}
\newcommand{\abs}[1]{\left\vert#1\right\vert}
\newcommand{\EPT}[1]{\mathrm{E}[#1]}
\newcommand{\me}{\mathrm{e}}
\newcommand{\mi}{\mathrm{i}}
\newcommand{\R}{\mathbb{R}}
\newcommand{\E}{\mathbb{E}}
\newcommand{\C}{\mathbb{C}}
\newcommand{\Q}{\mathbb{Q}}
\newcommand{\F}{\mathcal{F}}
\newcommand{\ii}{\mathrm{i}}
\newcommand{\Var}{\mathrm{Var}}
\newcommand{\cadlag}{c{\`a}dl{\`a}g}
\newcommand{\levy}{L{\'e}vy}
\newcommand{\poincare}{Poincar{\'e}}
\newcommand{\prob}{\mathbb{P}}
\newcommand{\expectation}{\mathbb{E}}
\newcommand{\norm}[1]{\left\lVert#1\right\rVert}
\newcommand{\loc}{\text{loc}}
\newcommand{\KerE}[1]{\mathsf{G}\left(#1\right)}
\newcommand{\KerEt}[1]{\mathsf{G}^\tau\left(#1\right)}
%\titlespacing{\section}{0pt}{0.2ex}{0ex}
%\titlespacing{\subsection}{0pt}{0.3ex}{0.3ex}
%\linespread{0.96}





\begin{document}


\title{Rough Heston model as the scaling limit of bivariate cumulative heavy-tailed INAR processes: Weak-error bounds and option pricing
}

\author{%
Yingli Wang\thanks{School of Mathematics, Shanghai University of Finance and Economics, Shanghai, People's Republic of China; naturesky1994@gmail.com.}%
\and
Zhenyu Cui\thanks{School of Business, Stevens Institute of Technology, Hoboken, NJ 07030, United States of America; zcui6@stevens.edu.}
\and
Lingjiong Zhu\thanks{Corresponding author.  Department of Mathematics, Florida State University, Tallahassee, Florida, United States of America; zhu@math.fsu.edu. All code (including FFT acceleration and option pricers) is available at \protect\url{https://github.com/gagawjbytw/INAR-rough-Heston}.}%}
}

\date{\today}

\maketitle

%\tableofcontents

\begin{abstract}
We study nearly unstable bivariate cumulative heavy-tailed INAR($\infty$) processes and show that, under a one-factor parameterization and a suitable scaling, they converge to the rough Heston model. This yields a discrete-time microstructural route to the joint price-variance dynamics and gives explicit formulas linking the INAR asymmetry parameters to the leverage correlation and diffusion scale of the limiting volatility process.
On the pricing side, we derive the exact finite-$\tau$ transform recursion and reduce it, in the diffusive scaling regime, to a quadratic discrete Volterra equation. We then compare this discrete equation with the continuous fractional Riccati equation from the rough Heston model. Under an admissible-strip assumption and local-in-frequency bounds, we obtain weak-error estimates for the truncated Carr--Madan pricing functional on bounded frequency windows of the form $C_1\tau^{-\alpha}+C_2(\alpha)\tau^{-(1-\alpha)}$, where the second branch comes from the discrete-to-continuous Volterra comparison. The coefficient $C_2(\alpha)$ collects the vanishing contributions arising from both the weakly singular baseline quadrature and the discrete-to-continuous resolvent comparison, and satisfies $C_2(\alpha)\to0$ as $\alpha\uparrow1^-$. 
%This is a bounded-frequency result for the truncated pricing functional rather than a full weak-error theorem for the entire European call price.
We also develop an FFT-accelerated CDQ simulator with $\mathcal O(\tau\log^2\tau)$ complexity per path and use it to price European and path-dependent options, examine the classical limit $\alpha=1$, and illustrate implied-volatility diagnostics.
\end{abstract}

\noindent%
{\it Keywords:} Rough volatility, INAR($\infty$), scaling limit, FFT-based simulation, implied volatility surface, weak-error bounds for option pricing.\\

\noindent%
{\it MSC:}  60G22, 60H35, 91G20, 62M10, 60F17.

\vfill

%%%%%%%%%%%%%%%%%%%%%%%%%%%%%%%%%%%%%%%%%%%%%%%%%%%%%%%%%%%%%%%%%%%%%%%%%%%%%%%%%%%%%%%%%%%%%%%%%%%%%%%%%%%%%%


%\tableofcontents

\section{Introduction}

The rough Heston model of \cite{el2019characteristic} is a stochastic-volatility model with discounted asset price and variance dynamics
\begin{align*}
&dS_t = S_t\sqrt{V_t}\,dW_t,
\\
&V_t = V_0 + \frac{1}{\Gamma(\alpha)}\int_0^t (t-s)^{\alpha-1}\gamma(\theta - V_s)\,ds
      + \frac{1}{\Gamma(\alpha)}\int_0^t (t-s)^{\alpha-1}\gamma\nu \sqrt{V_s}\,dB_s,
\end{align*}
where $\gamma,\theta,\nu,V_0$ are positive constants, $(W_{t})_{t\geq 0}$ and $(B_{t})_{t\geq 0}$ are standard Brownian motions with correlation $\rho$, and $\alpha\in(1/2,1)$ governs the roughness of the volatility path. It may be viewed as a rough-volatility extension of the classical Heston model \cite{heston1993closed}, motivated by empirical evidence that volatility is rough \cite{gatheral2018volatility}. Early pricing evidence under rough volatility and its implications for short-maturity skew appear in \cite{BayerFrizGatheral2015PricingRough}, while \cite{KellerResselMajid2020Comparison} clarify how rough and classical Heston specifications affect the volatility surface.

A central theoretical breakthrough is due to \cite{el2019characteristic}, who connect nearly unstable Hawkes processes to fractional volatility models and derive the characteristic function of the rough Heston log-price through a fractional Riccati equation. Their work builds on the scaling-limit theory for nearly unstable Hawkes processes developed in \cite{jaisson2015limit,jaisson2016rough}. This Hawkes-based route provides a continuous-time microstructural foundation for rough Heston. A natural question is whether one can obtain a parallel foundation in discrete time.

A convenient discrete framework is provided by INAR($\infty$) processes with Poisson thinning. If $(X_n)$ denotes the count process, then
\[
X_n\mid\mathcal F_{n-1}\sim \mathrm{Poisson}(\lambda_n),\qquad
\lambda_n=\mu+\sum_{s=1}^{n-1}\alpha_{n-s}X_s,
\]
and the cumulative process $N_n=\sum_{s=1}^n X_s$ may be viewed as a discrete Hawkes process.

Hawkes processes were introduced by \cite{hawkes1971spectra} and have since become a standard model for self-exciting event data. Functional limit theorems for Hawkes processes in various regimes are developed, for example, in \cite{bacry2013some,horst2026functional}. Their discrete-time analogs, cumulative INAR($\infty$) processes, provide an equally natural framework for count data observed on a time grid. The structural relation between INAR($\infty$) and Hawkes processes is well known: \cite{kirchner2016hawkes} develops the INAR($\infty$) framework and proves a discrete-to-continuous approximation theorem toward Hawkes processes. For broader discussions of Poisson autoregressive models and their links with INAR and Hawkes processes, see \cite{fokianos2021multivariate,huang2023nonlinear}. Discrete Hawkes models have also appeared in applications; see, for example, \cite{xu2022self}.

For rough-volatility scaling limits in discrete time, the key precursor is \cite{wang2026scaling}. They show that a heavy-tailed nearly unstable cumulative INAR($\infty$) process converges to a fractional CIR process by adapting the program of \cite{jaisson2016rough}. In particular, they consider kernels with heavy tail
\[
  \alpha_n\sim \frac{K}{n^{1+\alpha}},
  \qquad n\to\infty,
\]
for $\alpha\in(1/2,1)$, and obtain the volatility-side limit. However, their result is univariate: it identifies the rough variance dynamics, but it does not yet yield the joint price--variance model required for rough Heston option pricing.

This paper extends that discrete scaling picture from the univariate volatility setting to a bivariate price--variance framework. We study nearly unstable heavy-tailed bivariate cumulative INAR($\infty$) processes, prove that they converge to the full rough Heston model, and derive the associated microscopic-to-macroscopic parameter mapping. In this way, we obtain a discrete-time microstructural foundation for both volatility and price, complementing the continuous-time Hawkes-based viewpoint in \cite{el2018microstructural,horst2022microstructure}.

Beyond the scaling limit itself, we also study option pricing. Starting from the exact finite-$\tau$ transform recursion of the INAR approximation, we derive a second-order discrete Volterra reduction and compare it with the continuous fractional Riccati equation of rough Heston. This yields weak-error bounds for the \emph{truncated} Carr--Madan pricing functional on bounded frequency windows. The result is local in frequency rather than a full weak-error theorem for the entire European call price, because the high-frequency tail of the Fourier inversion is not analyzed in the present paper.

On the numerical side, simulation of rough Heston has been studied extensively. Existing approaches include Euler-type and strong approximations \cite{liang2017strong,RICHARD2021109,richard2023discrete}, multifactor approximations \cite{abi2019multifactor}, fast Monte Carlo and Fourier-based methods \cite{ma2022fast,callegaro2021fast,boyarchenko2025fast}, CTMC approximations \cite{yang2024}, weak simulation schemes tailored to rough Heston \cite{BayerBreneis2024WeakRoughHeston}, and related semimartingale approximations for control problems \cite{ma2023optimal}. On the weak-error side, several recent works treat rough-volatility models driven by explicit fractional-Brownian-motion inputs rather than rough Heston itself; see \cite{Gassiat2023WeakRoughVol,BayerHallTempone2021WeakLinearRoughVol,FrizSalkeldWagenhofer2024WeakRoughVol}. Our contribution is different in scope: we work from a discrete microstructural approximation and obtain a unified simulator that applies to both European and path-dependent contracts, and also to the classical limit.

The main contributions of the paper are as follows.
\begin{enumerate}
  \item \textit{A bivariate INAR microstructural limit for the rough Heston model.} We prove the joint convergence of a nearly unstable heavy-tailed bivariate cumulative INAR($\infty$) model to the rough Heston dynamics; see Theorems~\ref{theorem1} and \ref{maintheorem}. This yields a discrete-time microstructural foundation for both price and variance, together with an explicit mapping from the INAR parameters to the rough Heston coefficients and leverage correlation.

  \item \textit{A transform-based weak-error analysis on bounded frequency windows.} Starting from the exact finite-$\tau$ transform recursion, we derive a second-order discrete Volterra reduction and compare it quantitatively with the continuous fractional Riccati equation. This leads to weak-error bounds for the \emph{truncated} Carr--Madan pricing functional on bounded frequency windows of the form $C_1\tau^{-\alpha}+C_2(\alpha)\tau^{-(1-\alpha)}$ under an admissible-strip assumption and local-in-frequency control. Since $\alpha\in(1/2,1)$, the slower-decaying term is the $\tau^{-(1-\alpha)}$ branch. Its coefficient $C_2(\alpha)$ is inherited from the vanishing coefficients that appear in the weakly singular baseline approximation and in the discrete-to-continuous Volterra/resolvent comparison, and satisfies $C_2(\alpha)\to0$ as $\alpha\uparrow1^-$.

  \item \textit{FFT-accelerated simulation across rough and classical regimes.} We build a CDQ--FFT implementation of the INAR($\infty$) recursion with $\mathcal O(\tau\log^2\tau)$ cost per path. The same simulator architecture prices European and path-dependent contracts and also applies to the classical Heston case $\alpha=1$ when one uses the degenerate Markovian kernel from Remark~\ref{rem:alpha-one-kernel}.

  \item \textit{A unified numerical study across payoff classes.} We benchmark the scheme on European options, illustrate its use for Asian, lookback, and barrier contracts, compare implied-volatility slices with several recent rough Heston schemes, and examine the model-implied volatility surface.
\end{enumerate}

We also use the simulator to study the implied-volatility surface and compare the rough case $\alpha<1$ with the classical benchmark $\alpha=1$; see Figure~\ref{fig:iv_surface}. Under the parameter choices used in our experiments, the rough specification produces a steeper short-maturity at-the-money skew together with an approximately power-law-like decay, whereas the classical model yields a much flatter skew.

\paragraph{Roadmap}
Section~\ref{sec:INAR_to_rH} develops the discrete-time microstructural foundation: we introduce the bivariate cumulative INAR($\infty$) model, state the assumptions and parameterization, derive the renewal representation, and prove the scaling limit to the rough Heston model, together with the parameter and correlation mapping. Section~\ref{sec:weak-euro} develops the transform-based weak-error analysis for truncated Carr--Madan pricing on bounded frequency windows. Section~\ref{sec:numerical} presents the FFT-accelerated INAR simulator and the core numerical experiments. Section~\ref{sec:conclusion} concludes. Longer proofs and supporting technical results are collected in the appendices.


%%%%%%%%%%%%%%%%%%%%%%%%%%%%%%%%%%%%%%%%%%%%%%%%%%%%%%%%%%%%%%%%%%%%%%%%%%%%%%%%%%%%%%%%%%%
\section{From INAR\texorpdfstring{$(\infty)$}{(infty)} Processes to the Rough Heston Model}\label{sec:INAR_to_rH}

%%%%%%%%%%%%%%%%%%%%%%%%%%%%%%%%%%%%%%%%%%%%%%%%

\subsection{A discrete-time microstructural route to the rough Heston model}
Let $\iota^{(1)}, \iota^{(2)}, \iota^{(3)}, \iota^{(4)}$ denote four positive $\ell^1$ sequences.
 We consider a sequence of bivariate cumulative INAR($\infty$) processes $(N^{\tau,+}, N^{\tau,-})$, indexed by a positive integer $\tau \in \{1, 2, 3, \dots\}$ that tends to infinity. For $n \in \mathbb{N}$, their intensity is defined as
\begin{align*}
    \begin{pmatrix}
        \lambda_n^{\tau,+} \\
        \lambda_n^{\tau,-}
    \end{pmatrix}
    := \hat{\mu}_{\tau}(n) 
    \begin{pmatrix}
        1 \\
        1
    \end{pmatrix}
    + a_{\tau} \sum_{s=1}^{n-1} 
    \begin{pmatrix}
        \iota_{n-s}^{(1)} & \iota_{n-s}^{(3)} \\
        \iota_{n-s}^{(2)} & \iota_{n-s}^{(4)}
    \end{pmatrix}
    \begin{pmatrix}
        X_s^{\tau,+} \\
        X_s^{\tau,-}
    \end{pmatrix},
\end{align*}
where $\hat{\mu}_{\tau}(n)$ is the time-dependent baseline intensity, to be specified in Definition \ref{def:params}, and $a_{\tau}\in(0,1)$ denotes the decay parameter. Based on this framework, we model the tick-by-tick transaction price $P_n^{\tau}$ as
\[
    P_n^{\tau} = N_n^{\tau,+} - N_n^{\tau,-}.
\]
Next, we present a set of assumptions for our discrete model, analogous to the continuous-time setup of \cite{el2019characteristic}.

\begin{assumption}[Model Specifications]\label{ass:main}
Our model for the sequence of bivariate INAR($\infty$) processes is based on the following set of assumptions, which are analogous to those of \cite{el2019characteristic} for the continuous-time Hawkes process.
\begin{enumerate}[(i)]
    \item \textit{(Stability and Near-Instability)}
    Let 
    \[
    A:=\begin{pmatrix}
    \|\iota^{(1)}\|_1 & \|\iota^{(3)}\|_1\\
    \|\iota^{(2)}\|_1 & \|\iota^{(4)}\|_1
    \end{pmatrix},
    \quad\text{and let } \mathcal S(A) \text{ denote its spectral radius.}
    \]
    The coefficients $\iota^{(i)}$ are independent of $\tau$. We fix the normalization
    $\mathcal S(A)=1$.
    Then, for each $\tau$, $\mathcal S(a_\tau A)=a_\tau<1$.
    For the asymptotic analysis we work in the near-unstable regime, 
    \[
    a_\tau \uparrow 1,\qquad 1-a_\tau= \gamma\tau^{-\alpha},\qquad 
    \mu_\tau=\mu\tau^{\alpha-1},
    \]
    as $\tau\rightarrow\infty$, where $\alpha\in(\frac12,1),\ \gamma>0$.

    \item \textit{(No Statistical Arbitrage)} The influence kernels are balanced \emph{componentwise} so that past trades do not induce a predictable drift asymmetry between buys and sells. Precisely,
    \[
        \iota^{(1)}_n+\iota^{(3)}_n=\iota^{(2)}_n+\iota^{(4)}_n\quad\text{for all }n\ge1.
    \]
    Under the parameterization in Definition~\ref{def:params}, this is ensured by the choice of the constant matrix $\chi$ whose two rows coincide, which in turn yields $\lambda^{\tau,+}_s\equiv \lambda^{\tau,-}_s$.

    \item \textit{(Liquidity Asymmetry)} The cross-impact of trades is asymmetric. We assume that a past sell order has a stronger influence on future buy orders than a past buy order has on future sell orders:
    \[
        \iota^{(3)}=\beta \iota^{(2)}, \quad \text{for some constant } \beta>1.
    \]
    This implies $\iota^{(4)}=\iota^{(1)}+(\beta-1)\iota^{(2)}$.

    \item \textit{(Heavy-Tailed Kernels)} The influence of past trades decays slowly, following a power law. For some positive constant $C$, we assume:
    \[
        n^\alpha \sum_{s=n}^\infty \left( \iota_s^{(1)}+\beta\iota_s^{(2)} \right)\rightarrow C,\ \text{as}\ n\rightarrow \infty.
    \]
\end{enumerate}
\end{assumption}

\begin{remark}[Intuition behind the assumptions.]
\begin{itemize}
    \item \textit{Assumption~\ref{ass:main}(i)} places the model in the critical regime where cumulative self-excitation has a nondegenerate continuous-time limit.
    \item \textit{Assumption~\ref{ass:main}(ii)} rules out a predictable directional drift in the order flow.
    \item \textit{Assumption~\ref{ass:main}(iii)} is the microscopic source of the negative price-volatility correlation in the limit.
    \item \textit{Assumption~\ref{ass:main}(iv)} creates the long memory that produces rough volatility with Hurst parameter $H=\alpha-\frac12<\frac12$.
\end{itemize}
\end{remark}

Assumption~\ref{ass:main} records the structural conditions that encode the microstructural interpretation of the bivariate model. For the scaling-limit and pricing analysis below, we now work with a tractable separable one-factor specialization of that setup. This parameterization is chosen so that the balanced-impact, asymmetry, near-instability, and heavy-tail requirements are built in explicitly, closely following the framework of \cite{el2019characteristic}.

\begin{definition}[Model Parameterization]\label{def:params}
Let $\beta>1$, $1/2<\alpha<1$, $\gamma>0$, $\xi_0>0$ and $\mu>0$ be constants. The parameters of the INAR($\infty$) process sequence are specified as follows for each integer $\tau \in \{1, 2, 3, \dots\}$:
\begin{itemize}
    \item \textit{Decay and Intensity Parameters:} The decay parameter $a_\tau$ and the exogenous order intensity $\mu_\tau$ are set to:
    \begin{equation}\label{a:tau:mu:tau}
        a_\tau:=1-\gamma \tau^{-\alpha}, \quad \text{and} \quad \mu_{\tau}:=\mu\tau^{\alpha-1}.
    \end{equation}
    
    \item \textit{Influence Kernels:} The influence matrix kernel $K_n^\tau$ is given by the product of a scalar kernel $\varphi_n^\tau$ and a constant matrix $\chi$:
    \[
        K_{n}^{\tau}:=\varphi_n^{\tau} \chi, \quad \text{where} \quad \chi:=\frac1{\beta+1}\left(
        \begin{matrix}
          1 & \beta\\
          1 & \beta
        \end{matrix}
        \right).
    \]
    The scalar kernel $\varphi_n^\tau$ is a rescaled version of a base kernel $\varphi_n$:
    \[
        \varphi_n^{\tau}:=a_{\tau}\varphi_n,
    \]
    where the base kernel $(\varphi_n)_{n\ge 1}$ is defined as:
    \[
      \varphi_n:=
      \begin{cases}
        1-\frac1{\Gamma(1-\alpha)}, &n=1,\\
        \frac1{\Gamma(1-\alpha)}\left( \frac1{(n-1)^\alpha}-\frac1{n^\alpha}\right), &n\ge2.
      \end{cases}
    \]
    
    \item \textit{Baseline Intensity:} The time-dependent baseline intensity $\hat\mu_{\tau}(n)$ is defined as:
    \[
      \hat \mu_{\tau}(n):=\mu_{\tau}+\xi_0 \mu_{\tau}\left( \frac1{1-a_{\tau}}\left( 1-\sum_{s=1}^{n-1}\varphi_s^{\tau} \right)-\sum_{s=1}^{n-1}\varphi_s^{\tau} \right).
    \]
\end{itemize}
\end{definition}

\begin{remark}[Classical case $\alpha=1$]\label{rem:alpha-one-kernel}
Definition~\ref{def:params} is used for the rough regime $\alpha\in(1/2,1)$ and is not meant to be interpreted literally at $\alpha=1$. For the classical Heston benchmark in Section~\ref{sec:numerical}, we instead use the discrete Markovian kernel $\varphi_1:=1$ and $\varphi_n:=0$ for $n\ge2$, so that $\varphi_n^\tau=a_\tau\varphi_n$ and only the lag-one term remains in the recursion. This is the discrete kernel used throughout the classical-limit numerics.
\end{remark}

\begin{remark}  
  This baseline intensity follows \cite[Definition~2.1]{el2019characteristic}, but with a discrete kernel $(\varphi_n)_{n\ge1}$ chosen so that $\sum_{n=1}^\infty \varphi_n=1$ and
  \[
    n^\alpha \sum_{k=n+1}^\infty \varphi_k \to \frac{1}{\Gamma(1-\alpha)}
    \quad \text{as} \quad n \to \infty.
  \]
\end{remark}

To analyze the limit, we first express the intensity $\lambda_n^{\tau,+}$ as the solution to a discrete renewal-type equation. We start from the original definition of the intensity vector $\boldsymbol{\lambda}_n^{\tau}$:
\begin{align*}
    \begin{pmatrix}
        \lambda_n^{\tau,+} \\
        \lambda_n^{\tau,-}
    \end{pmatrix}
    = \hat{\mu}_{\tau}(n) 
    \begin{pmatrix}
        1 \\
        1
    \end{pmatrix}
    + \sum_{s=1}^{n-1} K_{n-s}^{\tau}
    \begin{pmatrix}
        X_s^{\tau,+} \\
        X_s^{\tau,-}
    \end{pmatrix} 
    = \hat{\mu}_{\tau}(n) 
    \begin{pmatrix}
        1 \\
        1
    \end{pmatrix}
    + \sum_{s=1}^{n-1} \varphi_{n-s}^{\tau} \chi
    \begin{pmatrix}
        X_s^{\tau,+} \\
        X_s^{\tau,-}
    \end{pmatrix},
\end{align*}
where we used the parameterization from Definition~\ref{def:params}. Let us focus on the first component, $\lambda_n^{\tau,+}$. Substituting the form of matrix $\chi$, we get:
\[
  \lambda_n^{\tau,+} = \hat{\mu}_{\tau}(n) + \frac{1}{1+\beta} \sum_{s=1}^{n-1} \varphi_{n-s}^{\tau} \left(X_s^{\tau,+} + \beta X_s^{\tau,-}\right).
\]
Now, let $\mathbf M_n^{\tau}=(M_n^{\tau,+},M_n^{\tau,-})^\top=\mathbf N_n^{\tau}-\sum_{s=1}^{n}\boldsymbol\lambda_s^{\tau}$ be the martingale associated with the counting process $\mathbf N_n^{\tau}$. Thus, $X_s^{\tau,\pm} = \lambda_s^{\tau,\pm} + (M_s^{\tau,\pm} - M_{s-1}^{\tau,\pm})$. Substituting this back into the expression for $\lambda_n^{\tau,+}$, and noting that $\lambda_s^{\tau,+} = \lambda_s^{\tau,-}$ (due to the structure of $\chi$), we obtain a discrete renewal equation for $\lambda_n^{\tau,+}$:
\begin{align*}
  \lambda_n^{\tau,+} 
  &= \hat{\mu}_{\tau}(n) 
  \\
  &\quad+ \frac{1}{1+\beta} \sum_{s=1}^{n-1} \varphi_{n-s}^{\tau} \left( \lambda_s^{\tau,+} + \left(M_s^{\tau,+} - M_{s-1}^{\tau,+}\right) + \beta \lambda_s^{\tau,-} + \beta\left(M_s^{\tau,-} - M_{s-1}^{\tau,-}\right) \right) \\
  &= \hat{\mu}_{\tau}(n) + \sum_{s=1}^{n-1} \varphi_{n-s}^{\tau} \lambda_s^{\tau,+} \\
  &\qquad+ \frac{1}{1+\beta} \sum_{s=1}^{n-1} \varphi_{n-s}^{\tau} \left( \left(M_s^{\tau,+} - M_{s-1}^{\tau,+}\right) + \beta\left(M_s^{\tau,-} - M_{s-1}^{\tau,-}\right) \right).
\end{align*}
This is a discrete renewal equation. Its solution can be expressed using the discrete renewal kernel $\psi^{\tau}$, which is defined as $\psi^{\tau} = \sum_{k=1}^\infty (\varphi^{\tau})^{*k}$, where $*$ denotes the discrete convolution operation. Since $\left\|\varphi^{\tau}\right\|_1 = a_\tau < 1$, the series converges and $\psi^\tau$ is well-defined. As established in Lemma~9 of \cite{wang2026scaling}, the solution is given by:
\begin{equation}\label{eq:lambda-decomp}
  \lambda_n^{\tau,+} = \hat{\mu}_\tau(n) + \sum_{s=1}^{n-1} \psi_{n-s}^\tau \hat{\mu}_\tau(s) + \frac{1}{1 + \beta} \sum_{s=1}^{n-1} \psi_{n-s}^\tau \left( M_s^{\tau,+} - M_{s-1}^{\tau,+} + \beta M_s^{\tau,-} - \beta M_{s-1}^{\tau,-} \right).
\end{equation}

%%%%%%%%%%%%%%%%%%%%%%%%%%%%%%%%%%%%%
\subsection{The rough limits of cumulative INAR\texorpdfstring{$(\infty)$}{(infty)} processes}
First, we consider how to choose the right scaling constant. We have chosen
in \eqref{a:tau:mu:tau} that
$\mu_{\tau}=\mu \tau^{\alpha-1}$,
where $\mu$ is some positive constant. To obtain a nondegenerate limit, similar to \cite{wang2026scaling}, fix $T>0$ and for $t\in[0,T]$ consider $(1-a_{\tau})\lambda_{\lfloor t\tau\rfloor}^{\tau,+}/\mu_{\tau}$, and define the renormalized intensity
\[
  C_t^{\tau}:=\frac{1-a_{\tau}}{\mu_{\tau}}\lambda_{\lfloor t\tau\rfloor}^{\tau,+}.
\]
After some calculations, this can be re-written as
\begin{equation}\label{eq:Ctau}
\begin{aligned}
  C_t^{\tau}
  &=\frac{1-a_{\tau}}{\mu_{\tau}}\hat\mu_{\tau}(\lfloor t\tau \rfloor)
   +(1-a_{\tau})\sum_{s=1}^{\lfloor t\tau \rfloor-1}\psi_{\lfloor t\tau \rfloor-s}^{\tau}\frac{\hat\mu_{\tau}(s)}{\mu_{\tau}}
   \nonumber
   \\
   &\qquad+ \kappa \tau(1-a_{\tau}) \sum_{s=1}^{\lfloor t\tau \rfloor-1}\psi_{\lfloor t\tau \rfloor-s}^{\tau}\sqrt{C_{s/\tau}^{\tau}}
  \left( B_{(s+1)/\tau}^{\tau}-B_{s/\tau}^{\tau} \right),
\end{aligned}
\end{equation}
where 
\begin{equation}
    \kappa:=\sqrt{\frac{1+\beta^2}{\gamma\mu(1+\beta)^2}},
\end{equation}
and for any $t\in[0,T]$,
\begin{equation}\label{eq:inar-mart}
B_t^{\tau}:=\frac1{\sqrt{\tau}}\sum_{s=1}^{\lfloor t\tau \rfloor-1}\frac{M_s^{\tau,+}-M_{s-1}^{\tau,+}+\beta(M_s^{\tau,-}-M_{s-1}^{\tau,-})}{\sqrt{\lambda_s^{\tau,+}+\beta^2\lambda_s^{\tau,-}}},
\end{equation}
which will converge weakly to a standard Brownian motion $B$ as $\tau\rightarrow\infty$ by the martingale central limit theorem; see \cite{wang2026scaling}. 
We first make the vector notation explicit. 
For $n\ge 0$, let 
\[
\mathbf N_n^{\tau}:=\begin{pmatrix}N_n^{\tau,+}\\ N_n^{\tau,-}\end{pmatrix}
\quad\text{and}\quad
\boldsymbol\lambda_n^{\tau}:=\mathbb E\left[\left.\mathbf N_n^{\tau}-\mathbf N_{n-1}^{\tau}\right|\mathcal F_{n-1}\right]
=\begin{pmatrix}\lambda_n^{\tau,+}\\ \lambda_n^{\tau,-}\end{pmatrix},
\]
where $(\mathcal F_n)_{n\ge0}$ is the natural filtration, 
$\mathbf N_n^{\tau}$ is the bivariate counting process up to step $n$, and 
$\boldsymbol\lambda_n^{\tau}$ is its (predictable) intensity vector. 
With this notation, fixing $T>0$, for $t\in[0,T]$, we define the rescaled cumulative counts, cumulative intensities, 
and the associated normalized martingale as
\[
    \mathbf Y_t^{\tau}
    :=\frac{1-a_{\tau}}{\tau^\alpha \mu}\mathbf N_{\lfloor t\tau \rfloor}^{\tau},\qquad
    \boldsymbol\Lambda_t^{\tau}
    :=\frac{1-a_{\tau}}{\tau^\alpha \mu}\sum_{s=1}^{\lfloor t\tau \rfloor}\boldsymbol\lambda_s^{\tau},\qquad 
    \mathbf Z_t^{\tau}
    :=\sqrt{\frac{\tau^\alpha\mu}{1-a_{\tau}}}\left(\mathbf Y_t^{\tau}-\boldsymbol\Lambda_t^{\tau}\right).
\]

A key technical input is that the rescaled renewal kernel associated with $\psi^\tau$ converges weakly to the Mittag--Leffler resolvent density
\[
  f_{\alpha,\gamma}(t)=\gamma t^{\alpha-1}E_{\alpha,\alpha}(-\gamma t^\alpha);
\]
see Appendix~\ref{supp:scaling-proofs}.

\iffalse
\begin{table}[ht]
\centering
\caption{Numerical comparison between the symbol-level coefficient
$|c_1(\alpha)|$ and the baseline coefficient
$\left|1-\Gamma(2-\alpha)^{-1}\right|$ as $\alpha\uparrow1^-$.}
\label{tab:c1-vs-baseline}
\small
\setlength{\tabcolsep}{10pt}
\begin{tabular}{ccc}
\toprule
$\alpha$ & $|c_1(\alpha)|$ & $\left|1-\dfrac{1}{\Gamma(2-\alpha)}\right|$ \\
\midrule
0.70 & 0.0713 & 0.1141 \\
0.90 & 0.0088 & 0.0510 \\
0.95 & 0.0022 & 0.0271 \\
0.99 & $9.1\times10^{-5}$ & 0.00573 \\
\bottomrule
\end{tabular}
\end{table}
\fi

\begin{proposition}\label{proposition1}
  Fix $T>0$. Under the parameterization in Definition~\ref{def:params}, the sequence
  $(\boldsymbol\Lambda^{\tau},\mathbf Y^{\tau},\mathbf Z^{\tau})$ is $C$-tight on $[0,T]$.
  Moreover, if $(\mathbf Y,\mathbf Z)$ is a limit point of
  $(\mathbf Y^{\tau},\mathbf Z^{\tau})$, then $\mathbf Z$ is a continuous martingale with
  \[
    [\mathbf Z,\mathbf Z]=\mathrm{diag}(\mathbf Y),
  \]
  that is,
$[Z^+,Z^+]_t=Y_t^+$,
$[Z^-,Z^-]_t=Y_t^-$
and $[Z^+,Z^-]_t=0$, for any $t\in[0,T]$.
\end{proposition}

\begin{proof}
{$C$-tightness of $(\boldsymbol\Lambda^\tau,\mathbf Y^\tau)$ is verified by applying Lemma~\ref{xu2023} componentwise, using the intensity-moment bounds $\mathbb{E}[\lambda_n^{\tau,+}]\lesssim\tau^{2\alpha-1}$ and $\mathbb{E}[(\lambda_n^{\tau,+})^2]\lesssim\tau^{4\alpha-2}$. For $\mathbf{Z}^\tau$, an analogous argument with the martingale increments establishes the required oscillation estimate.
The diagonal quadratic covariation $[Z^{\tau,+},Z^{\tau,+}]=Y^{\tau,+}$ follows from the conditional Poisson variance $\mathrm{Var}(X_s^{\tau,+}\mid\mathcal{F}_{s-1})=\lambda_s^{\tau,+}$, and the vanishing cross-variation $[Z^{\tau,+},Z^{\tau,-}]\to0$ follows from the independence of the two Poisson components.} Full details are in Appendix~\ref{supp:scaling-proofs}.
\end{proof}

Since the components of $(\mathbf Y^{\tau})$ and $(\boldsymbol\Lambda^{\tau})$ are pure jump processes, the classical Kolmogorov--Chentsov criterion is not applicable. We instead use a $C$-tightness criterion for c\`adl\`ag processes from \cite[Lemma 3.5]{horst2023convergence}; its precise statement and the verification of its hypotheses in the present setting are given in Appendix~\ref{supp:scaling-proofs}.

We apply this criterion on a fixed horizon $[0,T]$ to each component of our scaled processes; the required moment estimates are deferred to Appendix~\ref{supp:scaling-proofs}.

%\subsubsection{Convergence of $Y^\tau$ and $Z^\tau$}\label{app:tightness}
Before we prove the convergence of $Y^\tau$ and $Z^\tau$, let us first prove the following lemma.

\begin{lemma}\label{lem:tightness}
  Fix $T>0$. Then
  \[
  \sup_{t \in [0,T]} \left\|\boldsymbol{\Lambda}_t^\tau - \mathbf{Y}_t^\tau\right\|_1 \rightarrow 0 \text{ in probability},
  \]
as $\tau \rightarrow \infty$, where $\|\cdot\|_1$ denotes the $\ell^1$-norm on $\mathbb{R}^2$.
\end{lemma}

\begin{proof}
{Because $\Lambda_t^{\tau,+}=\Lambda_t^{\tau,-}$ holds exactly (both equal the same renewal convolution), the $\ell^1$-difference $\|\boldsymbol\Lambda_t^\tau-\mathbf{Y}_t^\tau\|_1$ equals $2|Y_t^{\tau,+}-\Lambda_t^{\tau,+}|$. The compensator--process gap satisfies $|Y_t^{\tau,+}-\Lambda_t^{\tau,+}|=(1-a_\tau)/(\tau^\alpha\mu)\cdot|M^{\tau,+}_{\lfloor t\tau\rfloor}|$, and the martingale maximal inequality together with $\mathbb{E}[(\lambda_n^{\tau,+})^2]\lesssim\tau^{4\alpha-2}$ gives $\sup_t\mathbb{E}[|M^{\tau,+}_{\lfloor t\tau\rfloor}|^2]\lesssim\tau^{4\alpha-1}$. Multiplying by $(1-a_\tau)^2/(\tau^\alpha\mu)^2\asymp\tau^{-4\alpha}$ yields $\sup_t|Y_t^{\tau,+}-\Lambda_t^{\tau,+}|\to0$ in probability.} Full details are in Appendix~\ref{supp:scaling-proofs}.
\end{proof}

Lemma~\ref{lem:tightness} implies the uniform convergence to zero in probability of $Y^{\tau,+} - \Lambda^{\tau,+}$.

Next, by noting that
$\Lambda_t^{\tau,+} = \Lambda_t^{\tau,-}$,
we obtain
\[
  \sup_{t \in [0,T]} \left|Y_t^{\tau,+} - Y_t^{\tau,-}\right| \rightarrow 0,
\]
as $\tau \rightarrow \infty$. Therefore, if a subsequence of $Y^{\tau,+}$ converges to some $Y$, then the corresponding subsequence of $Y^{\tau,-}$ converges to the same $Y$. We thus have the following proposition regarding the limit points of $Y^{\tau,+}$ and $Y^{\tau,-}$.

\begin{proposition}\label{prop2}
  If $(Y,Y,Z^+,Z^-)$ is a limit point for $(Y^{\tau,+},Y^{\tau,-},Z^{\tau,+},Z^{\tau,-})$, then this limit point admits a representation of the form
  \[
    Y_t=\int_0^t v_sds,\qquad Z_t^+=\int_0^t\sqrt{v_s}dB_s^1,\qquad Z_t^-=\int_0^t\sqrt{v_s}dB_s^2,
  \]
  for some two-dimensional Brownian motion $(B^1,B^2)$ with independent components and some rate process $v=(v_t)_{t\ge 0}$. In particular, the corresponding rate process $v$ satisfies
  \begin{equation}\label{finalequation1}
    v_t
    = \xi_0\left(1-F_{\alpha,\gamma}(t)\right)+F_{\alpha,\gamma}(t)
    + \sqrt{\frac{1+\beta^2}{\gamma\mu(1+\beta)^2}}
      \int_0^t f_{\alpha,\gamma}(t-s)\sqrt{v_s}\,dW_s^V,
  \end{equation}
  where
    \[
    W^V:=\frac{B^1+\beta B^2}{\sqrt{1+\beta^2}},
  \]
  and for $t>0$,
    \[
    f_{\alpha,\gamma}(t):=\gamma t^{\alpha-1}E_{\alpha,\alpha}(-\gamma t^\alpha),
    \qquad
    F_{\alpha,\gamma}(t):=\int_0^t f_{\alpha,\gamma}(u)\,du
    =1-E_{\alpha}(-\gamma t^\alpha),
  \]
  with the convention $F_{\alpha,\gamma}(0)=0$. Here $E_{\alpha,\beta}(z)=\sum_{k=0}^\infty \frac{z^k}{\Gamma(\alpha k+\beta)}$ is the two-parameter Mittag--Leffler function (and $E_\alpha=E_{\alpha,1}$). Equivalently, $\mathcal{L}\{f_{\alpha,\gamma}\}(z)=\frac{\gamma}{\gamma+z^\alpha}$ for $z>0$.
  Furthermore, for any $\epsilon>0$, the process $v$ has H\"older regularity $\alpha-\frac12-\epsilon$.
\end{proposition}

\begin{proof}
{Fix a convergent subsequence $(\mathbf{Y}^{\tau_k},\mathbf{Z}^{\tau_k})\to(\mathbf{Y},\mathbf{Z})$. By Lemma~\ref{lem:tightness} the limit satisfies $Y^+=Y^-=:Y$, and by Proposition~\ref{proposition1} the limiting martingale $\mathbf{Z}$ has diagonal quadratic variation. The martingale representation theorem then gives the Brownian form $(Z^+,Z^-)=(\int\sqrt{v_s}\,dB^1,\int\sqrt{v_s}\,dB^2)$. The discrete renewal sum in \eqref{eq:Ctau} is then passed to the limit: by the weak kernel convergence of Lemma~\ref{lem:kernel_convergence}, $\frac{1-a_\tau}{\tau}\sum_s\psi_{\lfloor t\tau\rfloor-s}^\tau(\cdot)$ converges to $\int_0^t f_{\alpha,\gamma}(t-s)(\cdot)\,ds$, yielding the Volterra equation \eqref{finalequation1}.} Full details are in Appendix~\ref{supp:scaling-proofs}.
\end{proof}

We then state the main result of the paper. 

\begin{theorem}[Scaling limit for cumulative intensities and martingale fluctuations]\label{theorem1}
  Fix $T>0$. Under Assumption~\ref{ass:main} and the model parameterization specified in Definition~\ref{def:params}, as $\tau\rightarrow \infty$, the sequence of processes $(\boldsymbol\Lambda_t^{\tau},\mathbf Y_t^{\tau},\mathbf Z_t^{\tau})_{t\in[0,T]}$ converges in law under the Skorokhod topology on $D([0,T],\mathbb{R}^6)$ to $(\boldsymbol\Lambda,\mathbf Y,\mathbf Z)$, where the limit process is described by a rate process $v = (v_t)_{t\ge 0}$ as follows:
  \[
    \boldsymbol\Lambda_t=\mathbf Y_t=\int_0^t v_s ds\left(  \begin{matrix}
     1\\1
    \end{matrix}\right), \quad \text{and} \quad \mathbf Z_t=\int_0^t\sqrt{v_s}\left(\begin{matrix}
    dB_s^1\\dB_s^2
    \end{matrix}\right),
  \]
  and the rate process $v$ is the unique solution of the following stochastic Volterra equation:
  \begin{equation}\label{eq:vt}
    v_t=\xi_0+\frac1{\Gamma(\alpha)}\int_0^t(t-s)^{\alpha-1}\gamma(1-v_s)ds+\gamma \sqrt{\frac{1+\beta^2}{\gamma \mu(1+\beta)^2}}\frac1{\Gamma(\alpha)}\int_0^t(t-s)^{\alpha-1}\sqrt{v_s}dW_s^V,
  \end{equation}
  where $W^V:=\frac{B^1+\beta B^2}{\sqrt{1+\beta^2}}$. Furthermore, for any $\epsilon>0$, the process $v$ has H\"older regularity $\alpha-1/2-\epsilon$. 
\end{theorem}

\begin{proof}
{The proof proceeds in three stages; full details are in Appendix~\ref{supp:scaling-proofs}.

\medskip
\noindent\textit{Stage 1: $C$-tightness and martingale structure (Proposition~\ref{proposition1}).}
We apply the $C$-tightness criterion of Lemma~\ref{xu2023} componentwise to $(Y^{\tau,+}, Y^{\tau,-}, Z^{\tau,+}, Z^{\tau,-})$. The key moment estimates are $\mathbb{E}[\lambda_n^{\tau,+}]\lesssim\tau^{2\alpha-1}$ and $\mathbb{E}[(\lambda_n^{\tau,+})^2]\lesssim\tau^{4\alpha-2}$, obtained from the renewal decomposition \eqref{eq:lambda-decomp} together with the pointwise resolvent bound $\psi_j^\tau\le C_T j^{\alpha-1}$ (Lemma~\ref{lem:psi-pointwise}). These estimates yield the required $h^{b_i}/\tau^{a_i}$ bound with $\varrho_*>1$. Condition~(i) of Lemma~\ref{xu2023} (small jumps) follows from the rescaling factor $(1-a_\tau)/(\tau^\alpha\mu)\to0$. The limiting martingale $\mathbf{Z}$ inherits the diagonal quadratic variation $[Z^+,Z^+]_t=Y_t^+$, $[Z^-,Z^-]_t=Y_t^-$, $[Z^+,Z^-]_t=0$.

\medskip
\noindent\textit{Stage 2: Characterization of limit points (Proposition~\ref{prop2}).}
From Lemma~\ref{lem:tightness}, the two intensity components $\Lambda^{\tau,+}$ and $\Lambda^{\tau,-}$ collapse uniformly, so every limit point takes the form $(Y, Y, Z^+, Z^-)$. The martingale representation theorem, applied to the diagonal quadratic variation from Stage~1, yields $Z^+_t=\int_0^t\sqrt{v_s}\,dB^1_s$ and $Z^-_t=\int_0^t\sqrt{v_s}\,dB^2_s$ for independent Brownian motions $(B^1,B^2)$. Passing the discrete renewal decomposition to the limit via the weak kernel convergence of Lemma~\ref{lem:kernel_convergence} — specifically, $\mathsf{S}^\tau(t)\,dt\Rightarrow f_{\alpha,\gamma}(t)\,dt$ — converts the discrete convolution sum into the fractional Volterra integral, yielding \eqref{finalequation1}.

\medskip
\noindent\textit{Stage 3: Uniqueness and convergence.}
Standard pathwise Gronwall estimates for stochastic Volterra equations of the form \eqref{eq:vt} yield uniqueness of the solution $v$; see Appendix~\ref{supp:scaling-proofs}. Since all limit points satisfy the same equation and the solution is unique, the whole sequence converges.}
\end{proof}

%%%%%%%%%%%%%%%%%%%%%%%%%%%%%
\subsection{Parameter mapping and correlation for the macroscopic limit}

Throughout the limit, we define the Brownian motions driving price and variance via the orthogonal decomposition
\begin{align}
W_t^X := \frac{1}{\sqrt{2}}\left(B_t^1-B_t^2\right),
\qquad
W_t^V := \frac{B_t^1+\beta B_t^2}{\sqrt{1+\beta^2}},\label{eq:wtxv}
\end{align}
with $(B^1,B^2)$ a two-dimensional standard Brownian motion with independent components. A direct calculation gives the correlation
\begin{equation}\label{eq:rho-def}
  d\left\langle W^X,W^V\right\rangle_t
  =\frac{1-\beta}{\sqrt{2(1+\beta^2)}}dt
  =:\rho dt,
\end{equation}
where $\rho\in\left(-\frac1{\sqrt2},0\right)$ is consistent with the leverage effect induced by the liquidity asymmetry $\beta>1$.

Moreover, moving from the rate process $v$ (for intensities) to the variance process
\[
  V_t := \theta v_t,
\]
rescales the drift to $\gamma(\theta-V_t)$ and the diffusion coefficient accordingly. Matching the diffusion scale obtained from the INAR martingale CLT yields the identity
\begin{equation}\label{eq:mu-nu-map}
  \nu = \sqrt{\frac{\theta(1+\beta^2)}{\gamma\mu(1+\beta)^2}}
  \qquad\Longleftrightarrow\qquad
  \mu = \frac{\theta(1+\beta^2)}{\gamma\nu^2(1+\beta)^2}.
\end{equation}
We will henceforth \emph{define} $\nu$ via \eqref{eq:mu-nu-map}, so that, with $V=\theta v$ and the correlation $\rho$ in \eqref{eq:rho-def}, the limiting variance dynamics takes the rough Heston form in \eqref{eq:V}.

Then we can also build microscopic processes converging to the log-price. %like of \cite{el2019characteristic}. 
For $\theta>0$, fixing $T>0$, and $t\in[0,T]$, define the (already properly rescaled) microscopic log-price process $P^\tau=(P_t^\tau)_{t\in[0,T]}$ as
\[
  P_t^{\tau}
  :=\sqrt{\frac{\theta}{2}}\sqrt{\frac{1-a_{\tau}}{\mu \tau^\alpha}}
    \left( N_{\lfloor t \tau \rfloor}^{\tau,+}-N_{\lfloor t \tau \rfloor}^{\tau,-} \right)
   -\frac{\theta}{2}\frac{1-a_{\tau}}{\mu \tau^\alpha}
    N_{\lfloor t \tau \rfloor}^{\tau,+}
  =\sqrt{\frac{\theta}{2}}\left(Z_{t}^{\tau,+}-Z_{t}^{\tau,-}\right)
   -\frac{\theta}{2}Y_{t}^{\tau,+}.
\]

From Theorem~\ref{theorem1}, $(P^{\tau})_{t\in[0,T]}$ converges in law (in the Skorokhod topology on $D([0,T])$) to the macroscopic price process $(P_t)_{t\in[0,T]}$ given by
\[
 P_t := \sqrt{\frac\theta2}\int_0^t\sqrt{v_s}\left(dB_s^1-dB_s^2\right)-\frac\theta2\int_0^t v_sds.
\]
Let $V_t:=\theta v_t$ and $W_t^X:=\frac{1}{\sqrt{2}}(B_t^1-B_t^2)$. 
It follows from the equation for $v_t$ in \eqref{eq:vt} that $V_t$ represents the variance process. 
Then we have the following theorem.

\begin{theorem}[Price-process consequence: rough Heston limit]\label{maintheorem}
  Fix $T>0$. Under the same conditions as Theorem \ref{theorem1}, as $\tau\to\infty$, the rescaled price processes $(P_t^{\tau})_{t\in[0,T]}$ converge in law (Skorokhod topology on $D([0,T])$) to
  \begin{equation}\label{eq:P}
    P_t=\int_0^t\sqrt{V_s}dW_s^X-\frac12\int_0^tV_sds,
  \end{equation}
  where $V_t=\theta v_t$ is the unique solution of the rough Volterra equation:
  \begin{equation}\label{eq:V}
    V_t=\theta\xi_0+\frac{\gamma}{\Gamma(\alpha)}\int_0^t (t-s)^{\alpha-1}\left(\theta-V_s\right)ds
    +\frac{\gamma\nu}{\Gamma(\alpha)}\int_0^t (t-s)^{\alpha-1}\sqrt{V_s}dW_s^V,
  \end{equation}
  where $\nu$ is defined from the INAR parameters via \eqref{eq:mu-nu-map} and $\left(W^X,W^V\right)$ is a two-dimensional Brownian motion defined in \eqref{eq:wtxv} whose correlation is
\begin{equation}\label{correlation}
    d\left\langle W^X,W^V\right\rangle_t=\rho dt,\qquad
    \rho=\frac{1-\beta}{\sqrt{2(1+\beta^2)}}.
  \end{equation}
\end{theorem}

\begin{proof}
{The rescaled price $P_t^\tau=\sqrt{\theta/2}(Z_t^{\tau,+}-Z_t^{\tau,-})-(\theta/2)Y_t^{\tau,+}$ is a continuous linear functional of $(Y^{\tau,+},Z^{\tau,+},Z^{\tau,-})$. Theorem~\ref{theorem1} gives joint convergence of these three processes, so the continuous-mapping theorem yields $P^\tau\to P$ in the Skorokhod topology. Substituting the limit $Z^+_t-Z^-_t=\int_0^t\sqrt{v_s}(dB^1_s-dB^2_s)$ and identifying $W^X=\frac{1}{\sqrt{2}}(B^1-B^2)$, together with the parameter maps \eqref{eq:rho-def} and \eqref{eq:mu-nu-map}, converts the limit into the rough Heston form \eqref{eq:P}--\eqref{eq:V}.} Full details are in Appendix~\ref{supp:scaling-proofs}.
\end{proof}

\begin{remark}
   From \eqref{correlation} and the condition that $\beta > 1$, the correlation coefficient $\rho$ between the two Brownian motions must lie in the interval $\left(-\frac{1}{\sqrt{2}}, 0\right)$, which is consistent with \cite{el2019characteristic}.
\end{remark}

%%%%%%%%%%%%%%%%%%%%%%%%%%%%%%%%%%%%

\section{Weak-Error Bounds for Truncated Carr--Madan Pricing}\label{sec:weak-euro}

In this section we analyze the weak error for the truncated Carr--Madan pricing functional associated with the finite-$\tau$ INAR($\infty$) approximation. The logic of this section is as follows. We start from the exact finite-$\tau$ transform recursion, which takes the form of an exponential discrete Volterra recursion. On bounded vertical strips, this exact recursion can be reduced by a second-order expansion to a quadratic discrete Volterra equation. After a deterministic rescaling, the reduced recursion matches the coefficient structure of the rough Heston fractional Riccati equation. We then invert the discrete renewal operator and compare the resulting cumulative resolvents in discrete and continuous time. This resolvent comparison is the key technical step: once it is established, it transfers to the Riccati variables, then to the logarithmic transforms, then to the transforms themselves, and finally to the truncated Carr--Madan pricing functional on bounded frequency windows.

The slower branch $\tau^{-(1-\alpha)}$ has a specific analytic origin that we make explicit here. The first source is the weighted baseline-sum approximation, which isolates the weakly singular quadrature defect coming from the baseline term:

\begin{proposition}[Weighted baseline-sum approximation]\label{prop:baseline-riemann}
Fix $T>0$ and let $N_\tau:=\lfloor\tau T\rfloor$. With
\[
  \varepsilon_\tau:=1-a_\tau=\gamma\tau^{-\alpha},
\]
for every $f\in C^\alpha([0,T])$, there exists a constant $C_T>0$, depending
only on the fixed model parameters and on $T$, but independent of $\tau$ and
$f$, such that
\begin{equation}\label{eq:baseline-riemann}
\begin{aligned}
  &\left|
    \frac{\varepsilon_\tau}{\mu}
    \sum_{m=1}^{N_\tau}\hat\mu_\tau(m)f\!\left(T-\frac{m}{\tau}\right)
    -
    \gamma\int_0^T f(s)\,ds
    -
    \xi_0 I^{1-\alpha}f(T)
  \right|
  \\
  &\le
  C_T\|f\|_{C^\alpha([0,T])}\tau^{-\alpha}
  +
  C_{\mathrm{base}}(\alpha)\|f\|_{C^\alpha([0,T])}\tau^{-(1-\alpha)},
\end{aligned}
\end{equation}
where
\[
  I^{1-\alpha}f(T)
  :=
  \frac{1}{\Gamma(1-\alpha)}
  \int_0^T(T-s)^{-\alpha}f(s)\,ds,
\]
and the coefficient $C_{\mathrm{base}}(\alpha)\ge0$ satisfies
$C_{\mathrm{base}}(\alpha)\to0$ as $\alpha\uparrow1^-$.
\end{proposition}

\begin{proof}
See Appendix~\ref{supp:weak-error-proofs}.
\end{proof}

The second source is the following small-$s$ expansion of the kernel transform,
which is the direct analytic input to the resolvent comparison in
Section~\ref{sec:renewal-comparison}.

\begin{lemma}[Small-$s$ expansion of the kernel transform]\label{lem:phi-expansion}
For the kernel $(\varphi_n)_{n\ge1}$ from Definition~\ref{def:params}, let $\hat\varphi(s):=\sum_{n\ge1}\varphi_n e^{-sn}$ for $s>0$. Then, as $s\downarrow0$,
\begin{equation}\label{eq:phi-expansion-main}
  1-\hat\varphi(s)=s^\alpha+c_1(\alpha)s+\mathcal O(s^{1+\alpha}),
\end{equation}
where
\begin{equation}\label{eq:c1-alpha-main}
  c_1(\alpha):=1+\frac{\zeta(\alpha)}{\Gamma(1-\alpha)}.
\end{equation}
Moreover, $c_1(\alpha)\to0$ as $\alpha\uparrow1^-$.
\end{lemma}

\begin{proof}
See Appendix~\ref{app:phi-expansion}.
\end{proof}

\begin{remark}
The vanishing of both $c_1(\alpha)$ and $C_{\mathrm{base}}(\alpha)$ as
$\alpha\uparrow1^-$ is the direct analytic reason why $C_2(\alpha)\to0$ in the
weak-error bound. Proposition~\ref{prop:baseline-riemann} identifies the
baseline quadrature contribution, while Lemma~\ref{lem:phi-expansion}
identifies the symbol-level correction in the kernel transform. In the Laplace
domain, the first discrete-to-continuous correction to the cumulative resolvent
is proportional to $c_1(\alpha)\tau^{-(1-\alpha)}$; once
$c_1(\alpha)\to0$, this correction shrinks to zero and the two-branch rate
degenerates to the single $\tau^{-\alpha}$ branch of the classical
($\alpha=1$) case.
\end{remark}

These two inputs feed the discrete/continuous resolvent comparison developed in Section~\ref{sec:renewal-comparison}; the technical details appear in Appendix~\ref{supp:weak-error-proofs}.

Throughout this section, unless stated otherwise, the parameters defining the bounded frequency window and the pricing functional are regarded as fixed. In particular, when discussing truncated Carr--Madan pricing below, $\eta$, $R$, $T$, $K$, and $S_0$ are fixed, and all constants are allowed to depend on these fixed quantities.


We first state the common analytic-strip assumption under which all subsequent transform and pricing comparisons are carried out.

\begin{assumption}[Admissible transform strip]\label{ass:strip}
Fix the maturity $T>0$ considered in this section. We assume that there exists a nonempty open vertical strip
\[
  \mathcal S_T:=\left\{z\in\C:\underline\eta_T<\mathrm{Re}\ z<\overline\eta_T\right\}
\]
such that the continuous transform $\phi(z,T)$ and the discrete transforms $\Phi^\tau(z,T)$ are finite and analytic on $\mathcal S_T$ for all sufficiently large $\tau$. All weak-error statements below are formulated for vertical lines $\mathrm{Re}\ z=\eta$ contained in this common admissible strip. Establishing such a strip directly from the discrete dynamics is outside the scope of the present paper.
\end{assumption}

{
\begin{assumption}[Bounded-strip Riccati regularity]\label{ass:riccati-bound}
Fix $\eta\in\R$ and $R>0$ with $\Gamma_{\eta,R}\subset\mathcal S_T$. We assume that the continuous fractional Riccati solution satisfies the bounded-strip supremum bound
\[
  M_{\eta,R,T}
  :=
  \sup_{z\in\Gamma_{\eta,R}}\sup_{0\le t\le T}|h(t,z)|
  <\infty.
\]
\end{assumption}

\begin{remark}
Together, Assumptions~\ref{ass:strip} and~\ref{ass:riccati-bound} constitute the two analytic hypotheses under which all subsequent transform and pricing comparisons are carried out. Assumption~\ref{ass:strip} provides the common domain of analyticity; Assumption~\ref{ass:riccati-bound} ensures that the continuous Riccati solution remains uniformly bounded on the relevant frequency strip, which is the regularity input needed to control the nonlinear comparison in Proposition~\ref{prop:Volterra-rate-fixed}. Parameter-specific numerical checks for both assumptions are reported in Subsection~\ref{subsec:numerical-assumption-checks}.
\end{remark}
}

Throughout, $T>0$ is fixed, $N_\tau:=\lfloor \tau T\rfloor$, $t_n:=n/\tau$, and $k:=\log(K/S_0)$ denotes log-moneyness. We write
\begin{equation}\label{eq:ctau-dtau-weak}
  c_\tau:=\sqrt{\frac{\theta(1-a_\tau)}{2\mu\tau^\alpha}},
  \qquad
  d_\tau:=\frac{\theta(1-a_\tau)}{2\mu\tau^\alpha},
\end{equation}
so that the microscopic log-price at maturity can be written as
\begin{equation}\label{eq:Ptau-weak}
  P_T^\tau
  =c_\tau\left(N_{N_\tau}^{\tau,+}-N_{N_\tau}^{\tau,-}\right)-d_\tau N_{N_\tau}^{\tau,+}.
\end{equation}

%%%%%%%%%%%%%%%%%%%%%%%%%%%%%%%%%%%%%%%%%%%%%%%%

\subsection{Continuous fractional Riccati equation}

Let $P_T=\log(S_T/S_0)$ denote the rough Heston log-price at maturity $T$.
By \cite{el2019characteristic}, its log-price transform admits the representation
\begin{equation}\label{eq:phi-cont-weak}
\phi(z,T):=\E\!\left[e^{zP_T}\right]
  =\exp\!\left(\gamma\theta\, I^1 h(\cdot,z)(T)+V_0\, I^{1-\alpha}h(\cdot,z)(T)\right),
\end{equation}
for complex $z$ in the analyticity strip under consideration, where, for $r\in(0,1]$,
\[
  I^r f(t):=\frac{1}{\Gamma(r)}\int_0^t (t-s)^{r-1}f(s)\,ds,
\]
and $h(\cdot,z)$ solves the fractional Riccati equation
\begin{equation}\label{eq:riccati-cont-weak}
  D^\alpha h(t,z)
  =\frac12(z^2-z)+\gamma(\rho\nu z-1)h(t,z)
   +\frac{(\gamma\nu)^2}{2}\left(h(t,z)\right)^2,
  \qquad I^{1-\alpha}h(0,z)=0.
\end{equation}
For real $u$, the characteristic function is recovered as $\phi(\ii u,T)$.
Here, the fractional derivative $D^\alpha$ is defined by
\[
  D^\alpha f(t)
  :=\frac{1}{\Gamma(1-\alpha)}
    \frac{d}{dt}\int_0^t (t-s)^{-\alpha}f(s)\,ds.
\]
Equivalently, applying $I^\alpha$ to \eqref{eq:riccati-cont-weak}, one obtains the Volterra form
\begin{equation}\label{eq:riccati-volterra-weak}
  h(t,z)
  =\frac{1}{\Gamma(\alpha)}\int_0^t
    (t-s)^{\alpha-1}
    \left[
      \frac12(z^2-z)
      +\gamma(\rho\nu z-1)h(s,z)
      +\frac{(\gamma\nu)^2}{2}\left(h(s,z)\right)^2
    \right]ds.
\end{equation}

%%%%%%%%%%%%%%%%%%%%%%%%%%%%%%%%%%%%%%%%%%%%%%%%

\subsection{Exact discrete transform recursion}\label{subsec:exact-discrete-transform}

In this subsection we derive the \emph{exact} finite-$\tau$ transform recursion for the rescaled INAR($\infty$) log-price. The key observation is that, because the offspring mechanism is Poisson, the bivariate INAR($\infty$) process admits a discrete-time branching--immigration representation, entirely analogous in spirit to the cluster representation used for Hawkes processes. In the present setting, however, the clusters evolve on the integer time grid, and the relevant recursion is therefore a \emph{discrete} Volterra recursion.

Throughout this subsection, we use the following convention for the bivariate influence matrix: the \emph{rows} correspond to the child type (future arrival type), while the \emph{columns} correspond to the parent type (source event type). Under the parameterization of Definition~\ref{def:params},
\[
  \chi=\frac1{1+\beta}
  \begin{pmatrix}
    1 & \beta\\
    1 & \beta
  \end{pmatrix},
\]
so that a ``$+$''-parent produces ``$+$''- and ``$-$''-children with the same lag-$k$ mean $\varphi_k^\tau/(1+\beta)$, whereas a ``$-$''-parent produces ``$+$''- and ``$-$''-children with the same lag-$k$ mean $\beta\varphi_k^\tau/(1+\beta)$.


We write
\[
  \mathcal D_\tau(T):=\left\{z\in\mathbb C:\ \E\left[\left|e^{zP_T^\tau}\right|\right]<\infty\right\}
\]
for the finite-$\tau$ moment domain at maturity $T$.

For $z\in\mathcal D_\tau(T)$, define the one-step exponential weights
\begin{equation}\label{eq:Theta-pm-exact}
  \Theta_+^\tau(z):=z(c_\tau-d_\tau),
  \qquad
  \Theta_-^\tau(z):=-zc_\tau,
\end{equation}
where $c_\tau$ and $d_\tau$ are defined in \eqref{eq:ctau-dtau-weak}.
These are the direct contributions of one ``$+$''-arrival and one ``$-$''-arrival, respectively, to the exponent $zP_T^\tau$.

Next, introduce the \emph{parent-type aggregated lag coefficients}
\begin{equation}\label{eq:q-pm-exact}
  q_k^{\tau,+}:=\frac{\varphi_k^\tau}{1+\beta},
  \qquad
  q_k^{\tau,-}:=\frac{\beta\varphi_k^\tau}{1+\beta},
  \qquad k\ge 1.
\end{equation}
Here $q_k^{\tau,+}$ corresponds to a ``$+$''-type parent and $q_k^{\tau,-}$ to a ``$-$''-type parent. These coefficients are \emph{not} indexed by child type: rather, they give the common lag-$k$ Poisson mean with which a parent of the specified type produces each of the two child types.

We now introduce the cluster transforms. Consider a single immigrant born at the current time and let its descendants evolve over the next $n$ time steps according to the bivariate Poisson branching mechanism of the INAR($\infty$) model. Let
\[
  \mathcal M_n^{\tau,+}(z)
  \quad\text{and}\quad
  \mathcal M_n^{\tau,-}(z)
\]
denote the transforms of the total contribution of such a cluster to the exponent $zP_T^\tau$, when the initial immigrant is of type $+$ and $-$, respectively. Since the ancestor itself contributes exactly one ``$+$''-jump or one ``$-$''-jump at its birth time, we have the initial conditions
\begin{equation}\label{eq:M0-exact}
  \mathcal M_0^{\tau,+}(z)=e^{\Theta_+^\tau(z)},
  \qquad
  \mathcal M_0^{\tau,-}(z)=e^{\Theta_-^\tau(z)}.
\end{equation}

The reason the quantity $\mathcal M_{n-k}^{\tau,+}(z)+\mathcal M_{n-k}^{\tau,-}(z)-2$ appears below is the following. For a fixed parent type and lag $k$, the numbers of ``$+$''-children and ``$-$''-children are independent Poisson random variables with the \emph{same} mean $q_k^{\tau,\pm}$. Hence the lag-$k$ continuation factor is the product of the two Poisson cluster factors:
\begin{align*}
  &\exp\!\left(q_k^{\tau,\pm}(\mathcal M_{n-k}^{\tau,+}(z)-1)\right)\,
  \exp\!\left(q_k^{\tau,\pm}(\mathcal M_{n-k}^{\tau,-}(z)-1)\right)\\
  &=
  \exp\!\left(q_k^{\tau,\pm}(\mathcal M_{n-k}^{\tau,+}(z)+\mathcal M_{n-k}^{\tau,-}(z)-2)\right).
\end{align*}
This is exactly the aggregated form used in the proposition below.

The next proposition gives the exact cluster recursion in this aggregated representation and the exact transform of the finite-$\tau$ INAR approximation.

\begin{proposition}[Exact exponential discrete Volterra recursion]\label{prop:exact-transform-weak}
For each $\tau\ge1$ and each $z\in\mathcal D_\tau(T)$, let $q_k^{\tau,+}$ and $q_k^{\tau,-}$ be the parent-type coefficients from \eqref{eq:q-pm-exact}, corresponding respectively to a ``$+$''-parent and a ``$-$''-parent. Then the cluster transforms $\mathcal M_n^{\tau,+}(z)$ and $\mathcal M_n^{\tau,-}(z)$ satisfy, for every $n\ge1$,
\begin{align}
  \mathcal M_n^{\tau,+}(z)
  &=\exp\!\left(
      \Theta_+^\tau(z)
      +\sum_{k=1}^n q_k^{\tau,+}
       \left(\mathcal M_{n-k}^{\tau,+}(z)+\mathcal M_{n-k}^{\tau,-}(z)-2\right)
    \right), \label{eq:Mplus-exact}
  \\
  \mathcal M_n^{\tau,-}(z)
  &=\exp\!\left(
      \Theta_-^\tau(z)
      +\sum_{k=1}^n q_k^{\tau,-}
       \left(\mathcal M_{n-k}^{\tau,+}(z)+\mathcal M_{n-k}^{\tau,-}(z)-2\right)
    \right). \label{eq:Mminus-exact}
\end{align}
If we define
\begin{equation}\label{eq:Gn-exact}
  G_n^\tau(z):=\mathcal M_n^{\tau,+}(z)+\mathcal M_n^{\tau,-}(z)-2,
\end{equation}
then $G_n^\tau(z)$ satisfies the closed recursion
\begin{equation}\label{eq:G-rec-exact}
  G_n^\tau(z)
  =\exp\!\left(
      \Theta_+^\tau(z)+\sum_{k=1}^n q_k^{\tau,+}G_{n-k}^\tau(z)
    \right)
   +\exp\!\left(
      \Theta_-^\tau(z)+\sum_{k=1}^n q_k^{\tau,-}G_{n-k}^\tau(z)
    \right)-2.
\end{equation}
Moreover, the exact transform of the INAR approximation is
\begin{equation}\label{eq:Phi-exact}
  \Phi^\tau(z,T):=\E\!\left[e^{zP_T^\tau}\right]
  =\exp\!\left(\sum_{m=1}^{N_\tau}\hat\mu_\tau(m)\,G_{N_\tau-m}^\tau(z)\right).
\end{equation}
\end{proposition}

\begin{proof}
{The bivariate INAR($\infty$) process admits a cluster representation: each immigrant generates a Poisson-branching cluster, and the log-price transform factors over these clusters. Tracking the contribution of a single cluster of type $+$ or $-$ over $n$ residual steps leads directly to the recursion \eqref{eq:Mplus-exact}--\eqref{eq:Mminus-exact}. Summing over the independent immigrants whose baseline-intensity weights are $\hat\mu_\tau(m)$ yields the outer exponential formula \eqref{eq:Phi-exact}.} Full details are in Appendix~\ref{supp:weak-error-proofs}.
\end{proof}

%%%%%%%%%%%%%%%%%%

\subsection{Second-order reduction on bounded frequency windows}
\label{subsec:second-order-reduction}

We now reduce the exact exponential recursion
\eqref{eq:G-rec-exact} to a second-order discrete recursion on each bounded vertical strip. Fix $\eta\in\mathbb R$ and $R>0$, and define
\[
  \Gamma_{\eta,R}:=\{\eta+\ii \xi:\ |\xi|\le R\},
  \qquad
  M_{\eta,R}:=\sup_{z\in\Gamma_{\eta,R}}|z|=\sqrt{\eta^2+R^2}.
\]
For the remainder of this subsection, all constants may depend on $(\eta,R,T)$, but not on $(n,\tau,z)$.
Throughout, we assume that $\Gamma_{\eta,R}$ is contained in the common admissible strip $\mathcal S_T$ from Assumption~\ref{ass:strip}; in particular, $\Gamma_{\eta,R}\subset\mathcal D_\tau(T)$ for all sufficiently large~$\tau$, so that all moment-generating function manipulations below are valid on this segment.

Recalling the exponents appearing in \eqref{eq:G-rec-exact}, define
\begin{equation}\label{eq:AB-def-new33}
  A_n^\tau(z)
  :=\Theta_+^\tau(z)+\sum_{k=1}^n q_k^{\tau,+}G_{n-k}^\tau(z),
  \qquad
  B_n^\tau(z)
  :=\Theta_-^\tau(z)+\sum_{k=1}^n q_k^{\tau,-}G_{n-k}^\tau(z),
\end{equation}
where
\[
  \Theta_+^\tau(z)=z(c_\tau-d_\tau),
  \qquad
  \Theta_-^\tau(z)=-zc_\tau,
\]
and
\[
  q_k^{\tau,+}=\frac{\varphi_k^\tau}{1+\beta},
  \qquad
  q_k^{\tau,-}=\frac{\beta\varphi_k^\tau}{1+\beta}.
\]
By definition,
\begin{equation}\label{eq:G-from-AB-new33}
  G_n^\tau(z)=e^{A_n^\tau(z)}+e^{B_n^\tau(z)}-2.
\end{equation}

It is convenient to introduce the symmetric and antisymmetric combinations
\begin{equation}\label{eq:SD-def-new33}
  S_n^\tau(z):=A_n^\tau(z)+B_n^\tau(z),
  \qquad
  D_n^\tau(z):=A_n^\tau(z)-B_n^\tau(z).
\end{equation}
Using
\[
  q_k^{\tau,+}+q_k^{\tau,-}=\varphi_k^\tau,
  \qquad
  q_k^{\tau,+}-q_k^{\tau,-}=\frac{1-\beta}{1+\beta}\,\varphi_k^\tau,
\]
together with
\[
  \Theta_+^\tau(z)+\Theta_-^\tau(z)=-zd_\tau,
  \qquad
  \Theta_+^\tau(z)-\Theta_-^\tau(z)=z(2c_\tau-d_\tau),
\]
we obtain the exact recursions
\begin{align}
  S_n^\tau(z)
  &= -zd_\tau+\sum_{k=1}^n \varphi_k^\tau\,G_{n-k}^\tau(z),
  \label{eq:S-rec-new33}
  \\
  D_n^\tau(z)
  &= z(2c_\tau-d_\tau)
     +\frac{1-\beta}{1+\beta}\sum_{k=1}^n \varphi_k^\tau\,G_{n-k}^\tau(z).
  \label{eq:D-rec-new33}
\end{align}
Since
\[
  A_n^\tau(z)=\frac{S_n^\tau(z)+D_n^\tau(z)}{2},
  \qquad
  B_n^\tau(z)=\frac{S_n^\tau(z)-D_n^\tau(z)}{2},
\]
it follows from \eqref{eq:G-from-AB-new33} that
\begin{equation}\label{eq:G-SD-new33}
  G_n^\tau(z)
  =
  e^{(S_n^\tau(z)+D_n^\tau(z))/2}
  +
  e^{(S_n^\tau(z)-D_n^\tau(z))/2}
  -2.
\end{equation}

With $\varepsilon_\tau:=1-a_\tau=\gamma\tau^{-\alpha}$, the first step is a
bounded-strip bootstrap on the exact recursion:
\[
  \sup_{0\le n\le N_\tau}\sup_{z\in\Gamma_{\eta,R}}
  \left(
    |G_n^\tau(z)|+|A_n^\tau(z)|+|B_n^\tau(z)|+|S_n^\tau(z)|+|D_n^\tau(z)|
  \right)
  \le C_{\eta,R,T}\,\varepsilon_\tau,
\]
for all sufficiently large $\tau$; this follows from a Gronwall-type argument on
\eqref{eq:G-rec-exact}, and full details are given in Appendix~\ref{supp:weak-error-proofs}.

\begin{proposition}[Second-order reduced discrete recursion]
\label{prop:reduced-recursion-new33}
For every fixed $\eta\in\mathbb R$ and $R>0$, uniformly for $0\le n\le N_\tau$ and
$z\in\Gamma_{\eta,R}$,
\begin{equation}\label{eq:G-reduced-new33}
\begin{aligned}
  G_n^\tau(z)
  &=
    \sum_{k=1}^n \varphi_k^\tau\,G_{n-k}^\tau(z)
    -zd_\tau
    + c_\tau^2 z^2 \\
  &\quad
    + \frac{1-\beta}{1+\beta}\,c_\tau z
      \sum_{k=1}^n \varphi_k^\tau\,G_{n-k}^\tau(z)
    + \frac{1+\beta^2}{2(1+\beta)^2}
      \left(\sum_{k=1}^n \varphi_k^\tau\,G_{n-k}^\tau(z)\right)^2
    + \mathcal R_n^\tau(z),
\end{aligned}
\end{equation}
where
\begin{equation}\label{eq:R-reduced-new33}
  \sup_{0\le n\le N_\tau}\sup_{z\in\Gamma_{\eta,R}}
  |\mathcal R_n^\tau(z)|
  \le C_{\eta,R,T}\,\varepsilon_\tau^3.
\end{equation}
\end{proposition}

\begin{proof}
{Using the bounded-strip smallness estimate above, expand $e^{A_n^\tau}+e^{B_n^\tau}-2$ to second order in the small exponents $A_n^\tau$ and $B_n^\tau$. The coefficient identities \eqref{eq:source-match}--\eqref{eq:quadratic-match} then identify the source, linear, and quadratic terms as in \eqref{eq:G-reduced-new33}, with the third-order remainder bounded by $C\varepsilon_\tau^3$.} Full details are in Appendix~\ref{supp:weak-error-proofs}.
\end{proof}

Proposition~\ref{prop:reduced-recursion-new33} is the point where the nonlinear discrete dynamics acquires the same source, linear, and quadratic structure as the fractional Riccati equation. From here on, the argument is no longer probabilistic: it becomes a deterministic comparison between a discrete renewal equation and its continuous Volterra counterpart.

%%%%%%%%%%%%%%%%%%%%%%%%%%%%%%%%%%%%%%%%%%%%%%%%

\subsection{Rescaling and match with the continuous fractional Riccati equation}
\label{subsec:rescaling-match}

Next, we want to show that, after a deterministic rescaling, the reduced discrete recursion matches the continuous Riccati coefficients up to a local truncation error.

We now compare the second-order reduced discrete recursion
\eqref{eq:G-reduced-new33} with the continuous fractional Riccati equation
\eqref{eq:riccati-volterra-weak}. The main point is that, after a suitable deterministic rescaling, the reduced discrete recursion has the same leading-order source, linear, and quadratic coefficients as the continuous rough Heston Riccati equation.

Recall that
\[
  \varepsilon_\tau:=1-a_\tau=\gamma\tau^{-\alpha},
\]
and define the rescaled discrete variable
\begin{equation}\label{eq:H-rescaled-def}
  H_n^\tau(z):=\frac{\mu}{\theta\,\varepsilon_\tau}\,G_n^\tau(z),
  \qquad 0\le n\le N_\tau,\quad z\in\Gamma_{\eta,R}.
\end{equation}
Equivalently,
\[
  G_n^\tau(z)=\frac{\theta}{\mu}\,\varepsilon_\tau H_n^\tau(z).
\]

This normalization is chosen so that the source term in
\eqref{eq:G-reduced-new33} becomes exactly of order $\tau^{-\alpha}$ with the same coefficient as in the continuous Riccati equation. Indeed, substituting
$G_{n-k}^\tau(z)=\frac{\theta}{\mu}\,\varepsilon_\tau H_{n-k}^\tau(z)$
into \eqref{eq:G-reduced-new33} and dividing by $(\theta/\mu)\varepsilon_\tau$, we obtain
\begin{equation}\label{eq:Htau-rec-pre}
\begin{aligned}
  H_n^\tau(z)
  &= a_\tau\sum_{k=1}^n \varphi_k H_{n-k}^\tau(z)
     +\frac{-zd_\tau+c_\tau^2z^2}{(\theta/\mu)\varepsilon_\tau}
     +\frac{1-\beta}{1+\beta}\,c_\tau z\,
      a_\tau\sum_{k=1}^n \varphi_k H_{n-k}^\tau(z)
\\
  &\qquad
     +\frac{\theta}{\mu}\,
      \frac{1+\beta^2}{2(1+\beta)^2}\,
      \varepsilon_\tau
      \left(a_\tau\sum_{k=1}^n \varphi_k H_{n-k}^\tau(z)\right)^2
     +\widetilde{\mathcal R}_n^\tau(z),
\end{aligned}
\end{equation}
where
\begin{equation}\label{eq:Rtilde-def}
  \widetilde{\mathcal R}_n^\tau(z)
  :=\frac{\mu}{\theta\,\varepsilon_\tau}\mathcal R_n^\tau(z).
\end{equation}
By Proposition~\ref{prop:reduced-recursion-new33},
\begin{equation}\label{eq:Rtilde-bound}
  \sup_{0\le n\le N_\tau}\sup_{z\in\Gamma_{\eta,R}}
  |\widetilde{\mathcal R}_n^\tau(z)|
  \le C_{\eta,R,T}\,\varepsilon_\tau^2.
\end{equation}

We now identify the coefficients in \eqref{eq:Htau-rec-pre} term by term.

\paragraph{Source term}
Using
$c_\tau^2=\frac{\theta\varepsilon_\tau}{2\mu\tau^\alpha}$
and
$d_\tau=\frac{\theta\varepsilon_\tau}{2\mu\tau^\alpha}$,
we get the exact identity
\begin{equation}\label{eq:source-match}
  \frac{-zd_\tau+c_\tau^2z^2}{(\theta/\mu)\varepsilon_\tau}
  =\frac{1}{2\tau^\alpha}(z^2-z).
\end{equation}

\paragraph{Linear term}
Using
\[
  c_\tau=\sqrt{\frac{\theta\gamma}{2\mu}}\,\tau^{-\alpha},
\]
the discrete linear coefficient is
$\frac{1-\beta}{1+\beta}\sqrt{\frac{\theta\gamma}{2\mu}}\,z\,\tau^{-\alpha}$.
On the other hand, by \eqref{correlation} and \eqref{eq:mu-nu-map},
\[
  \rho=\frac{1-\beta}{\sqrt{2(1+\beta^2)}},
  \qquad
  \nu=\sqrt{\frac{\theta(1+\beta^2)}{\gamma\mu(1+\beta)^2}},
\]
and hence
\begin{equation}\label{eq:rho-nu-id}
  \gamma\rho\nu
  =
  \frac{1-\beta}{1+\beta}\sqrt{\frac{\gamma\theta}{2\mu}}.
\end{equation}
Therefore,
\begin{equation}\label{eq:linear-match}
  \frac{1-\beta}{1+\beta}\,c_\tau z
  =
  \gamma\rho\nu\,z\,\tau^{-\alpha}.
\end{equation}

\paragraph{Quadratic term}
Again using \eqref{eq:mu-nu-map} and $\varepsilon_\tau=\gamma\tau^{-\alpha}$,
\begin{equation}\label{eq:quadratic-match}
  \frac{\theta}{\mu}\frac{1+\beta^2}{2(1+\beta)^2}\,\varepsilon_\tau
  =
  \frac{(\gamma\nu)^2}{2}\tau^{-\alpha}.
\end{equation}

Combining \eqref{eq:Htau-rec-pre}--\eqref{eq:quadratic-match}, and using
$a_\tau=1-\gamma\tau^{-\alpha}$, we obtain the following form of the
reduced discrete recursion:
\begin{equation}\label{eq:Htau-rec-main}
\begin{aligned}
  H_n^\tau(z)
  &= a_\tau\sum_{k=1}^n \varphi_k H_{n-k}^\tau(z)
     +\frac{1}{2\tau^\alpha}(z^2-z)
\\
  &\quad
     +\gamma\rho\nu\,z\,\tau^{-\alpha}
      \sum_{k=1}^n \varphi_k H_{n-k}^\tau(z)
     +\frac{(\gamma\nu)^2}{2}\tau^{-\alpha}
      \left(\sum_{k=1}^n \varphi_k H_{n-k}^\tau(z)\right)^2
     +E_n^\tau(z),
\end{aligned}
\end{equation}
where the remainder $E_n^\tau(z)$ collects the difference between
\eqref{eq:Htau-rec-pre} and the leading-order approximation above. In particular, by \eqref{eq:Rtilde-bound}, the exact coefficient identities above, and the exact relation $a_\tau=1-\gamma\tau^{-\alpha}$,
\begin{equation}\label{eq:En-bound}
  \sup_{0\le n\le N_\tau}\sup_{z\in\Gamma_{\eta,R}}
  |E_n^\tau(z)|
  \le C_{\eta,R,T}\tau^{-2\alpha}.
\end{equation}

Equation \eqref{eq:Htau-rec-main} already shows that the reduced discrete recursion has the same source, linear, and quadratic coefficients as the continuous fractional Riccati equation, up to a local truncation error of order $\tau^{-2\alpha}$. The remaining task is to identify the discrete Volterra operator and compare it quantitatively with the continuous fractional integral in \eqref{eq:riccati-volterra-weak}.

To do so, we next rewrite \eqref{eq:Htau-rec-main} in renewal form. Let
$\psi^\tau:=\sum_{m\ge1}(\varphi^\tau)^{*m}$
be the discrete renewal kernel associated with $\varphi^\tau=a_\tau\varphi$, and define
\[
  \mathsf{S}_k^\tau:=\frac{\tau(1-a_\tau)}{a_\tau}\psi_k^\tau,
  \qquad k\ge1.
\]
A key technical input is that the measures associated with $\mathsf{S}^\tau$ converge weakly to the Mittag--Leffler resolvent density
\[
  f_{\alpha,\gamma}(t)=\gamma t^{\alpha-1}E_{\alpha,\alpha}(-\gamma t^\alpha);
\]
see Appendix~\ref{supp:scaling-proofs}.
This is exactly the resolvent kernel that arises in the continuous rough Heston Riccati equation.

%%%%%%%%%%%%%%%%%%%%%%%%%%%%%%%%%%%%%%%%%%%%%%%%

\subsection{Renewal inversion and Laplace-domain comparison mechanisms}
\label{sec:renewal-comparison}

The goal of this subsection is to rewrite the reduced discrete recursion in
resolvent form and to identify the comparison mechanism between the discrete and
continuous Volterra structures at the level of cumulative resolvents. In this
subsection we isolate the exact discrete renewal identities, the pointwise
renewal bounds, and the Laplace-domain consistency estimates that identify the
mechanism behind the comparison. 

Introduce the quadratic polynomial
\begin{equation}\label{eq:Fz-def}
  F_z(x):=\frac12(z^2-z)+\gamma\rho\nu z\,x+\frac{(\gamma\nu)^2}{2}x^2,
  \qquad x\in\C.
\end{equation}
Then \eqref{eq:Htau-rec-main} can be rewritten as
\begin{equation}\label{eq:Htau-rec-main-again}
  H_n^\tau(z)
  =
  a_\tau\sum_{k=1}^n\varphi_k H_{n-k}^\tau(z)
  +\tau^{-\alpha}F_z\!\left(\sum_{k=1}^n\varphi_k H_{n-k}^\tau(z)\right)
  +E_n^\tau(z),
\end{equation}
where
\begin{equation}\label{eq:En-bound-again}
  \sup_{0\le n\le N_\tau}\sup_{z\in\Gamma_{\eta,R}}
  |E_n^\tau(z)|
  \le C_{\eta,R,T}\tau^{-2\alpha}.
\end{equation}

Using the reduced recursion together with the bounded-strip smallness, one may replace the averaged state variable by $H_n^\tau$ itself up to a lower-order error; see Appendix~\ref{supp:weak-error-proofs}. Inverting the discrete renewal operator then rewrites the reduced recursion in discrete resolvent form, while the continuous fractional Riccati equation admits an equivalent resolvent form with kernel $f_{\alpha,\gamma}$; see Appendix~\ref{supp:weak-error-proofs}.

We now introduce the cumulative resolvents associated with the discrete and
continuous renewal structures.

\begin{definition}[Discrete and continuous cumulative resolvents]
\label{def:cum-resolvents}
The discrete cumulative resolvent is defined by
\[
  \mathcal G_n^\tau
  :=
  \frac{1-a_\tau}{\gamma a_\tau}
  \sum_{k=1}^n\psi_k^\tau,
  \qquad n\ge0,
\]
with $\mathcal G_0^\tau=0$.
The continuous cumulative resolvent associated with $f_{\alpha,\gamma}$ is
defined by
\[
  \mathcal G(t)
  :=
  \frac1\gamma\left(1-E_\alpha(-\gamma t^\alpha)\right),
  \qquad t\ge0.
\]
Equivalently,
$\mathcal G(t)=\frac1\gamma\int_0^t f_{\alpha,\gamma}(s)\,ds$.
\end{definition}

A direct generating-function calculation gives the closed discrete renewal
equation
\begin{equation}\label{eq:cum-resolvent-renewal-mainpaper}
  \mathcal G_n^\tau
  =
  a_\tau\sum_{k=1}^n\varphi_k\,\mathcal G_{n-k}^\tau
  +\tau^{-\alpha}\sum_{k=1}^n\varphi_k,
  \qquad n\ge1.
\end{equation}
If $\mathcal G^\tau(t):=\mathcal G_n^\tau$ for
$t\in[n/\tau,(n+1)/\tau)$, then its Laplace transform is
\begin{equation}\label{eq:cum-resolvent-laplace-mainpaper}
  \widehat{\mathcal G^\tau}(z)
  =
  \frac{1-a_\tau}{\gamma z}\,
  \frac{\hat\varphi(z/\tau)}{1-a_\tau\hat\varphi(z/\tau)},
  \qquad z>0,
\end{equation}
while the continuous cumulative resolvent satisfies
\begin{equation}\label{eq:cum-resolvent-laplace-cont-mainpaper}
  \widehat{\mathcal G}(z)=\frac{1}{z(z^\alpha+\gamma)},
  \qquad z>0.
\end{equation}
These formulas make the comparison mechanism explicit. Inserting
Lemma~\ref{lem:phi-expansion} into
\eqref{eq:cum-resolvent-laplace-mainpaper} yields the symbol-level estimate
\[
  \widehat{\mathcal G^\tau}(z)-\widehat{\mathcal G}(z)
  = \mathcal O\!\left(c_1(\alpha)\tau^{-(1-\alpha)}\right),
\]
uniformly on compact $z$-sets, so the first discrete-to-continuous correction
vanishes as $\alpha\uparrow1^-$. The derivation is carried out in Appendix~\ref{supp:weak-error-proofs}.

%%%%%%%%%%%%%%%%%%

\subsection{Riccati comparison on bounded strips}
\label{subsec:transform-bound}

For fixed $(\eta,R)$, all constants may depend on $(\eta,R,T)$, but not on $(n,\tau,z)$. On each bounded vertical strip where the continuous fractional Riccati solution remains bounded, the continuous solution enjoys the time-regularity estimates needed below, and the corresponding regularity bounds for the Riccati nonlinearity $g_z(t)=F_z(h(t,z))$ are standard; see Appendix~\ref{supp:weak-error-proofs}. Parameter-specific bounded-strip Riccati checks for the numerical experiments are reported in Subsection~\ref{subsec:numerical-assumption-checks}.

The weakly singular quadrature constant generating the coefficient $c_1(\alpha)$, the corresponding quadrature-defect decomposition for $\mathcal G'$, and a pointwise renewal bound for the discrete resolvent increments are likewise proved in Appendix~\ref{supp:weak-error-proofs}. These are the inputs that turn the symbol-level correction into the following uniform time-domain comparison of the cumulative resolvents.

\begin{proposition}[Uniform cumulative-resolvent comparison]
\label{prop:cum-resolvent-uniform-comparison}
Assume $\alpha\in(1/2,1)$ and let
\[
  \mathcal G_n^\tau
  :=
  \frac{1-a_\tau}{\gamma a_\tau}\sum_{k=1}^n\psi_k^\tau,
\qquad
\text{and}
\qquad
  \mathcal G(t)
  :=
  \frac1\gamma\left(1-E_\alpha(-\gamma t^\alpha)\right).
\]
Then there exists a coefficient $C_{\mathrm{res}}(\alpha)\ge0$ and a constant $C_T>0$, independent of $\tau$, such that
\begin{equation}\label{eq:cum-resolvent-uniform-comparison}
  \sup_{0\le n\le N_\tau}
  |\mathcal G_n^\tau-\mathcal G(t_n)|
  \le
  C_T\tau^{-\alpha}+C_{\mathrm{res}}(\alpha)\,\tau^{-(1-\alpha)},
\end{equation}
with $C_{\mathrm{res}}(\alpha)\to0$ as $\alpha\uparrow1^-$. 
\end{proposition}

\begin{proof}
{The proof starts from the explicit formulas
\eqref{eq:cum-resolvent-laplace-mainpaper} and
\eqref{eq:cum-resolvent-laplace-cont-mainpaper}. In the Laplace domain,
$\widehat{\mathcal{G}}^\tau(z)\approx\widehat{\mathcal{G}}(z)$ up to a
correction proportional to $c_1(\alpha)(z/\tau)^{1-\alpha}$. Inverting this
Laplace-domain correction in discrete time produces the
$C_{\mathrm{res}}(\alpha)\tau^{-(1-\alpha)}$ term, with
$C_{\mathrm{res}}(\alpha)\propto c_1(\alpha)\to0$. The $C_T\tau^{-\alpha}$ term
comes from the regular Riemann sum approximation of $\mathcal{G}$.} Full
details are in Appendix~\ref{supp:weak-error-proofs}.
\end{proof}

We now introduce the discrete and continuous resolvent operators acting on a test
function $g$:
\[
  (\mathcal Q_\tau g)_n
  :=
  \frac{1}{\gamma\tau}\sum_{j=1}^{n-1}\mathsf{S}_{n-j}^\tau g(t_j),
  \qquad 1\le n\le N_\tau,
\]
and
\[
  (\mathcal Q g)(t_n)
  :=
  \frac1\gamma\int_0^{t_n}f_{\alpha,\gamma}(t_n-s)g(s)\,ds.
\]

\begin{proposition}[Transfer from cumulative-resolvent comparison to operator comparison]
\label{prop:resolvent-transfer-from-cumulative}
Let $g:[0,T]\to\C$ be absolutely continuous and satisfy
\begin{align}
  |g(t)-g(0)| &\le C_g t^\alpha,
  &&0\le t\le T,\label{eq:g-assump-1}\\
  |\partial_t g(t)| &\le C_g t^{\alpha-1},
  &&0<t\le T.\label{eq:g-assump-2}
\end{align}
Then there exists a coefficient $C_{\mathrm{op}}(\alpha)\ge0$ and a constant
$C_T'>0$, depending only on $(\alpha,\gamma,T)$, such that
\begin{equation}\label{eq:resolvent-operator-compare-main}
  \sup_{1\le n\le N_\tau}
  |(\mathcal Q_\tau g)_n-(\mathcal Q g)(t_n)|
  \le
  C_T'(|g(0)|+C_g)\tau^{-\alpha}
  +
  C_{\mathrm{op}}(\alpha)(|g(0)|+C_g)\tau^{-(1-\alpha)},
\end{equation}
where $C_{\mathrm{op}}(\alpha)\to0$ as $\alpha\uparrow1^-$.
\end{proposition}

\begin{proof}
See Appendix~\ref{supp:weak-error-proofs}.
\end{proof}

Applying Proposition~\ref{prop:resolvent-transfer-from-cumulative} to the
continuous Riccati nonlinearity yields the nonlinear resolvent estimate needed
below. The final comparison between $H^\tau$ and $h$ is then closed with a
discrete weakly singular Gr\"onwall estimate; see Appendix~\ref{supp:weak-error-proofs}.

\begin{proposition}[Bounded-strip Riccati comparison]
\label{prop:Volterra-rate-fixed}
{Fix $\eta\in\R$ and $R>0$. Under Assumptions~\ref{ass:strip} and~\ref{ass:riccati-bound},} there exists a coefficient $C_{\mathrm{Ric}}(\alpha)\ge0$ and a constant
$C_{\eta,R,T}>0$ such that
\begin{equation}\label{eq:Volterra-rate-fixed}
  \sup_{0\le n\le N_\tau}\sup_{z\in\Gamma_{\eta,R}}
  \left|H_n^\tau(z)-h(t_n,z)\right|
  \le
  C_{\eta,R,T}\tau^{-\alpha}+C_{\mathrm{Ric}}(\alpha)\tau^{-(1-\alpha)},
\end{equation}
with $C_{\mathrm{Ric}}(\alpha)\to0$ as $\alpha\uparrow1^-$.
\end{proposition}

\begin{proof}
{The argument has four steps; full technical details are in Appendix~\ref{supp:weak-error-proofs}.

\medskip
\noindent\textit{Step 1: Reformulation in resolvent form.}
The averaged-state replacement estimate in Appendix~\ref{supp:weak-error-proofs} shows that
$|H_n^\tau - \sum_{k=1}^n\varphi_k H_{n-k}^\tau|\le C\tau^{-\alpha}$, so the averaged
state variable can be replaced by $H_n^\tau$ itself up to a $\tau^{-2\alpha}$
error. The discrete and continuous resolvent reformulations in Appendix~\ref{supp:weak-error-proofs} then
rewrite the reduced recursion and the continuous Riccati equation in terms of
the cumulative resolvents $\mathcal{G}_n^\tau$ and $\mathcal{G}(t)$ from
Definition~\ref{def:cum-resolvents}.

\medskip
\noindent\textit{Step 2: Cumulative-resolvent comparison (origin of the two branches).}
Proposition~\ref{prop:cum-resolvent-uniform-comparison} provides the estimate
\[
  \sup_{0\le n\le N_\tau}|\mathcal{G}_n^\tau-\mathcal{G}(t_n)|
  \le C_T\tau^{-\alpha}+C_{\mathrm{res}}(\alpha)\tau^{-(1-\alpha)}.
\]
The $\tau^{-\alpha}$ branch originates from the regular-Riemann-sum error in the baseline approximation. The $C_{\mathrm{res}}(\alpha)\tau^{-(1-\alpha)}$ branch originates from the symbol-level correction $c_1(\alpha)$ identified in Lemma~\ref{lem:phi-expansion}: in the Laplace domain, $(1-\hat\varphi(s))^{-1}$ deviates from the continuous Mittag-Leffler kernel by $c_1(\alpha)s^{1-\alpha}$, which in time domain produces the $\tau^{-(1-\alpha)}$ rate. Since $c_1(\alpha)\to0$ as $\alpha\uparrow1^-$, so does $C_{\mathrm{res}}(\alpha)$.

\medskip
\noindent\textit{Step 3: Transfer to the resolvent operators.}
Proposition~\ref{prop:resolvent-transfer-from-cumulative} transfers the
cumulative-resolvent comparison to the associated operators $\mathcal{Q}_\tau g$
and $\mathcal{Q}g$. Under Assumption~\ref{ass:riccati-bound}, the continuous
Riccati solution $h(\cdot,z)$ is uniformly $C^\alpha([0,T])$ on
$\Gamma_{\eta,R}$. Hence the operator bound applies to the nonlinearity
$g_z(t)=F_z(h(t,z))$.

\medskip
\noindent\textit{Step 4: Closing by discrete weakly singular Gr\"{o}nwall.}
Setting $e_n^\tau:=H_n^\tau(z)-h(t_n,z)$, the preceding steps yield a discrete
integral inequality of the form
$|e_n^\tau|\le C\tau^{-\alpha}+C_{\mathrm{res}}(\alpha)\tau^{-(1-\alpha)}+C_{\eta,R,T}\sum_{j<n}\Delta\mathcal{G}_{n-j}^\tau|e_j^\tau|$.
The discrete weakly singular Gronwall estimate in Appendix~\ref{supp:weak-error-proofs} closes this
inequality and propagates the two-branch rate through the entire comparison.}
\end{proof}

We now transfer the bounded-strip Riccati comparison from
Proposition~\ref{prop:Volterra-rate-fixed} to the exponent and transform levels.
The first transfer step is the logarithmic-transform estimate:

\begin{lemma}[Exponent comparison on bounded strips]
\label{lem:exponent-comparison}
For every fixed $\eta\in\R$ such that the vertical line $\mathrm{Re}\ z=\eta$ lies in
the admissible strip from Assumption~\ref{ass:strip}, and every $R>0$, there
exists a coefficient $C_{\log}(\alpha)\ge0$ and a constant $C_{\eta,R,T}>0$
such that, for all sufficiently large $\tau$,
\begin{equation}\label{eq:log-transform-bound}
  \sup_{z\in\Gamma_{\eta,R}}
  |\log\Phi^\tau(z,T)-\log\phi(z,T)|
  \le
  C_{\eta,R,T}\tau^{-\alpha}
  +
  C_{\log}(\alpha)\tau^{-(1-\alpha)},
\end{equation}
with $C_{\log}(\alpha)\to0$ as $\alpha\uparrow1^-$.
\end{lemma}

\begin{proof}
See Appendix~\ref{supp:weak-error-proofs}.
\end{proof}

The transform comparison below then follows by exponentiation. The slower
branch remains inherited from the combination of the weighted baseline
quadrature, the kernel-symbol correction, and the ensuing
discrete/continuous resolvent comparison.

\begin{proposition}[Bounded-strip transform comparison]
\label{prop:transform-bound-fixed}
{For every fixed $\eta\in\R$ such that the vertical line $\mathrm{Re}\,z=\eta$ lies in the admissible strip from Assumption~\ref{ass:strip}, and every $R>0$ with $\Gamma_{\eta,R}\subset\mathcal S_T$, under Assumptions~\ref{ass:strip} and~\ref{ass:riccati-bound},} there exists a coefficient $C_{\Phi}(\alpha)\ge0$ and a constant $C_{\eta,R,T}>0$ such that, for all sufficiently large
$\tau$,
\begin{equation}\label{eq:transform-bound-fixed}
  \sup_{z\in\Gamma_{\eta,R}}
  \left|
    \Phi^\tau(z,T)-\phi(z,T)
  \right|
  \le
  C_{\eta,R,T}\tau^{-\alpha}+C_{\Phi}(\alpha)\tau^{-(1-\alpha)},
\end{equation}
Equivalently,
\begin{equation}\label{eq:transform-bound-xi}
  \sup_{|\xi|\le R}
  \left|
    \Phi^\tau(\eta+\ii \xi,T)-\phi(\eta+\ii \xi,T)
  \right|
  \le
  C_{\eta,R,T}\tau^{-\alpha}+C_{\Phi}(\alpha)\tau^{-(1-\alpha)},
\end{equation}
with $C_{\Phi}(\alpha)\to0$ as $\alpha\uparrow1^-$.
\end{proposition}

\begin{proof}
{By Lemma~\ref{lem:exponent-comparison},
$|\log\Phi^\tau(z,T)-\log\phi(z,T)|\le C\tau^{-\alpha}+C_{\log}(\alpha)\tau^{-(1-\alpha)}$
uniformly on $\Gamma_{\eta,R}$. Because both logarithmic transforms are
uniformly bounded on $\Gamma_{\eta,R}$ for large $\tau$, exponentiating
transfers the estimate directly to the transforms $|\Phi^\tau-\phi|$; one may
therefore take $C_{\Phi}(\alpha):=C_{\log}(\alpha)$.} Full
details are in Appendix~\ref{supp:weak-error-proofs}.
\end{proof}

%%%%%%%%%%%%%%%%%%%%%%%%%%%%%%%%

\subsection{Truncated Carr--Madan pricing error}
\label{subsec:pricing-error}

In this subsection, we transfer the bounded-strip transform comparison from
Proposition~\ref{prop:transform-bound-fixed} to the corresponding truncated
Carr--Madan pricing functional. This yields a bounded-frequency weak-error bound
for the low-frequency contribution to the Carr--Madan representation of the
European call price.

Recall that
\[
  P_T^\tau=\log(S_T^\tau/S_0),
  \qquad
  P_T=\log(S_T/S_0),
\]
and write
\[
  \Phi^\tau(z,T):=\E\left[e^{zP_T^\tau}\right],
  \qquad
  \phi(z,T):=\E\left[e^{zP_T}\right].
\]
Let
\[
  k:=\log(K/S_0)
\]
denote the log-moneyness.

Fix a damping parameter $\eta>1$ such that the vertical line
$\mathrm{Re}\,z=\eta$ lies in the admissible strip from
Assumption~\ref{ass:strip}. For $R>0$, define the truncated Carr--Madan
functionals
\begin{align}
  \frac{C_R^\tau(T,K)}{S_0}
  &:=
  \frac{e^{-(\eta-1)k}}{2\pi}
  \int_{-R}^{R}
    e^{-\ii \xi k}
    \frac{\Phi^\tau(\eta+\ii \xi,T)}
         {(\eta+\ii \xi)(\eta+\ii \xi-1)}
  \,d\xi,
  \label{eq:CRtau-def-trunc}
  \\
  \frac{C_R(T,K)}{S_0}
  &:=
  \frac{e^{-(\eta-1)k}}{2\pi}
  \int_{-R}^{R}
    e^{-\ii \xi k}
    \frac{\phi(\eta+\ii \xi,T)}
         {(\eta+\ii \xi)(\eta+\ii \xi-1)}
  \,d\xi.
  \label{eq:CR-def-trunc}
\end{align}
These are the low-frequency truncations of the usual Carr--Madan inversion
formula for the discrete and continuous models, respectively.

\begin{remark}
The quantities $C_R^\tau(T,K)$ and $C_R(T,K)$ are not the full call prices.
They are the prices obtained by truncating the Carr--Madan Fourier inversion to
the bounded frequency window $|\xi|\le R$. Passing from $C_R^\tau,C_R$ to the
full prices $C^\tau,C$ requires a separate high-frequency tail analysis, which
is outside the scope of the present bounded-strip comparison.
\end{remark}

We now state the direct pricing consequence of the bounded-strip transform
comparison.

\begin{proposition}[Truncated pricing error on bounded frequency windows]
\label{prop:pricing-local-trunc}
Fix $\eta>1$ in the admissible strip from Assumption~\ref{ass:strip} and let $R>0$.
Then, for all sufficiently large $\tau$,
\[
  \left|C_R^\tau(T,K)-C_R(T,K)\right|
  \le
  C_1\tau^{-\alpha}+C_2(\alpha)\tau^{-(1-\alpha)},
\]
for some constants $C_1>0$ and $C_2(\alpha)\ge0$, independent of $\tau$, with
$C_2(\alpha)\to0$ as $\alpha\uparrow1^-$.
\end{proposition}

\begin{proof}
{By Proposition~\ref{prop:transform-bound-fixed}, for all $z=\eta+\ii\xi$ with $|\xi|\le R$,
\[
  |\Phi^\tau(\eta+\ii\xi,T)-\phi(\eta+\ii\xi,T)|
  \le C_{\eta,R,T}\tau^{-\alpha}+C_\Phi(\alpha)\tau^{-(1-\alpha)}.
\]
The Carr--Madan kernel $(\eta+\ii\xi)^{-1}(\eta+\ii\xi-1)^{-1}$ is bounded in magnitude on $|\xi|\le R$ by a constant $\kappa_{\eta,R}>0$ (since $\eta>1$ keeps the denominator away from zero). Inserting the pointwise bound into the definitions \eqref{eq:CRtau-def-trunc}--\eqref{eq:CR-def-trunc} and integrating over $[-R,R]$ gives
\begin{align*}
  |C_R^\tau(T,K)-C_R(T,K)|
  &\le
  \frac{e^{-(\eta-1)k}}{2\pi}\int_{-R}^{R}
    \frac{|\Phi^\tau(\eta+\ii\xi,T)-\phi(\eta+\ii\xi,T)|}
         {|(\eta+\ii\xi)(\eta+\ii\xi-1)|}
  \,d\xi \\
  &\le \frac{R\,\kappa_{\eta,R}\,e^{-(\eta-1)k}}{\pi}
   \left(C_{\eta,R,T}\tau^{-\alpha}+C_\Phi(\alpha)\tau^{-(1-\alpha)}\right),
\end{align*}
which is exactly the stated bound with $C_1=\frac{R\,\kappa_{\eta,R}\,e^{-(\eta-1)k}}{\pi}C_{\eta,R,T}$ and $C_2(\alpha)=\frac{R\,\kappa_{\eta,R}\,e^{-(\eta-1)k}}{\pi}C_\Phi(\alpha)$. Since $C_\Phi(\alpha)\to0$ as $\alpha\uparrow1^-$, so does $C_2(\alpha)$.} The proof is provided in Appendix~\ref{supp:weak-error-proofs}.
\end{proof}

A uniform version over compact moneyness sets is recorded in Appendix~\ref{app:proof-pricing-local-trunc-uniform}.

\begin{remark}
The bounded-window weak error therefore inherits the same two-branch structure
\[
  \tau^{-\alpha}
  \qquad\text{and}\qquad
  C_2(\alpha)\,\tau^{-(1-\alpha)},
  \qquad C_2(\alpha)\to0\ \text{as }\alpha\uparrow1^-.
\]
The second term is the slower-decaying branch. Analytically, it is generated by the weakly singular baseline quadrature, the kernel-symbol correction, and the resulting discrete/continuous cumulative-resolvent comparison.

This is a low-frequency pricing result on a fixed bounded frequency window. A
full weak-error estimate for European call prices would require, in addition, a
quantitative control of the high-frequency tails in the Carr--Madan inversion
formula and a rule coupling the truncation level $R$ with the discretization
parameter $\tau$.
\end{remark}

%%%%%%%%%%%%%%%%

\section{Numerical Results}\label{sec:numerical}

Section~\ref{sec:weak-euro} establishes weak-error bounds only for the truncated Carr--Madan functional on fixed bounded frequency windows. The numerical experiments below serve two distinct purposes. First, for European pricing they illustrate the practical behavior of the microstructural approximation in parameter regimes where Fourier-based benchmarks are available. Second, for Asian, lookback, and barrier contracts they demonstrate the broader applicability of the simulator, without claiming an error theory of the same strength as in the bounded-window European analysis. Accordingly, the European comparisons should be read as complementary to specialized weak schemes rather than as a claim of uniform superiority.

\subsection{Payoff conventions and implementation workflow}
\label{subsec:numerical-workflow}

\paragraph{Pricing framework and risk-neutral measure}
We work under the risk-neutral measure $\mathbb Q$ associated with the money-market num\'eraire $B_t^{0}=\exp\left(\int_0^t r_udu\right)$ where $(r_t)_{t\in[0,T]}$ represents the interest rate process.
Let the discounted price be $\tilde S_t := S_t/B_t^0$. For notational convenience, we keep writing $S_t$ for $\tilde S_t$.
Under this convention, $S$ is a $\mathbb Q$-martingale and evolves as
\[
dS_t = S_t \sqrt{V_t} dW_t^X,
\]
while $V_t$ follows the rough Heston dynamics specified below, and $d\left\langle W^X,W^V\right\rangle_t=\rho dt$.
Equivalently, one may view $S_t$ as the discounted price process under the corresponding pricing measure.

In the numerics, the cumulative INAR($\infty$) model is used as a simulation device for the limiting risk-neutral dynamics. Option prices are computed by simulating the approximating paths and evaluating the corresponding payoffs. For simplicity and consistency with the cited literature, we take the risk-free rate to be $r=0$, so no discount factor appears in the reported prices. The same implementation is used across European, Asian, lookback, and barrier contracts.

We consider the following four payoff classes:
\begin{enumerate}
    \item \textit{European options}. The payoff depends only on the terminal price:
    \[
        H_T = \begin{cases}
        \max(S_T - K, 0) & \text{for a call option},\\
        \max(K - S_T, 0) & \text{for a put option}.
        \end{cases}
    \]

    \item \textit{Arithmetic Asian options}. The payoff depends on the arithmetic average over the simulation grid:
    \[
        H_T = \begin{cases}
        \max(\bar{S}_T - K, 0) & \text{for a call option},\\
        \max(K-\bar{S}_T, 0) & \text{for a put option},
        \end{cases}
    \]
    where
    \begin{equation}\label{Asian:defn}
    \bar{S}_T = \frac{1}{M+1}\sum_{i=0}^{M} S_{i\Delta t},\qquad\text{with }\Delta t = T/M.
    \end{equation}
    In the literature one often studies the continuous approximation
    $\frac{1}{T}\int_0^T S_u\,du$. In simulation we work directly with the
    discrete average \eqref{Asian:defn}, including the initial point.

    \item \textit{Lookback options}. The payoff depends on the running maximum
    $M_T := \max_{0\leq t\leq T} S_t$ or running minimum
    $m_T := \min_{0\leq t\leq T} S_t$:
    \[
        H_T = \begin{cases}
        \max(M_T-K,0) & \text{for a lookback call option},\\
        \max(K-m_T,0) & \text{for a lookback put option}.
        \end{cases}
    \]

    \item \textit{Barrier options}. The payoff depends on whether the path hits
    a barrier level $B$ before maturity. For an up-and-in barrier call,
    \[
        H_T = \begin{cases}
        \max(S_T - K, 0) & \text{if } \max_{0\leq t\leq T} S_t \geq B,\\
        0 & \text{otherwise},
        \end{cases}
    \]
    while for a down-and-out barrier put,
    \[
        H_T = \begin{cases}
        \max(K - S_T, 0) & \text{if } \min_{0\leq t\leq T} S_t > B,\\
        0 & \text{otherwise}.
        \end{cases}
    \]
\end{enumerate}

We simulate the \emph{INAR($\infty$)} recursion on $[0,T]$ with zero pre-history ($X_n\equiv0$ for $n\le0$) via a non-circular linear convolution (CDQ--FFT). At time $n\le\lfloor\tau T\rfloor$, the update only involves lags $1,\dots,n-1$. Hence, for any $p\ge n-1$ (in particular $p\ge\lfloor\tau T\rfloor$), the INAR($p$) recursion coincides \emph{pathwise} with our INAR($\infty$) update on $[0,T]$. In other words, our implementation does not fix $p$; effectively it uses all available lags at each step, and therefore no extra ``truncation'' error from $p\to\infty$ arises on the finite horizon.

In practice one works with a finite value of $\tau$. The effect of this discretization level on numerical accuracy is examined separately in the convergence analysis below; at the implementation level, the key point is that the CDQ--FFT scheme reduces the per-path cost from the naive $\mathcal{O}(\tau^2)$ to $\mathcal{O}(\tau\log^2\tau)$ while retaining numerical stability.

The implementation revolves around two basic ingredients. First, we use a
helper kernel
\[
f_{\alpha}(t) =
\begin{cases}
0 & \text{if } t \leq 0, \\
1 - \frac{1}{\Gamma(1-\alpha)} & \text{if } t = 1, \\
\frac{1}{\Gamma(1-\alpha)} \left(\frac{1}{(t-1)^{\alpha}} - \frac{1}{t^{\alpha}}\right) & \text{if } t > 1,
\end{cases}
\]
to build the discrete heavy-tailed weights. Second, the path generator records,
on the fly, the functionals needed for the different option classes: terminal
price for European payoffs, running averages for Asian payoffs, running extrema
for lookback payoffs, and barrier-hit indicators for barrier payoffs.

\subsubsection{Enhanced parallel Monte Carlo implementation}

We employ a parallel Monte Carlo simulation that distributes the computational
workload across CPU cores. Before giving the algorithm, we first make explicit
the core $\mathcal{O}(\tau^2)$ update that CDQ accelerates. At each integer time
step $n$, the two processes $X_n^+$ and $X_n^-$ are sampled independently from
a Poisson distribution with intensity
\begin{equation}\label{eq:lambda-update}
  \lambda_n = \hat{\mu}_\tau(n) + \sum_{k=1}^{n-1} w_k Y_{n-k},
  \qquad
  Y_n := X_n^+ + \beta X_n^-.
\end{equation}
The factor $\frac{1}{1+\beta}$ in the kernel
$w_k = \frac{a_\tau}{1+\beta} f_\alpha(k)$ arises from collapsing the bivariate
$(X^+,X^-)$ system to an effective univariate intensity: since
$\mathbb{E}[Y_n \mid \mathcal{F}_{n-1}] = (1+\beta)\lambda_n$, each historical
value $Y_{n-k}$ must be divided by $(1+\beta)$ to recover the contribution to
$\lambda_n$. Computing \eqref{eq:lambda-update} sequentially costs
$\mathcal{O}(n)$ per step and $\mathcal{O}(\tau^2)$ per path. The CDQ scheme
reduces this to $\mathcal{O}(\tau\log^2\tau)$ via divide-and-conquer and FFT
convolution. All path functionals are computed on the simulated grid over
$[0,T]$.

{% Algorithm 1 is placed here, inside Section 4.1.1, as a properly anchored float.
}\begin{algorithm}[ht]
\caption{Parallel INAR($\infty$) simulation with CDQ--FFT convolution}\label{algo}
\SetAlgoLined
\SetKwFunction{CDQFunc}{CDQ}
\SetKwProg{Fn}{Function}{}{end}
\small
\KwIn{Model parameters $(\alpha, \gamma, \rho, \nu, \theta, V_0, S_0)$ and simulation parameters $(\tau, n_{\text{sims}})$}
\KwOut{Price estimates and 95\% confidence intervals for European, Asian, lookback, and barrier options}

\textit{Preprocessing:}\\
Derive parameters $\beta \leftarrow \text{solve\_beta}(\rho)$, $\mu \leftarrow \text{solve\_mu}(\nu, \theta, \beta, \gamma)$, and $\xi_0 \leftarrow V_0/\theta$\\
Derive INAR($\infty$) time-scaled parameters $a_\tau, \mu_\tau$ from $(\gamma, \mu, \tau, \alpha)$\\
Pre-compute the fractional kernel $w_k = \frac{a_\tau}{1+\beta} f_{\alpha}(k)$ for $k=1, \dots, \tau$\\
Pre-compute the baseline intensity $\hat{\mu}_\tau(t)$ for $t=1, \dots, \tau$\\
Pre-compute the price-path coefficients
$c_\tau \leftarrow \sqrt{\dfrac{\theta(1-a_\tau)}{2\mu\tau^\alpha}}$ and
$d_\tau \leftarrow \dfrac{\theta(1-a_\tau)}{2\mu\tau^\alpha}$\\
Identify the CDQ block lengths required by the recursion, and pre-compute the corresponding FFT plans and transformed kernel blocks\\
Build the fixed CDQ operation list (leaf/direct/FFT blocks)

\textit{Parallel simulation:}\\
Determine the number of available CPU threads $N_{\text{threads}}$\\
Partition $n_{\text{sims}}$ simulations among threads\\
\For{each thread in parallel}{
    Initialize a thread-specific random number generator and Poisson sampler\\
    Allocate once the reusable work arrays $\lambda^{\mathrm{hist}}$, $Y$, and the FFT buffer, together with thread-local sums and sums of squares for all payoffs\\
    \For{each simulation assigned to the thread}{
        Reset the reusable workspace and initialize the running summaries $N^\pm$, $S_{\mathrm{avg}}$, $S_{\max}$, $S_{\min}$, and the barrier flags\\
        \For{each pre-computed CDQ operation in execution order}{
            \eIf{the operation is a leaf at time $t$}{
                $\lambda \leftarrow \max\{0,\hat{\mu}_\tau(t) + \lambda^{\mathrm{hist}}_t\}$\\
                Sample $X_t^+, X_t^- \sim \mathrm{Poisson}(\lambda)$ independently\\
                Set $Y_t \leftarrow X_t^+ + \beta X_t^-$\\
                Update the running totals $N_t^\pm \leftarrow N_{t-1}^\pm + X_t^\pm$ and the price\\
                $S_t \leftarrow S_0\exp\!\left(c_\tau(N_t^+ - N_t^-) - d_\tau N_t^+\right)$\\
                Update the running average, extrema, and barrier indicators
            }{
                \eIf{the block is small}{
                    Propagate the left-to-right convolution directly
                }{
                    Reuse the pre-computed FFT plan and transformed kernel block to convolve the left data with the kernel and add the contribution to the right half of $\lambda^{\mathrm{hist}}$
                }
            }
        }
        Compute payoffs for all option types from the path summaries\\
        Update the thread-local payoff sums and sums of squares
    }
}

\textit{Postprocessing:}\\
Aggregate the thread-local sums and sums of squares across all threads\\
For each option type, compute the mean and 95\% confidence interval from the aggregated statistics\\
\Return{A structure containing all option price estimates and their confidence intervals}
\end{algorithm}

\subsection{Parameter-specific checks of the analytical assumptions}
\label{subsec:numerical-assumption-checks}

The weak-error analysis is stated under Assumption~\ref{ass:strip} and under
the bounded-strip Riccati regularity condition~\ref{ass:riccati-bound}. Since
these assumptions are not proved in full
generality in the present paper, we record here the parameter-specific checks
used for the numerical experiments. The purpose of this subsection is precisely
to show the reader that the reported experiments are run in parameter regimes
for which the relevant assumptions are numerically supported. These checks do
not constitute general proofs.

\paragraph{Admissible-strip check}
For the rough-parameter sets used in the numerical section, we numerically
checked the line $\mathrm{Re}\ z=2$ at maturity $T=1$. On the continuous side, solving
\eqref{eq:riccati-cont-weak} at $z=2$ gave finite values, with
\[
  h(1,2)\approx 0.999
\]
for the parameter set
\[
  (\alpha,\gamma,\rho,\nu,\theta,V_0)=(0.62,0.1,-0.681,0.331,0.3156,0.0392),
\]
and
\[
  h(1,2)\approx 0.688,\ 0.697,\ 0.715,\ 0.721
\]
for the slice-comparison parameter set
\[
  (\gamma,\rho,\nu,\theta,V_0,T)=(0.3,-0.7,1,0.02/0.3,0.02,1),
\]
at $\alpha=0.55,0.62,0.80,0.95$. On the discrete side, evaluating
\eqref{eq:G-rec-exact}--\eqref{eq:Phi-exact} at $z=2$ yielded finite transforms
for the discretization levels used in the experiments:
\[
  \log \Phi^\tau(2,1)\approx 0.05593,\ 0.05574,\ 0.05564,\ 0.05560
\]
for $\tau=40,80,160,320$ in the first parameter set, and
\[
  \log \Phi^{160}(2,1)\approx 0.02321,\ 0.02307,\ 0.02276,\ 0.02254
\]
for $\alpha=0.55,0.62,0.80,0.95$ in the slice-comparison set. These
computations provide parameter-specific evidence that the line $\mathrm{Re}\ z=2$ is
admissible for the reported experiments, but they are not a general proof of
Assumption~\ref{ass:strip}.

\paragraph{Bounded-strip Riccati check}
{The following computations verify Assumption~\ref{ass:riccati-bound} (the bounded-strip Riccati regularity condition stated in Section~\ref{sec:weak-euro}) for the parameter sets used in the experiments.}
For the same parameter sets, we numerically solved the fractional Riccati
equation on a grid of times $t\in[0,T]$ and frequencies
$z=\eta+\ii\xi$ with $\eta=2$, $|\xi|\le 20$, and $T=1$, and recorded
$\max_{0\le t\le T}|h(t,z)|$. For the rough-parameter set
\[
  (\alpha,\gamma,\rho,\nu,\theta,V_0)=(0.62,0.1,-0.681,0.331,0.3156,0.0392),
\]
we obtained
\[
  \max_{|\xi|\le 20}\max_{0\le t\le 1}|h(t,2+\ii\xi)| \approx 184.2933,
\]
while for the slice-comparison parameter set
\[
  (\gamma,\rho,\nu,\theta,V_0,T)=(0.3,-0.7,1,0.02/0.3,0.02,1),
\]
we obtained
\[
  \max_{|\xi|\le 20}\max_{0\le t\le 1}|h(t,2+\ii\xi)|
  \approx 53.3162,\ 54.1052,\ 56.4927,\ 59.0652
\]
for $\alpha=0.55,0.62,0.80,0.95$. These computations support the bounded-strip
Riccati regularity condition for the reported parameter sets only and do not
replace a general proof.

%%%%%%%%%%%%%%%%%%%%%%%%%%%%%%%%%%%%%%%%%%%%%%%%

\subsection{Rough case \texorpdfstring{$(\alpha=0.62)$}{(alpha=0.62)}}

We now turn to the rough specification used for the main pricing experiments.
The same FFT-accelerated INAR($\infty$) simulator from
Subsection~\ref{subsec:numerical-workflow} is used for all payoff classes on
$[0,T]$. Here we focus on the numerical evidence most directly tied to the
analytical and modeling claims of the paper.

\begin{table}[ht]
\centering
\caption{Model parameters for the rough-case numerical experiments.}
\small
\setlength{\tabcolsep}{6pt}
\begin{tabular}{llll}
\toprule
Parameter & Value & Type & Derivation \\
\midrule
$\alpha$ & 0.62 & Direct & - \\
$\gamma$ & 0.1 & Direct & - \\
$\rho$ & $-0.681$ & Direct & - \\
$\nu$ & 0.331 & Direct & vol-of-vol in the $V$-equation \\
$\theta$ & 0.3156 & Direct & - \\
$V_0$ & 0.0392 & Direct & - \\
$S_0$ & 100 & Direct & - \\
\midrule
$\beta$ & 27.5583 & Derived & From $\displaystyle \rho=\frac{1-\beta}{\sqrt{2(1+\beta^2)}}$ \\
$\mu$ & 26.8592 & Derived & From $\displaystyle \mu=\frac{\theta(1+\beta^2)}{\gamma\nu^2(1+\beta)^2}$ \\
$\xi_0$ & 0.124208 & Derived & From $\displaystyle V_0=\xi_0\theta$ \\
\bottomrule
\end{tabular}
\label{tab:all_params}
\end{table}

We now turn to the empirical effect of the discretization level $\tau$. For a European call option with strike $K=100$, Table~\ref{tab:tau_convergence} reports the estimated price, its 95\% confidence interval, and the corresponding CPU time for increasing values of $\tau$.

Table~\ref{tab:tau_convergence} suggests numerical stabilization and close numerical agreement with the benchmark:
as $\tau$ increases from $40$ to $320$, the estimates lie within narrow $95\%$
confidence intervals around the benchmark and stabilize. {While the point estimates are not monotone in $\tau$, the benchmark deviation is generally smaller at finer discretizations in this experiment, with the smallest deviation attained at $\tau=320$.}

Balancing accuracy and cost, we therefore use $\tau=320$ as a practical default:
further increases yield improvements that are within sampling error while the
computational time grows noticeably. We keep $\tau=320$ for the remaining experiments
unless otherwise stated.

\subsubsection{Convergence analysis and optimal discretization}
\label{sec:convergence_analysis}

We now examine the practical effect of choosing a finite discretization level $\tau$. To investigate the empirical accuracy-cost trade-off, we conducted a series of convergence tests for a European call option with strike $K=100$, using an increasing number of time steps from $\tau=40$ to $\tau=320$.

\begin{table}[ht]
    \centering
    \caption{Convergence of the European call option price ($K=100$) with respect to the discretization parameter $\tau$. All simulations were run with 500,000 paths.}
    \label{tab:tau_convergence}
    \small
    \begin{tabularx}{\linewidth}{@{}>{\centering\arraybackslash}p{0.17\linewidth}>{\centering\arraybackslash}X>{\centering\arraybackslash}X>{\centering\arraybackslash}X@{}}
        \toprule
        \textit{Parameter $\tau$} & \textit{\shortstack{Estimated\\ Price}} & \textit{\shortstack{Deviation from\\ Benchmark}} & \textit{\shortstack{Computation Time\\ (seconds)}}\\
        \midrule
        40    & 9.4844 & $+0.0107$ & 0.54s\\
        80    & 9.4923 & $+0.0186$ & 1.28s\\
        160   & 9.4827 & $+0.0090$ & 3.30s\\ 
        320   & 9.4722 & $-0.0015$ & 7.63s\\
        \midrule
        \multicolumn{3}{@{}l}{\textit{Benchmark (\cite{callegaro2021fast})}} & \textit{9.4737} \\
        \bottomrule
    \end{tabularx}
\end{table}

Table~\ref{tab:all_options} reports a representative European benchmark under the rough specification.

\begin{table}[thb]
\centering
\caption{Representative European option prices under the rough Heston model ($\alpha=0.62$) using the INAR($\infty$) FFT-based simulator. All simulations use $\tau=320$ and $500,000$ paths. The values in parentheses are benchmark prices from \cite{callegaro2021fast}.}
\small
\setlength{\tabcolsep}{5pt}
\begin{tabular*}{\textwidth}{@{\extracolsep{\fill}}ccccc@{}}
\toprule
Strike & Call & 95\% CI & Put & 95\% CI \\
\midrule
80  & 22.1558(22.1366) & [22.0972, 22.2145] & 2.1141 & [2.0989, 2.1294] \\
90  & 14.9839(14.9672) & [14.9324, 15.0355] & 4.9422 & [4.9177, 4.9668] \\
100 & 9.4722(9.4737)  & [9.4470, 9.5326]  & 9.4481 & [9.4136, 9.4825] \\
110 & 5.6407(5.6234)  & [5.6070, 5.6744]  & 15.5990 & [15.5553, 15.6427] \\
120 & 3.1584(3.1424)  & [3.1331, 3.1837]  & 23.1167 & [23.0654, 23.1680] \\
\bottomrule
\end{tabular*}
\label{tab:all_options}
\end{table}

For European option pricing, the values in parentheses in Table~\ref{tab:all_options} are taken from \cite{callegaro2021fast}. As Table~\ref{tab:all_options} shows, the INAR($\infty$) estimates stay close to that Fourier benchmark across strikes.

%\subsubsection{Computational Performance}
On an Apple M4-based desktop, the optimized implementation completes the representative rough-volatility experiment in a few seconds while retaining tight confidence intervals.

\subsubsection{Comparison with existing \texorpdfstring{$\alpha<1$}{alpha<1} schemes on implied-volatility slices}
\label{sec:iv-slice-comparison}

We next consider a fixed-maturity implied-volatility slice and compare our INAR($\infty$)-based simulator with three representative alternatives from the recent literature in the regime $\alpha<1$. More precisely, we use the rough Heston parameter set
\[
\gamma=0.3,\ \theta=\frac{0.02}{0.3},\ \nu=1,\ \rho=-0.7,\ V_0=0.02,\ S_0=100,\ T=1,
\]
and the log-moneyness grid $k=\log(K/S_0)\in\{-0.5,-0.4,\ldots,0.5\}$. We compare the FFT-accelerated INAR($\infty$) scheme with a hybrid quadratic-exponential (HQE) weak scheme from the affine-forward-variance literature \cite{Gatheral2022,BayerBreneis2024WeakRoughHeston}, a multifactor/integrated Euler approximation in the spirit of \cite{AlfonsiKebaier2024}, and the inverse-Gaussian iVi scheme of \cite{AbiJaberAttal2025}. For all methods, we use $2^{24}$ Monte Carlo paths and a common time discretization of $160$ steps; for the multifactor approximation, we use $3$ factors. The reference slice is computed by Fourier inversion of the rough Heston characteristic function.

The resulting smile overlays are shown in Figure~\ref{fig:iv_slice_panel_alpha}. In each panel, the overall smile shape is captured well, but the slice-level accuracy differs materially across methods and roughness levels. The corresponding maximum absolute slice errors and wall-clock times are reported in Table~\ref{tab:iv_slice_panel_alpha}.

\begin{figure}[ht]
    \centering
    \includegraphics[width=0.92\textwidth]{./smile_overlay_panel_t_1_paper_smile.png}
    \caption{Implied-volatility slice comparison for $\alpha\in\{0.55,0.62,0.80,0.95\}$ at maturity $T=1$. In each panel, the black curve is the Fourier reference smile, the red curve is the INAR($\infty$)-CDQ-FFT slice, the blue and green curves are the iVi and HQE weak-scheme benchmarks, and the brown curve is the $3$-factor integrated multifactor Euler approximation. All methods use $2^{24}$ Monte Carlo paths.}
    \label{fig:iv_slice_panel_alpha}
\end{figure}

\begin{table}[ht]
\centering
\caption{Slice-level comparison for the multi-$\alpha$ experiment in Figure~\ref{fig:iv_slice_panel_alpha}. Each entry reports the maximum absolute IV-slice error on the 11-point log-moneyness grid, followed by the wall-clock time in seconds.}
\label{tab:iv_slice_panel_alpha}
\small
\setlength{\tabcolsep}{5pt}
\begin{tabular}{ccccc}
\toprule
$\alpha$ & Integrated (3 factors) & INAR-CDQ-FFT & iVi & HQE \\
\midrule
0.55 & $0.0229$ ($78.9$s) & $0.0080$ ($190.9$s) & $0.0014$ ($64.6$s) & $0.0006$ ($56.8$s) \\
0.62 & $0.0185$ ($28.5$s) & $0.0053$ ($85.0$s)  & $0.0014$ ($61.6$s) & $0.0006$ ($68.1$s) \\
0.80 & $0.0078$ ($27.1$s) & $0.0020$ ($111.7$s) & $0.0013$ ($67.6$s) & $0.0006$ ($76.3$s) \\
0.95 & $0.0550$ ($26.7$s) & $0.0009$ ($110.4$s) & $0.0021$ ($63.1$s) & $0.0010$ ($77.8$s) \\
\bottomrule
\end{tabular}
\end{table}

Several conclusions emerge from Figure~\ref{fig:iv_slice_panel_alpha} and Table~\ref{tab:iv_slice_panel_alpha}. First, in this experiment and at the common step budget used here, the specialized weak schemes HQE and iVi produce the smallest European-slice errors for the rougher cases $\alpha=0.55$ and $\alpha=0.62$, with HQE delivering the smallest error among the methods tested. Second, under this parameter set, the INAR($\infty$) approximation becomes more competitive as $\alpha$ increases: at $\alpha=0.80$ its slice error is close to the iVi/HQE benchmarks, and at $\alpha=0.95$ its maximum IV-slice error is slightly smaller than that of HQE and smaller than that of iVi. This is consistent with the refined weak-error decomposition, in which the coefficient of the discrete-to-continuous Volterra branch vanishes as $\alpha\uparrow1^-$. Under the present proof, one obtains at best an $\mathcal O(1-\alpha)$ vanishing rate, so that for moderate $\tau$ the $\tau^{-\alpha}$ branch may dominate the observed error. Third, the coarse three-factor integrated approximation is visibly less accurate in this experiment, especially near $\alpha=1$, which is consistent with a deterioration of the underlying low-factor kernel approximation in that regime. More precisely, when $\alpha=0.95$ (so $H=\alpha-\tfrac12=0.45$), the Volterra kernel is already close to the constant Heston kernel, while the present three-factor approximation still represents it only through a small number of strictly decaying exponential modes; in the absence of a zero-decay mode, this low-factor kernel approximation can become noticeably less accurate in that near-Heston regime.

These slice experiments therefore clarify the numerical position of the proposed scheme. On European implied-volatility slices, the INAR($\infty$) simulator is not uniformly the most efficient competitor against current specialized weak schemes. However, its slice accuracy improves substantially as the roughness parameter approaches the classical regime, while still relying on the same microstructural simulator that also handles the path-dependent contracts studied above.

%%%%%%%%%%%%%%%%%%%%%%%%%%%%%%%%%%%%%

\subsection{Non-rough case: Heston model \texorpdfstring{$(\alpha=1)$}{(alpha=1)}}

To validate our implementation, we compare the INAR($\infty$) scheme with the classical Heston benchmark at $\alpha=1$, using the degenerate Markovian kernel from Remark~\ref{rem:alpha-one-kernel}.

%\subsubsection{European Options}
Since the closed-form solution for the European option is available for the classical Heston model, we can directly compare the results from the INAR($\infty$) implementation to this benchmark. The results are given in Table~\ref{tab:european_comparison}. We choose $\tau=320$, $S_0=100$ and $K=80,90,100,110,120$ for all cases. 

\begin{table}[ht]
  \centering
  \caption{Comparison of European option prices between the closed-form solution and the INAR($\infty$) implementation}
  \small
  \setlength{\tabcolsep}{5pt}
  \begin{tabular*}{\textwidth}{@{\extracolsep{\fill}}lcccccc@{}}
  \toprule
  \multirow{2}{*}{Strike} & \multicolumn{2}{c}{Closed-form solution} & \multicolumn{4}{c}{%
  INAR($\infty$)} \\
  \cmidrule(lr){2-3}\cmidrule(lr){4-7}
   & Call & Put & Call & 95\% CI & Put & 95\% CI \\
  \midrule
  80  & 21.8822 & 1.8822 & 21.8834 & [21.8265, 21.9402] & 1.8752 & [1.8612, 1.8893] \\
  90  & 14.6187 & 4.6187 & 14.6202 & [14.5704, 14.6700] & 4.6121 & [4.5889, 4.6353] \\
  100 & 9.0983  & 9.0983  & 9.0942  & [9.0531, 9.1354]  & 9.0861  & [9.0530, 9.1192] \\
  110 & 5.2883  & 15.2883 & 5.2811  & [5.2490, 5.3131]  & 15.2729 & [15.2306, 15.3152] \\
  120 & 2.8849  & 22.8849 & 2.8807  & [2.8569, 2.9045]  & 22.8725 & [22.8227, 22.9223] \\
  \bottomrule
  \end{tabular*}
  \label{tab:european_comparison}
\end{table}

\subsubsection{Classical-Heston benchmark details}

The classical Heston benchmark is simulated on a discrete time grid with time step
$\Delta t = T/\tau = 1/320$. For each path, the asset price $S$ and variance $V$
are discretized as
\begin{align*}
    S_{i+1} &= S_i + r S_i \Delta t + \sqrt{\max(V_i, 0)} S_i \sqrt{\Delta t} Z_1, \\
    V_{i+1} &= \max\left(0, V_i + \gamma(\theta - V_i)\Delta t + \nu \sqrt{\max(V_i, 0)} \sqrt{\Delta t} Z_2\right),
\end{align*}
where $i=0, \dots, \tau-1$, and $(Z_1, Z_2)$ are correlated standard normal
random variables with correlation $\rho$. This corresponds to a diffusion-truncation
scheme with terminal projection: we use $\sqrt{\max(V_i,0)}$ in the diffusion
term and finally set $V_{i+1}\leftarrow \max\{V_{i+1},0\}$. To reduce Monte Carlo
variance, we employ antithetic variates. For each pair $(Z_1,Z_2)$, we simulate
a second path with $(-Z_1,-Z_2)$ and use the average payoff of the two paths.
For efficiency, the benchmark accumulates path summaries (terminal value,
running average, running extrema, and barrier-hit indicators) on the fly and
prices all strikes/payoffs in a single Monte Carlo pass. All benchmark prices
reported in the following use 500,000 paths, matching the INAR($\infty$) runs.

\begin{table}[ht]
  \centering
  \caption{Comparison of Asian option prices between the Euler--Maruyama benchmark and the INAR($\infty$) implementation.}
  \small
  \setlength{\tabcolsep}{5pt}
  \begin{tabular*}{\textwidth}{@{\extracolsep{\fill}}lcccccc@{}}
  \toprule
  \multirow{2}{*}{Strike} & \multicolumn{2}{c}{Euler--Maruyama} & \multicolumn{4}{c}{INAR($\infty$)} \\
  \cmidrule(lr){2-3}\cmidrule(lr){4-7}
   & Call & Put & Call & 95\% CI & Put & 95\% CI \\
  \midrule
  80  & 20.1541 & 0.1557 & 20.1648 & [20.1312, 20.1983] & 0.1570 & [0.1541, 0.1598] \\
  90  & 11.2797 & 1.2813 & 11.2907 & [11.2611, 11.3202] & 1.2829 & [1.2737, 1.2920] \\
  100 & 4.9135  & 4.9151  & 4.9280  & [4.9066, 4.9494]  & 4.9202  & [4.9016, 4.9388] \\
  110 & 1.6174  & 11.6190 & 1.6293  & [1.6167, 1.6418]  & 11.6215 & [11.5944, 11.6485] \\
  120 & 0.4044  & 20.4059 & 0.4129  & [0.4068, 0.4190]  & 20.4051 & [20.3733, 20.4369] \\
  \bottomrule
  \end{tabular*}
  \label{tab:asian_comparison}
\end{table}

\begin{table}[ht]
  \centering
  \caption{Comparison of lookback option prices between the Euler--Maruyama benchmark and the INAR($\infty$) implementation.}
  \small
  \setlength{\tabcolsep}{5pt}
  \begin{tabular*}{\textwidth}{@{\extracolsep{\fill}}lcccccc@{}}
  \toprule
  \multirow{2}{*}{Strike} & \multicolumn{2}{c}{Euler--Maruyama} & \multicolumn{4}{c}{INAR($\infty$)} \\
  \cmidrule(lr){2-3}\cmidrule(lr){4-7}
   & Call & Put & Call & 95\% CI & Put & 95\% CI \\
  \midrule
  80  & 38.4117 & 3.3814 & 38.4427 & [38.3980, 38.4874] & 3.3807 & [3.3630, 3.3984] \\
  90  & 28.4117 & 8.3496 & 28.4427 & [28.3980, 28.4874] & 8.3549 & [8.3279, 8.3819] \\
  100 & 18.4117 & 16.5131 & 18.4427 & [18.3980, 18.4874] & 16.5222 & [16.4900, 16.5544] \\
  110 & 10.4984 & 26.5131 & 10.5300 & [10.4903, 10.5697] & 26.5222 & [26.4900, 26.5544] \\
  120 & 5.6284  & 36.5131 & 5.6523  & [5.6211, 5.6834]  & 36.5222 & [36.4900, 36.5544] \\
  \bottomrule
  \end{tabular*}
  \label{tab:lookback_comparison}
\end{table}

\begin{table}[htb]
  \centering
  \caption{Comparison of barrier option prices between the Euler--Maruyama benchmark and the INAR($\infty$) implementation (up-barrier $=110$, down-barrier $=90$).}
  \small
  \setlength{\tabcolsep}{5pt}
  \begin{tabular*}{\textwidth}{@{\extracolsep{\fill}}lccc@{}}
  \toprule
  \multirow{2}{*}{Strike} & \multicolumn{1}{c}{Euler--Maruyama} & \multicolumn{2}{c}{INAR($\infty$)} \\
  \cmidrule(lr){2-2}\cmidrule(lr){3-4}
   & Up-In Call $\left(B=110\right)$ & Up-In Call $\left(B=110\right)$ & 95\% CI \\
  \midrule
  80  & 19.4200 & 19.4397 & [19.3788, 19.5006] \\
  90  & 13.8070 & 13.8234 & [13.7725, 13.8743] \\
  100 & 8.9723  & 8.9829  & [8.9416, 9.0241] \\
  110 & 5.2734  & 5.2811  & [5.2490, 5.3131] \\
  120 & 2.8737  & 2.8807  & [2.8569, 2.9045] \\
  \bottomrule
  \end{tabular*}
  \vspace{0.3cm}

  \begin{tabular*}{\textwidth}{@{\extracolsep{\fill}}lccc@{}}
  \toprule
  \multirow{2}{*}{Strike} & \multicolumn{1}{c}{Euler--Maruyama} & \multicolumn{2}{c}{INAR($\infty$)} \\
  \cmidrule(lr){2-2}\cmidrule(lr){3-4}
   & Down-Out Put $\left(B=90\right)$ & Down-Out Put $\left(B=90\right)$ & 95\% CI \\
  \midrule
  80  & 0.0000 & 0.0000 & [0.0000, 0.0000] \\
  90  & 0.0000 & 0.0000 & [0.0000, 0.0000] \\
  100 & 0.1495  & 0.1510  & [0.1485, 0.1535] \\
  110 & 0.9262  & 0.9308  & [0.9222, 0.9393] \\
  120 & 2.5193  & 2.5223  & [2.5054, 2.5392] \\
  \bottomrule
  \end{tabular*}
  \label{tab:barrier_comparison}
\end{table}

The results indicate that, in this experiment, the optimized INAR($\infty$)
implementation tracks the classical Heston benchmark closely when $\alpha=1$.
European and Asian option prices are numerically close across the two
implementations, and the lookback and barrier prices remain similarly close.

%\subsubsection{Computational Performance}
Thus, in the pure Heston regime the specialized Euler benchmark is faster: for the full option set reported above, the optimized Euler--Maruyama benchmark requires about 2.4 seconds, compared with about 13.1 seconds for the optimized INAR($\infty$) implementation. Beyond this timing comparison, our framework retains the practical advantage of using the same simulator architecture in both rough and classical volatility regimes. This unified treatment is primarily useful when one wants to compare different volatility assumptions within a single implementation framework, rather than to claim a uniform advantage for European pricing in the classical case.

%%%%%%%%%%%%%%%%%%%%%%%%%%%%%%%%%%%%%%%%%%%%%%%%

\subsection{Analysis of the implied volatility surface}

To validate the INAR($\infty$)-based simulator, we examine the implied-volatility (IV) surface produced by our method. A widely discussed feature in the rough-volatility literature is that rough models generate a steeper short-maturity at-the-money (ATM) skew than classical models, often together with an approximate power-law term structure over some maturity ranges. However, the empirical status of a simple power-law ATM skew in equity markets is more nuanced: recent studies show that such a behavior is not uniformly supported across maturities, and rough models need not outperform non-rough or Markovian alternatives when the entire implied-volatility surface is taken into account; see, for example, \cite{guyon2022does,abi2025volatility}. Accordingly, the numerical experiment below should be interpreted as a model-implied diagnostic for the rough Heston specification used in this paper, rather than as evidence that a universal market stylized fact is a pure power law.

Following the simulation setup in Section~\ref{sec:numerical}, we price European calls on a grid of maturities and log-moneyness $k=\log(K/S_0)\in[-0.2,0.2]$. For each $(T,k)$ we invert the Black-Scholes formula to obtain IVs. Figure~\ref{fig:iv_surface} reports the resulting surface for the rough specification ($\alpha=0.62$), computed with the same path count and time step used throughout the study. We consider maturities on the interval $[0,1]$, sampled at monthly grid points.

\begin{figure}[ht]
    \centering
    \includegraphics[width=0.7\textwidth]{ivs.png}
    \caption{Implied-volatility surface produced by the roughness specification ($\alpha=0.62$) with the INAR($\infty$)-FFT simulator. Grid: $T\in\{1/12,2/12,3/12,...,1\}$, $k\in[-0.2,0.2]$. All runs use a fixed time step ($\tau=320$ per year) and $10^6$ paths.}
    \label{fig:iv_surface}
\end{figure}

\subsubsection{ATM-skew diagnostics}

To make this model-implied behavior more explicit, Figure~\ref{fig:atm_skew_panel} (a) traces the finite-difference ATM skew across maturities, while Figure~\ref{fig:atm_skew_panel} (b) shows the corresponding ATM implied volatilities. In our rough Heston specification, the skew steepens as maturity shrinks, whereas the ATM level follows a smoother upward term structure. Fitting a power law to the absolute skew (Figure~\ref{fig:atm_skew_panel}, panel (c)) yields
$|\mathrm{skew}(T)|\approx cT^{-0.369}$ with $R^2\approx0.850$, which is consistent with the short-end pattern generated by the rough specification with $\alpha=0.62$ (so that $H=\alpha-\frac12\approx0.12$). We stress, however, that this is a property of the model under our chosen parameters, not a claim that observed market ATM skews universally follow a pure power law; see \cite{guyon2022does,abi2025volatility} for a related empirical discussion.

\begin{figure}[ht]
    \centering
    \includegraphics[width=0.88\textwidth]{skew.png}
    \caption{ATM diagnostics under the rough specification ($\alpha=0.62$). Panel (a) shows the steepening skew toward short maturities; panel (b) reports the smoother ATM level term structure; panel (c) illustrates the approximate power-law decay $|\mathrm{skew}(T)|\propto T^{H-1/2}$.}
    \label{fig:atm_skew_panel}
\end{figure}

Overall, under the chosen parameter set, the INAR($\infty$) approximation reproduces the familiar short-maturity rough-volatility patterns within a unified simulation framework. We interpret this as a model-implied diagnostic of the rough specification used here, rather than as evidence of a universal empirical law or of uniform superiority over specialized European-pricing schemes.

%%%%%%%%%%%%%%%%%%%%%%%%%%%%%%%%%%%%%%%%%%%%%%%%%%%%%%%%%%%%%%%%%%%

\section{Conclusion}\label{sec:conclusion}

We introduce a bivariate cumulative heavy-tailed INAR($\infty$) framework whose nearly unstable scaling limit is the rough Heston model of \cite{el2019characteristic}, extending to the bivariate price--variance setting the discrete cumulative INAR scaling picture developed in \cite{wang2026scaling}. This provides a discrete-time microstructural interpretation of the joint variance-price dynamics and makes the link between the asymmetry of the microscopic order-flow mechanism and the macroscopic leverage effect explicit through a concrete parameter mapping.

On the pricing side, we derived the exact finite-$\tau$ transform recursion, reduced it in the diffusive regime to a quadratic discrete Volterra equation, and compared it with the continuous rough Heston fractional Riccati equation. The resulting weak-error analysis is local in frequency: under the admissible-strip assumption stated in Section~\ref{sec:weak-euro}, we obtain quantitative bounds for the truncated Carr--Madan pricing functional on bounded frequency windows of the form $C_1\tau^{-\alpha}+C_2(\alpha)\tau^{-(1-\alpha)}$, where $C_2(\alpha)\to0$ as $\alpha\uparrow1^-$. This should be viewed as a rigorous low-frequency pricing result rather than as a complete weak-error theory for full European call prices, since the high-frequency tail of the inversion integral is not analyzed in the present paper.

From a computational perspective, the INAR($\infty$) representation leads to a unified simulation architecture for both European and path-dependent options. The CDQ--FFT implementation reduces the long-memory convolution cost to $\mathcal O(\tau\log^2\tau)$ per path, and the numerical experiments show that the resulting scheme is numerically stable across several payoff classes under the reported parameter sets, remains consistent with the classical Heston benchmark when $\alpha=1$, and reproduces the qualitative implied-volatility features generated by the chosen rough specification. At the same time, the experiments also indicate that on European implied-volatility slices our method should be viewed as complementary to specialized weak schemes rather than as a uniform replacement for them.

Several directions remain open. On the analytical side, it would be desirable to derive the common transform strip directly from the discrete dynamics and to extend the local weak-error result to full Fourier-based European pricing by adding explicit high-frequency tail estimates. On the numerical side, natural next steps include calibration, variance reduction, and statistical inference for the underlying microstructural INAR model, as well as a broader comparison with specialized rough-volatility schemes across different parameter regimes.

%%%%%%%%%%%%%%%%%%%%%%%%%%%%%%%%%%%%%%%%%
\section*{Acknowledgments}

The authors would like to thank the AE and two anonymous referees for helpful comments. Lingjiong Zhu is partially supported by the grants NSF DMS-2053454, NSF DMS-2208303. 

\bibliographystyle{plainnat}
\bibliography{bibtex}

\appendix

\section*{Overview}
These appendices collect the longer proofs and technical bridge results referenced in the main paper.

\section{Proofs for the Scaling-Limit Section}
\label{supp:scaling-proofs}

Throughout Appendices~\ref{supp:scaling-proofs} and \ref{supp:weak-error-proofs}, we use the following asymptotic notation, with the relevant limiting regime specified by the context. We write $f \sim g$ to denote asymptotic equivalence, i.e., $f/g\to 1$ in the stated limit. We use $f \lesssim g$ to indicate that $f \le C g$ for some positive constant $C$ independent of the asymptotic parameter in the stated regime. We write $x \asymp y$ to denote that there exist constants $0<c<C<\infty$
such that $cy \le x \le Cy$ uniformly in the stated regime.

%%%%%%%%%%%%%%%%%%%%%%%%%%%%%%%%%

\subsection{Lemma~\ref{xu2023}}

\begin{lemma}[A $C$-tightness criterion for c\`adl\`ag processes, {\cite[Lemma~3.5]{horst2023convergence}}]
\label{xu2023}
If
$\sup_{n\ge1}\mathbb{E}\left[\left|X_0^{(n)}\right|^{q_0}\right]<\infty$
for some $q_0>0$, then for every fixed horizon $T>0$ the sequence
$\left(X^{(n)}\right)_{n\ge1}$ is $C$-tight on $[0,T]$ if there exists a
constant $\vartheta>2$ such that the following two conditions hold.
\begin{enumerate}[(i)]
  \item
  \[
    \sup_{k=0,1,\ldots,\lfloor Tn^\vartheta\rfloor}\sup_{h\in[0,1/n^\vartheta]}
    \left|\Delta_hX_{k/n^\vartheta}^{(n)}\right|
    \rightarrow 0
    \quad\text{in probability as } n\rightarrow\infty.
  \]
  \item There exist constants $C>0$, $p\ge1$, $m\in\{1,2,\ldots\}$ and pairs
  $\{(a_i,b_i)\}_{i=1}^m$ satisfying
  \[
    a_i\ge0,\qquad b_i>0,\qquad
    \varrho_*:=\min_{1\le i\le m}\{b_i+a_i/\vartheta\}>1,
  \]
  such that for all $n\ge1$ and $h\in(0,1)$,
  \[
    \sup_{t\in[0,T]}\mathbb{E}\left[\left|\Delta_h X_t^{(n)}\right|^p\right]
    \le C\cdot\sum_{i=1}^m\frac{h^{b_i}}{n^{a_i}}.
  \]
\end{enumerate}
\end{lemma}

%%%%%%%%%%%%%%%%%%%%%%%%%%%%%%%%%

\subsection{Lemma~\ref{lem:kernel_convergence}}\label{app:kernel_convergence}

\begin{lemma}[Weak convergence of the renewal kernel]
\label{lem:kernel_convergence}
Let $(\psi_n^\tau)_{n\ge1}$ be the sequence of discrete renewal kernels associated
with $\varphi^\tau=a_\tau\varphi$, and define
\[
  \mathsf S^\tau(t)
  :=
  \frac{1-a_\tau}{a_\tau}\tau\,\psi^\tau_{\lfloor t\tau\rfloor},
  \qquad t\ge0.
\]
Then, as $\tau\to\infty$, the measure $\mathsf S^\tau(t)\,dt$ converges weakly on
$[0,\infty)$ to the probability measure with density
\[
  f_{\alpha,\gamma}(t)
  :=
  \gamma t^{\alpha-1}E_{\alpha,\alpha}(-\gamma t^\alpha),
  \qquad t>0.
\]
\end{lemma}

\begin{proof}

Fix $\alpha\in(\frac12,1)$, $\gamma>0$, and $z>0$. Let $(\varphi_n)_{n\ge1}$ be a probability kernel on $\mathbb N$ satisfying
\[
\sum_{n\ge1}\varphi_n=1,
\qquad
\overline\varphi(n):=\sum_{k>n}\varphi_k
\sim \frac{1}{\Gamma(1-\alpha)}n^{-\alpha}
\qquad (n\to\infty).
\]
For each $\tau>0$, let $a_\tau\in(0,1)$ be such that
\[
\tau^\alpha(1-a_\tau)\to\gamma,
\]
and define
\[
\varphi_n^\tau:=a_\tau\varphi_n,\qquad n\ge1.
\]
For $n\ge1$, let
\[
\psi_n^\tau:=\sum_{k\ge1}(\varphi^\tau)^{*k}_n
=\sum_{k\ge1}a_\tau^k(\varphi^{*k})_n
\]
be the discrete renewal kernel associated with $\varphi^\tau$, and set
$\psi_0^\tau:=0$.
Define the piecewise-constant function
\[
\mathsf S^\tau(t):=\frac{1-a_\tau}{a_\tau}\tau\,\psi^\tau_{\lfloor t\tau\rfloor},
\qquad t\ge0.
\]
Since $\psi_0^\tau=0$, this means $\mathsf S^\tau(t)=0$ on $[0,1/\tau)$.

We first check that $\mathsf S^\tau(t)\,dt$ is a probability measure on $[0,\infty)$. Indeed,
\begin{align*}
\int_0^\infty \mathsf S^\tau(t)\,dt
=
\sum_{n\ge0}\int_{n/\tau}^{(n+1)/\tau}\frac{1-a_\tau}{a_\tau}\tau\,\psi_n^\tau\,dt
=
\frac{1-a_\tau}{a_\tau}\sum_{n\ge0}\psi_n^\tau
=
\frac{1-a_\tau}{a_\tau}\sum_{n\ge1}\psi_n^\tau.
\end{align*}
Since
\[
\sum_{n\ge1}\psi_n^\tau
=
\sum_{k\ge1}\sum_{n\ge1}(\varphi^\tau)^{*k}_n
=
\sum_{k\ge1}\|\varphi^\tau\|_1^k
=
\sum_{k\ge1}a_\tau^k
=
\frac{a_\tau}{1-a_\tau},
\]
it follows that
\[
\int_0^\infty \mathsf S^\tau(t)\,dt=1.
\]

Now consider the Laplace transform
\[
\widehat{\mathsf S^\tau}(z):=\int_0^\infty e^{-zt}\mathsf S^\tau(t)\,dt.
\]
Since $\mathsf S^\tau$ is constant on each interval $[n/\tau,(n+1)/\tau)$, $n\ge0$, we obtain
\begin{align*}
\widehat{\mathsf S^\tau}(z)
&=
\frac{1-a_\tau}{a_\tau}\tau\sum_{n\ge0}\psi_n^\tau
\int_{n/\tau}^{(n+1)/\tau}e^{-zt}\,dt
\\
&=
\frac{1-a_\tau}{a_\tau}\tau\sum_{n\ge1}\psi_n^\tau
\frac{e^{-zn/\tau}-e^{-z(n+1)/\tau}}{z}
=
\frac{1-a_\tau}{a_\tau}\frac{1-e^{-z/\tau}}{z}\tau
\sum_{n\ge1}\psi_n^\tau e^{-zn/\tau}.
\end{align*}

Let
\[
\hat\varphi(s):=\sum_{n\ge1}\varphi_n e^{-sn},
\qquad s>0.
\]
Then
\[
\sum_{n\ge1}\psi_n^\tau e^{-zn/\tau}
=
\sum_{k\ge1}\left(a_\tau\hat\varphi(z/\tau)\right)^k
=
\frac{a_\tau\hat\varphi(z/\tau)}{1-a_\tau\hat\varphi(z/\tau)}.
\]
Therefore,
\begin{equation}\label{eq:laplace-Stau-proof}
\widehat{\mathsf S^\tau}(z)
=
(1-a_\tau)\frac{1-e^{-z/\tau}}{z}\tau\,
\frac{\hat\varphi(z/\tau)}{1-a_\tau\hat\varphi(z/\tau)}.
\end{equation}

Since
\[
\frac{\tau(1-e^{-z/\tau})}{z}\to1
\qquad (\tau\to\infty),
\]
it remains to identify the small-$s$ behavior of $\hat\varphi(s)$ as $s\downarrow0$.

By Abel summation,
\[
\hat\varphi(s)
=
\sum_{n\ge1}\varphi_n e^{-sn}
=
1-(1-e^{-s})\sum_{n\ge0}e^{-sn}\overline\varphi(n).
\]
Using the tail asymptotic
\[
\overline\varphi(n)\sim \frac{1}{\Gamma(1-\alpha)}n^{-\alpha},
\]
together with the standard discrete Karamata--Tauberian asymptotic
\[
\sum_{n\ge1}e^{-sn}n^{-\alpha}\sim \Gamma(1-\alpha)s^{\alpha-1}
\qquad (s\downarrow0),
\]
we obtain
\[
(1-e^{-s})\sum_{n\ge0}e^{-sn}\overline\varphi(n)\sim s^\alpha,
\qquad s\downarrow0.
\]
Hence,
\[
\hat\varphi(s)=1-s^\alpha+o(s^\alpha),
\qquad s\downarrow0.
\]

Substituting $s=z/\tau$, we get
\[
1-a_\tau\hat\varphi(z/\tau)
=
(1-a_\tau)+a_\tau(z/\tau)^\alpha+o(\tau^{-\alpha}).
\]
Since $(1-a_\tau)\sim \gamma\tau^{-\alpha}$ and $\hat\varphi(z/\tau)\to1$, \eqref{eq:laplace-Stau-proof} yields
\begin{align*}
\widehat{\mathsf S^\tau}(z)
&\to
\lim_{\tau\to\infty}
\frac{(1-a_\tau)\hat\varphi(z/\tau)}
{1-a_\tau\hat\varphi(z/\tau)}
\\
&=
\lim_{\tau\to\infty}
\frac{(1-a_\tau)\left(1-(z/\tau)^\alpha+o(\tau^{-\alpha})\right)}
{(1-a_\tau)+a_\tau(z/\tau)^\alpha+o(\tau^{-\alpha})}
=
\frac{\gamma}{\gamma+z^\alpha}.
\end{align*}

Now recall the classical Laplace transform identity
\[
\int_0^\infty e^{-zt}\,\gamma t^{\alpha-1}E_{\alpha,\alpha}(-\gamma t^\alpha)\,dt
=
\frac{\gamma}{\gamma+z^\alpha},
\qquad z>0.
\]
Thus the pointwise Laplace-transform limit is exactly the Laplace transform of the probability density
\[
f_{\alpha,\gamma}(t):=\gamma t^{\alpha-1}E_{\alpha,\alpha}(-\gamma t^\alpha),
\qquad t>0.
\]
Since each $\mathsf S^\tau(t)\,dt$ is a probability measure on $[0,\infty)$, the continuity theorem for Laplace transforms implies
\[
\mathsf S^\tau(t)\,dt
\Rightarrow
f_{\alpha,\gamma}(t)\,dt
\qquad (\tau\to\infty).
\]
This completes the proof.
\end{proof}

%%%%%%%%%%%%%%%%%%%%%%%%%%%%%%%%%%%%%%%%%%%%%%%%
\subsection{Lemma~\ref{lem:psi-pointwise}}\label{app:proof-psi-pointwise}

For later use, define the discrete cumulative resolvent by
\[
  \mathcal G_n^\tau
  :=
  \frac{1-a_\tau}{\gamma a_\tau}
  \sum_{k=1}^n\psi_k^\tau,
  \qquad n\ge0,
\]
with $\mathcal G_0^\tau=0$.

\begin{lemma}[Doney-type identity and pointwise resolvent bound]
\label{lem:psi-pointwise}
Define
\[
  R_n^\tau:=1_{\{n=0\}}+\psi_n^\tau,
  \qquad n\ge0.
\]
Then, for every $n\ge1$,
\begin{equation}\label{eq:doney-identity}
  nR_n^\tau
  =
  a_\tau\sum_{m=1}^n m\varphi_m(R^\tau*R^\tau)_{n-m}.
\end{equation}
Moreover, for every $T>0$, there exists $C_T>0$, independent of $\tau$, such
that
\begin{equation}\label{eq:psi-pointwise}
  \psi_n^\tau\le C_T n^{\alpha-1},
  \qquad 1\le n\le N_\tau.
\end{equation}
Equivalently,
\begin{equation}\label{eq:G-increment-pointwise}
  \Delta\mathcal G_n^\tau
  :=
  \mathcal G_n^\tau-\mathcal G_{n-1}^\tau
  =
  \frac{1-a_\tau}{\gamma a_\tau}\psi_n^\tau
  \le C_T\tau^{-\alpha}n^{\alpha-1},
  \qquad 1\le n\le N_\tau.
\end{equation}
\end{lemma}

\begin{proof}
Let
\[
  \mathcal R_\tau(q):=\sum_{n\ge0}R_n^\tau q^n
  =1+\sum_{n\ge1}\psi_n^\tau q^n
  =\frac{1}{1-a_\tau\Phi(q)}.
\]
Differentiating gives
\[
  q\mathcal R_\tau'(q)
  =
  a_\tau\,q\Phi'(q)\,\left(\mathcal R_\tau(q)\right)^2.
\]
Since
\[
  q\Phi'(q)=\sum_{m\ge1}m\varphi_m q^m,
\]
comparing coefficients of $q^n$ yields \eqref{eq:doney-identity}.

Now write $R^\tau=\delta_0+\psi^\tau$, where $\delta_0$ denotes the unit mass at $0$, i.e. $\delta_0(n)=1_{\{n=0\}}$. Expanding \eqref{eq:doney-identity} gives
\begin{equation}\label{eq:doney-expanded}
  n\psi_n^\tau
  =
  a_\tau n\varphi_n
  +2a_\tau\sum_{m=1}^{n-1}m\varphi_m\,\psi_{n-m}^\tau
  +a_\tau\sum_{m=1}^{n-1}m\varphi_m\,(\psi^\tau*\psi^\tau)_{n-m}.
\end{equation}
Because
\[
  \varphi_n\asymp n^{-1-\alpha},
  \qquad
  m\varphi_m\asymp m^{-\alpha},
\]
we prove by induction that $\psi_j^\tau\le C j^{\alpha-1}$ for $1\le j\le n-1$, with $C$ large enough.

{Under this induction hypothesis, the convolution $(\psi^\tau*\psi^\tau)_\ell\le C^2\sum_{r=1}^{\ell-1}r^{\alpha-1}(\ell-r)^{\alpha-1}$. This is a standard beta-integral convolution: splitting at $\lfloor\ell/2\rfloor$ shows both halves contribute $\mathcal O(\ell^{2\alpha-1})$, giving
\[
  (\psi^\tau*\psi^\tau)_\ell \le C_1 \ell^{2\alpha-1}.
\]
The linear and quadratic sums in \eqref{eq:doney-expanded} are then controlled by the same midpoint-split argument. For the linear term $\sum_{m=1}^{n-1}m^{-\alpha}(n-m)^{\alpha-1}$, the two halves give $n^{\alpha-1}\cdot\mathcal O(n^{1-\alpha})$ and $n^{-\alpha}\cdot\mathcal O(n^\alpha)$, both $\mathcal O(1)$. For the quadratic term $\sum_{m=1}^{n-1}m^{-\alpha}(n-m)^{2\alpha-1}$, the two halves give $\mathcal O(n^\alpha)$. Inserting into \eqref{eq:doney-expanded} yields $n\psi_n^\tau\le C_5 n^\alpha$, and hence $\psi_n^\tau\le C_5 n^{\alpha-1}$. Choosing $C\ge C_5$ closes the induction.}
Finally, since
\[
  \Delta\mathcal G_n^\tau
  =
  \frac{1-a_\tau}{\gamma a_\tau}\psi_n^\tau,
\]
\eqref{eq:G-increment-pointwise} is immediate from \eqref{eq:psi-pointwise}.
\end{proof}

%%%%%%%%%%%%%%%%%%%%%%%%%%%%%%%%%%%%%%%%%%%%%%%%%%%%%%%
\subsection{Proof of Proposition~\ref{proposition1}}\label{app:proposition1}
Fix $T>0$.

\medskip
\noindent
\textit{Step 1: uniform moment bounds for the intensities.}
We first establish the moment estimates for the intensities that will be used throughout the proof.

For $1\le s\le \lfloor \tau T\rfloor$, using Definition~\ref{def:params} and
\[
  1-\sum_{u=1}^{s-1}\varphi_u^\tau
  =(1-a_\tau)+a_\tau\sum_{u=s}^{\infty}\varphi_u,
\]
we obtain
\begin{align*}
  \hat\mu_\tau(s)
  &=\mu_\tau+\xi_0\mu_\tau\left(
      \frac{1}{1-a_\tau}\left((1-a_\tau)+a_\tau\sum_{u=s}^{\infty}\varphi_u\right)
      -a_\tau\sum_{u=1}^{s-1}\varphi_u
    \right) \\
  &\le C\mu_\tau\left(
      1+\frac{1}{1-a_\tau}\sum_{u=s}^{\infty}\varphi_u
    \right).
\end{align*}
For the kernel from Definition~\ref{def:params},
\[
  \sum_{u=s}^{\infty}\varphi_u
  =\frac{1}{\Gamma(1-\alpha)}(s-1)^{-\alpha},
  \qquad s\ge2,
\]
while $\sum_{u=1}^{\infty}\varphi_u=1$. Since
$\mu_\tau=\mu\tau^{\alpha-1}$ and $1-a_\tau=\gamma\tau^{-\alpha}$, it follows that
\begin{equation}\label{eq:muhat-uniform-bound-app-prop1}
  \sup_{1\le s\le \lfloor \tau T\rfloor}\hat\mu_\tau(s)
  \lesssim \tau^{2\alpha-1}.
\end{equation}

Next fix $n\le \lfloor \tau T\rfloor$. By \eqref{eq:lambda-decomp},
\[
  \lambda_n^{\tau,+}
  =
  \hat\mu_\tau(n)
  +
  \sum_{s=1}^{n-1}\psi_{n-s}^\tau \hat\mu_\tau(s)
  +
  \kappa_n^{\tau,+},
\]
where
\[
  \kappa_n^{\tau,+}
  :=
  \frac1{1+\beta}\sum_{s=1}^{n-1}\psi_{n-s}^{\tau}
  \left(
    M_s^{\tau,+}-M_{s-1}^{\tau,+}
    +\beta\left(M_s^{\tau,-}-M_{s-1}^{\tau,-}\right)
  \right).
\]

We first bound the deterministic convolution term $\sum_{s=1}^{n-1}\psi_{n-s}^\tau \hat\mu_\tau(s)$. From the explicit estimate obtained above,
\[
  \hat\mu_\tau(s)\lesssim \tau^{\alpha-1}+\tau^{2\alpha-1}s^{-\alpha},
\]
uniformly for $1\le s\le \lfloor \tau T\rfloor$. Writing $j=n-s$, we obtain
\begin{align}\label{for:two:terms}
  \sum_{s=1}^{n-1}\psi_{n-s}^\tau \hat\mu_\tau(s)
  \lesssim
  \tau^{\alpha-1}\sum_{j=1}^{n-1}\psi_j^\tau
  +\tau^{2\alpha-1}\sum_{j=1}^{n-1}\psi_j^\tau (n-j)^{-\alpha}.
\end{align}

For the first term on the right side of \eqref{for:two:terms},
\[
  \tau^{\alpha-1}\sum_{j=1}^{n-1}\psi_j^\tau
  \le
  \tau^{\alpha-1}\sum_{j\ge1}\psi_j^\tau
  =
  \tau^{\alpha-1}\frac{a_\tau}{1-a_\tau}
  \lesssim \tau^{2\alpha-1}.
\]

For the second term on the right side of \eqref{for:two:terms}, Lemma~\ref{lem:psi-pointwise} gives
\[
  \psi_j^\tau\le C_T j^{\alpha-1},
  \qquad 1\le j\le \lfloor \tau T\rfloor.
\]
Hence,
\[
  \sum_{j=1}^{n-1}\psi_j^\tau (n-j)^{-\alpha}
  \lesssim
  \sum_{j=1}^{n-1} j^{\alpha-1}(n-j)^{-\alpha}.
\]
We split this sum at $j=\lfloor n/2\rfloor$. For $1\le j\le \lfloor n/2\rfloor$,
we have $n-j\asymp n$, and thus
\[
  \sum_{j=1}^{\lfloor n/2\rfloor} j^{\alpha-1}(n-j)^{-\alpha}
  \lesssim
  n^{-\alpha}\sum_{j=1}^{\lfloor n/2\rfloor}j^{\alpha-1}
  \lesssim
  n^{-\alpha}n^\alpha
  \lesssim 1.
\]
For $\lfloor n/2\rfloor+1\le j\le n-1$, write $r=n-j$. Then $j\asymp n$, so that
\[
  \sum_{j=\lfloor n/2\rfloor+1}^{n-1} j^{\alpha-1}(n-j)^{-\alpha}
  \lesssim
  n^{\alpha-1}\sum_{r=1}^{\lceil n/2\rceil} r^{-\alpha}
  \lesssim
  n^{\alpha-1}n^{1-\alpha}
  \lesssim 1.
\]
Therefore,
\[
  \sum_{j=1}^{n-1}\psi_j^\tau (n-j)^{-\alpha}\lesssim 1,
\]
uniformly in $n\le \lfloor \tau T\rfloor$. Consequently,
\begin{equation}\label{eq:conv-muhat-bound-app-prop1}
  \sup_{1\le n\le \lfloor \tau T\rfloor}
  \sum_{s=1}^{n-1}\psi_{n-s}^\tau \hat\mu_\tau(s)
  \lesssim \tau^{2\alpha-1}.
\end{equation}

Taking expectations in \eqref{eq:lambda-decomp} and using that $\kappa_n^{\tau,+}$ has mean zero,
we deduce from \eqref{eq:muhat-uniform-bound-app-prop1} and
\eqref{eq:conv-muhat-bound-app-prop1} that
\begin{equation}\label{eq:Elambda-prop1}
  \sup_{1\le n\le \lfloor \tau T\rfloor}\E[\lambda_n^{\tau,+}]
  \lesssim \tau^{2\alpha-1}.
\end{equation}

We next bound the second moment. Since martingale increments at different times are orthogonal,
\begin{align*}
\E\!\left[\left(\kappa_n^{\tau,+}\right)^2\right]
  =
  \frac1{(1+\beta)^2}
  \sum_{s=1}^{n-1}(\psi_{n-s}^\tau)^2
  \E\!\left[
    \lambda_s^{\tau,+}+\beta^2\lambda_s^{\tau,-}
  \right]=
  \frac{1+\beta^2}{(1+\beta)^2}
  \sum_{s=1}^{n-1}(\psi_{n-s}^\tau)^2\E[\lambda_s^{\tau,+}],
\end{align*}
where we used $\lambda_s^{\tau,+}=\lambda_s^{\tau,-}$. Using \eqref{eq:Elambda-prop1} and
Lemma~\ref{lem:psi-pointwise},
\[
  \E\!\left[(\kappa_n^{\tau,+})^2\right]
  \lesssim
  \tau^{2\alpha-1}\sum_{j=1}^{n-1}(\psi_j^\tau)^2
  \lesssim
  \tau^{2\alpha-1}\sum_{j=1}^{\lfloor \tau T\rfloor}j^{2\alpha-2}
  \lesssim \tau^{4\alpha-2},
\]
since $2\alpha-2>-1$.

Returning to \eqref{eq:lambda-decomp}, the deterministic part
\[
  \hat\mu_\tau(n)+\sum_{s=1}^{n-1}\psi_{n-s}^\tau\hat\mu_\tau(s)
\]
is $\mathcal O(\tau^{2\alpha-1})$ uniformly in $n\le \lfloor \tau T\rfloor$, and hence
\begin{equation}\label{eq:Elambda2-prop1}
  \sup_{1\le n\le \lfloor \tau T\rfloor}\E[(\lambda_n^{\tau,+})^2]
  \lesssim \tau^{4\alpha-2}.
\end{equation}
The same bounds hold for the ``$-$'' component.

\medskip
\noindent
\textit{Step 2: $C$-tightness of $\mathbf Y^\tau$.}
We apply Lemma~\ref{xu2023} componentwise to $(Y^{\tau,+})_{\tau\ge1}$ and
$(Y^{\tau,-})_{\tau\ge1}$. Since $Y_0^{\tau,\pm}=0$, the initial-moment condition is trivial.

We check the lemma for $Y^{\tau,+}$; the proof for $Y^{\tau,-}$ is identical.

\smallskip
\noindent
\underline{Condition (i).}
Fix $\vartheta>2$. Since $Y^{\tau,+}$ is piecewise constant on the grid
$\{n/\tau:n\ge0\}$,
\[
  \sup_{k=0,\dots,\lfloor T\tau^\vartheta\rfloor-1}
  \sup_{h\in[0,\tau^{-\vartheta}]}
  \left|\Delta_hY_{k/\tau^\vartheta}^{\tau,+}\right|
  \le
  \frac{1-a_\tau}{\tau^\alpha\mu}
  \sup_{1\le s\le \lfloor \tau T\rfloor}X_s^{\tau,+}.
\]
Hence it suffices to show
\[
  \frac{1-a_\tau}{\tau^\alpha\mu}
  \sup_{1\le s\le \lfloor \tau T\rfloor}X_s^{\tau,+}\to0
  \qquad\text{in probability.}
\]
By the union bound and Markov's inequality,
\begin{align*}
  \mathbb P\!\left(
    \frac{1-a_\tau}{\tau^\alpha\mu}
    \sup_{1\le s\le \lfloor \tau T\rfloor}X_s^{\tau,+}>\varepsilon
  \right)
  &\le
  \sum_{s=1}^{\lfloor \tau T\rfloor}
  \frac{\E[(X_s^{\tau,+})^2]}
       {\left(\frac{\tau^\alpha\mu}{1-a_\tau}\varepsilon\right)^2}.
\end{align*}
Since $X_s^{\tau,+}\mid \mathcal F_{s-1}\sim \mathrm{Poisson}(\lambda_s^{\tau,+})$,
\[
  \E[(X_s^{\tau,+})^2]
  =
  \E[\lambda_s^{\tau,+}]
  +\E[(\lambda_s^{\tau,+})^2].
\]
Using \eqref{eq:Elambda-prop1}, \eqref{eq:Elambda2-prop1}, and
\[
  \frac{\tau^\alpha\mu}{1-a_\tau}\asymp \tau^{2\alpha},
\]
we obtain
\[
  \sum_{s=1}^{\lfloor \tau T\rfloor}
  \frac{\E[(X_s^{\tau,+})^2]}
       {\left(\frac{\tau^\alpha\mu}{1-a_\tau}\varepsilon\right)^2}
  \lesssim \tau^{-1}\to0.
\]
Thus condition (i) holds.

\smallskip
\noindent
\underline{Condition (ii).}
Take $p=2$ and use the pairs
$(a_1,b_1)=(0,2)$ and $(a_2,b_2)=(1,1)$.
Then
\[
  \varrho_*=\min\{2,\,1+1/\vartheta\}>1.
\]
For $t\in[0,T]$ and $h\in(0,1)$,
\begin{align*}
  \E\!\left[\left|\Delta_hY_t^{\tau,+}\right|^2\right]
  &=
  \E\!\left[
    \left(
      \frac{1-a_\tau}{\tau^\alpha\mu}
      \sum_{i=\lfloor \tau t\rfloor+1}^{\lfloor \tau(t+h)\rfloor}X_i^{\tau,+}
    \right)^2
  \right] \\
  &\le
  \left(\frac{1-a_\tau}{\tau^\alpha\mu}\right)^2
  \left(\lfloor \tau(t+h)\rfloor-\lfloor \tau t\rfloor\right)
  \sum_{i=\lfloor \tau t\rfloor+1}^{\lfloor \tau(t+h)\rfloor}
  \E\left[\left(X_i^{\tau,+}\right)^2\right].
\end{align*}
Using \eqref{eq:Elambda-prop1}--\eqref{eq:Elambda2-prop1}, we have
\[
  \sup_{1\le i\le \lfloor \tau T\rfloor}\E\left[\left(X_i^{\tau,+}\right)^2\right]
  \lesssim \tau^{4\alpha-2},
\]
and therefore
\[
  \E\!\left[\left|\Delta_hY_t^{\tau,+}\right|^2\right]
  \lesssim
  \left(\frac{1-a_\tau}{\tau^\alpha}\right)^2
  (\tau h+1)^2\tau^{4\alpha-2}
  \lesssim h^2+\frac{h}{\tau}.
\]
Hence condition (ii) of Lemma~\ref{xu2023} holds. Thus
$(Y^{\tau,+})_{\tau\ge1}$ is $C$-tight. The same proof applies to
$(Y^{\tau,-})_{\tau\ge1}$, so $\mathbf Y^\tau$ is $C$-tight.

\medskip
\noindent
\textit{Step 3: $C$-tightness of $\boldsymbol\Lambda^\tau$.}
By Lemma~\ref{lem:tightness},
\[
  \sup_{t\in[0,T]}\|\boldsymbol\Lambda_t^\tau-\mathbf Y_t^\tau\|_1\to0
  \qquad\text{in probability.}
\]
Since $\mathbf Y^\tau$ is $C$-tight, it follows that $\boldsymbol\Lambda^\tau$ is also
$C$-tight.

\medskip
\noindent
\textit{Step 4: $C$-tightness of $\mathbf Z^\tau$.}
We verify the $C$-tightness of $Z^{\tau,+}$ and $Z^{\tau,-}$ via a
martingale tightness criterion (e.g.\ \cite[Proposition~12]{rebolledo1980central}):
it suffices to show that their predictable quadratic variation processes are
$C$-tight and that their jumps vanish in probability.

By definition,
\[
  Z_t^{\tau,+}
  =
  \sqrt{\frac{\tau^\alpha\mu}{1-a_\tau}}
  \left(Y_t^{\tau,+}-\Lambda_t^{\tau,+}\right)
  =
  \sqrt{\frac{1-a_\tau}{\tau^\alpha\mu}}\,M_{\lfloor \tau t\rfloor}^{\tau,+},
\]
and similarly
\[
  Z_t^{\tau,-}
  =
  \sqrt{\frac{1-a_\tau}{\tau^\alpha\mu}}\,M_{\lfloor \tau t\rfloor}^{\tau,-}.
\]
Hence $Z^{\tau,+}$ and $Z^{\tau,-}$ are square-integrable c\`adl\`ag martingales. Moreover,
\[
  \langle Z^{\tau,+},Z^{\tau,+}\rangle_t
  =
  \frac{1-a_\tau}{\tau^\alpha\mu}
  \sum_{s=1}^{\lfloor \tau t\rfloor}\lambda_s^{\tau,+}
  =
  \Lambda_t^{\tau,+},
\]
and similarly
\[
  \langle Z^{\tau,-},Z^{\tau,-}\rangle_t
  =
  \Lambda_t^{\tau,-}.
\]
Since $(\boldsymbol\Lambda^\tau)_{\tau\ge1}$ is $C$-tight, both sequences
$(\langle Z^{\tau,+},Z^{\tau,+}\rangle)_{\tau\ge1}$ and
$(\langle Z^{\tau,-},Z^{\tau,-}\rangle)_{\tau\ge1}$ are $C$-tight.
It remains to show that
\[
  \sup_{t\in[0,T]}\left|\Delta Z_t^{\tau,+}\right|\xrightarrow{\mathbb P}0,
  \qquad
  \sup_{t\in[0,T]}\left|\Delta Z_t^{\tau,-}\right|\xrightarrow{\mathbb P}0.
\]
Since
\[
  \Delta Z_t^{\tau,+}
  =
  \sqrt{\frac{1-a_\tau}{\tau^\alpha\mu}}
  \left(M_{\lfloor \tau t\rfloor}^{\tau,+}-M_{\lfloor \tau t\rfloor-1}^{\tau,+}\right),
\]
we have
\[
  \sup_{t\in[0,T]}\left|\Delta Z_t^{\tau,+}\right|
  \le
  \sqrt{\frac{1-a_\tau}{\tau^\alpha\mu}}
  \sup_{1\le s\le \lfloor \tau T\rfloor}
  \left|X_s^{\tau,+}-\lambda_s^{\tau,+}\right|.
\]
We use the fourth-moment Markov inequality. For every $\varepsilon>0$,
\begin{align*}
  \mathbb P\!\left(
    \sup_{t\in[0,T]}\left|\Delta Z_t^{\tau,+}\right|>\varepsilon
  \right)
  &\le
  \sum_{s=1}^{\lfloor \tau T\rfloor}
  \mathbb P\!\left(
    \left|X_s^{\tau,+}-\lambda_s^{\tau,+}\right|
    >
    \varepsilon\sqrt{\frac{\tau^\alpha\mu}{1-a_\tau}}
  \right) \\
  &\le
  \sum_{s=1}^{\lfloor \tau T\rfloor}
  \frac{
    \E\!\left[\left|X_s^{\tau,+}-\lambda_s^{\tau,+}\right|^4\right]
  }{
    \varepsilon^4\left(\frac{\tau^\alpha\mu}{1-a_\tau}\right)^2
  }.
\end{align*}

Now conditionally on $\mathcal F_{s-1}$,
\[
  X_s^{\tau,+}\mid \mathcal F_{s-1}\sim \mathrm{Poisson}(\lambda_s^{\tau,+}),
\]
and the centered fourth moment of a Poisson random variable is
\[
  \E\!\left[
    \left(X_s^{\tau,+}-\lambda_s^{\tau,+}\right)^4 \mid \mathcal F_{s-1}
  \right]
  =
  \lambda_s^{\tau,+}+3(\lambda_s^{\tau,+})^2.
\]
Taking expectations and using \eqref{eq:Elambda-prop1} and \eqref{eq:Elambda2-prop1},
we obtain
\[
  \sup_{1\le s\le \lfloor \tau T\rfloor}
  \E\!\left[\left|X_s^{\tau,+}-\lambda_s^{\tau,+}\right|^4\right]
  \lesssim
  \tau^{2\alpha-1}+\tau^{4\alpha-2}
  \lesssim \tau^{4\alpha-2}.
\]
Since
\[
  \left(\frac{\tau^\alpha\mu}{1-a_\tau}\right)^2\asymp \tau^{4\alpha},
\]
it follows that
\begin{align*}
  \mathbb P\!\left(
    \sup_{t\in[0,T]}\left|\Delta Z_t^{\tau,+}\right|>\varepsilon
  \right)
  &\lesssim
  \lfloor \tau T\rfloor \cdot \tau^{4\alpha-2}\cdot \tau^{-4\alpha}
  \lesssim \tau^{-1}\to0.
\end{align*}
Thus,
\[
\sup_{t\in[0,T]}\left|\Delta Z_t^{\tau,+}\right|\to0
  \qquad\text{in probability.}
\]
The same proof applies to $Z^{\tau,-}$.

Therefore, by \cite[Proposition~12]{rebolledo1980central},
$(Z^{\tau,+})_{\tau\ge1}$ and $(Z^{\tau,-})_{\tau\ge1}$ are $C$-tight.
Hence $(\mathbf Z^\tau)_{\tau\ge1}$ is $C$-tight on $[0,T]$. Moreover,
\[
  \sup_{t\in[0,T]}\left\|\Delta \mathbf Z_t^\tau\right\|_1
  \le
  \sup_{t\in[0,T]}\left|\Delta Z_t^{\tau,+}\right|
  +
  \sup_{t\in[0,T]}\left|\Delta Z_t^{\tau,-}\right|
  \to 0
\]
in probability.

\medskip
\noindent
\textit{Step 5: identification of the quadratic variation of any limit point.}
Let $(\mathbf Y,\mathbf Z)$ be a limit point of $(\mathbf Y^\tau,\mathbf Z^\tau)$.

For the ``$+$'' component,
\[
  [Z^{\tau,+},Z^{\tau,+}]_{n/\tau}
  =
  \frac{1-a_\tau}{\tau^\alpha\mu}
  \sum_{s=1}^{n}(X_s^{\tau,+}-\lambda_s^{\tau,+})^2,
\]
so that
\[
  U_n^{\tau,+}
  :=
  [Z^{\tau,+},Z^{\tau,+}]_{n/\tau}
  -
  \langle Z^{\tau,+},Z^{\tau,+}\rangle_{n/\tau}
  =
  \frac{1-a_\tau}{\tau^\alpha\mu}
  \sum_{s=1}^{n}
  \left((X_s^{\tau,+}-\lambda_s^{\tau,+})^2-\lambda_s^{\tau,+}\right).
\]
Then $(U_n^{\tau,+})_{n\ge0}$ is a martingale, and by orthogonality of martingale
increments,
\begin{align*}
  \E\!\left[\left(U_{\lfloor t\tau\rfloor}^{\tau,+}\right)^2\right]
  &=
  \left(\frac{1-a_\tau}{\tau^\alpha\mu}\right)^2
  \sum_{s=1}^{\lfloor t\tau\rfloor}
  \E\!\left[
    \Var\!\left(
      (X_s^{\tau,+}-\lambda_s^{\tau,+})^2\mid\mathcal F_{s-1}
    \right)
  \right].
\end{align*}
Since $X_s^{\tau,+}\mid\mathcal F_{s-1}\sim\mathrm{Poisson}(\lambda_s^{\tau,+})$,
\[
  \Var\!\left((X_s^{\tau,+}-\lambda_s^{\tau,+})^2\mid\mathcal F_{s-1}\right)
  =
  2(\lambda_s^{\tau,+})^2+\lambda_s^{\tau,+}.
\]
Therefore,
\[
  \E\!\left[\left(U_{\lfloor t\tau\rfloor}^{\tau,+}\right)^2\right]
  \lesssim
  \left(\frac{1-a_\tau}{\tau^\alpha\mu}\right)^2
  \sum_{s=1}^{\lfloor t\tau\rfloor}
  \left(\E[(\lambda_s^{\tau,+})^2]+\E[\lambda_s^{\tau,+}]\right).
\]
Using \eqref{eq:Elambda-prop1} and \eqref{eq:Elambda2-prop1}, we get
\[
  \E\!\left[\left(U_{\lfloor t\tau\rfloor}^{\tau,+}\right)^2\right]
  \lesssim \tau^{-1},
\]
uniformly for $t\in[0,T]$. Hence, by Doob's inequality,
\[
  \sup_{t\in[0,T]}
  \left|
    [Z^{\tau,+},Z^{\tau,+}]_t-\Lambda_t^{\tau,+}
  \right|
  \to0
  \qquad\text{in probability.}
\]
Combining this with Lemma~\ref{lem:tightness},
\[
  \sup_{t\in[0,T]}
  \left|
    [Z^{\tau,+},Z^{\tau,+}]_t-Y_t^{\tau,+}
  \right|
  \to0
  \qquad\text{in probability.}
\]
Passing to the limit yields
\[
  [Z^+,Z^+]_t=Y_t^+.
\]
The same argument gives
\[
  [Z^-,Z^-]_t=Y_t^-.
\]

For the cross-variation,
\[
\left[Z^{\tau,+},Z^{\tau,-}\right]_{n/\tau}
  =
  \frac{1-a_\tau}{\tau^\alpha\mu}
  \sum_{s=1}^{n}
  \left(X_s^{\tau,+}-\lambda_s^{\tau,+}\right)\left(X_s^{\tau,-}-\lambda_s^{\tau,-}\right).
\]
Because $X_s^{\tau,+}$ and $X_s^{\tau,-}$ are conditionally independent given
$\mathcal F_{s-1}$, the process
\[
  V_n^\tau:=[Z^{\tau,+},Z^{\tau,-}]_{n/\tau}
\]
is a martingale. Moreover,
\begin{align*}
  \E\!\left[\left(V_{\lfloor t\tau\rfloor}^\tau\right)^2\right]
  &=
  \left(\frac{1-a_\tau}{\tau^\alpha\mu}\right)^2
  \sum_{s=1}^{\lfloor t\tau\rfloor}
  \E\!\left[
    \lambda_s^{\tau,+}\lambda_s^{\tau,-}
  \right].
\end{align*}
Since $\lambda_s^{\tau,+}=\lambda_s^{\tau,-}$ and \eqref{eq:Elambda2-prop1} holds,
\[
  \E\!\left[\left(V_{\lfloor t\tau\rfloor}^\tau\right)^2\right]
  \lesssim \tau^{-1}.
\]
Hence, by Doob's inequality,
\[
  \sup_{t\in[0,T]}|[Z^{\tau,+},Z^{\tau,-}]_t|
  \to0
  \qquad\text{in probability,}
\]
and therefore
\[
  [Z^+,Z^-]_t=0.
\]

Since $(\mathbf Z^\tau)_{\tau\ge1}$ is $C$-tight, every limit point $\mathbf Z$ is continuous.
Therefore $\mathbf Z$ is a continuous martingale with
\[
  [\mathbf Z,\mathbf Z]=\mathrm{diag}(\mathbf Y).
\]
This completes the proof.
%%%%%%%%%%%%%%%%%%%%%%%%%%%%%%%%%

\subsection{Proof of Lemma~\ref{lem:tightness}}\label{app:tightness}
Fix $T>0$. It suffices to prove
\[
  \sup_{t \in [0,T]} \left|\Lambda_t^{\tau,+} - Y_t^{\tau,+}\right| \to 0
  \qquad\text{in probability as }\tau\to\infty,
\]
since the proof for the ``$-$'' component is identical.

Since
\[
  Y_t^{\tau,+}-\Lambda_t^{\tau,+}
  =\frac{1-a_\tau}{\tau^\alpha\mu}\,M_{\lfloor t\tau\rfloor}^{\tau,+},
\]
Doob's $L^2$ inequality yields
\begin{align*}
  \E\left[\sup_{t\in[0,T]}|Y_t^{\tau,+}-\Lambda_t^{\tau,+}|^2\right]
  &\le \left(\frac{1-a_\tau}{\tau^\alpha\mu}\right)^2
  \E\left[\sup_{0\le n\le \lfloor\tau T\rfloor}|M_n^{\tau,+}|^2\right]\\
  &\le 4\left(\frac{1-a_\tau}{\tau^\alpha\mu}\right)^2
  \E\left[|M_{\lfloor\tau T\rfloor}^{\tau,+}|^2\right].
\end{align*}
Because $M^{\tau,+}$ is a square-integrable martingale,
\[
  \E\left[\left|M_{\lfloor\tau T\rfloor}^{\tau,+}\right|^2\right]
  =\E\left[\left\langle M^{\tau,+}\right\rangle_{\lfloor\tau T\rfloor}\right]
  =\sum_{s=1}^{\lfloor\tau T\rfloor}\E[\lambda_s^{\tau,+}].
\]
Using the uniform estimate \eqref{eq:Elambda-prop1}, we have
\[
  \sup_{1\le s\le \lfloor\tau T\rfloor}\E[\lambda_s^{\tau,+}]
  \lesssim \tau^{2\alpha-1},
\]
and therefore
\[
  \sum_{s=1}^{\lfloor\tau T\rfloor}\E[\lambda_s^{\tau,+}]
  \lesssim \tau^{2\alpha}.
\]
Hence,
\[
  \E\left[\sup_{t\in[0,T]}\left|Y_t^{\tau,+}-\Lambda_t^{\tau,+}\right|^2\right]
  \lesssim
  \left(\frac{1-a_\tau}{\tau^\alpha\mu}\right)^2\tau^{2\alpha}
  \lesssim \tau^{-2\alpha}\to0.
\]
Thus,
\[
  \sup_{t\in[0,T]}\left|Y_t^{\tau,+}-\Lambda_t^{\tau,+}\right|
  \to 0
  \qquad\text{in }L^2,
\]
and hence also in probability. This completes the proof.
%%%%%%%%%%%%%%%%%%%%%%%%%%%%%%%%%

\subsection{Proof of Proposition~\ref{prop2}}\label{app:prop2}

Fix $T>0$. Let
\[
\left(Y^{\tau_k,+},Y^{\tau_k,-},Z^{\tau_k,+},Z^{\tau_k,-}\right)
\]
be a subsequence converging in law on $D([0,T],\mathbb R^4)$ to a limit point
\[
(Y^+,Y^-,Z^+,Z^-).
\]
By Proposition~\ref{proposition1}, every limit point is continuous. Hence, by the
Skorokhod representation theorem, after possibly enlarging the probability space
and still denoting the subsequence by $\tau_k$, we may assume that
\begin{equation}\label{eq:prop2-skorokhod}
\sup_{t\in[0,T]}
\left(
\left|Y_t^{\tau_k,+}-Y_t^+\right|
+
\left|Y_t^{\tau_k,-}-Y_t^-\right|
+
\left|Z_t^{\tau_k,+}-Z_t^+\right|
+
\left|Z_t^{\tau_k,-}-Z_t^-\right|
\right)\to0
\qquad\text{a.s.}
\end{equation}
as $k\to\infty$.

\medskip
\noindent
\textit{Step 1: the two limit components of $\mathbf Y^\tau$ coincide.}

Recall that
\[
\Lambda_t^{\tau,+}=\Lambda_t^{\tau,-},
\qquad t\in[0,T].
\]
Hence, by Lemma~\ref{lem:tightness},
\begin{align*}
\sup_{t\in[0,T]}\left|Y_t^{\tau,+}-Y_t^{\tau,-}\right|
  &\le
  \sup_{t\in[0,T]}\left|Y_t^{\tau,+}-\Lambda_t^{\tau,+}\right|
  +
  \sup_{t\in[0,T]}\left|Y_t^{\tau,-}-\Lambda_t^{\tau,-}\right|
  \\
  &\to0
  \qquad\text{in probability.}
\end{align*}
Passing to the almost sure Skorokhod realization along the chosen subsequence,
we obtain
\[
  Y_t^+=Y_t^-
  \qquad\text{for all }t\in[0,T],\ \text{a.s.}
\]
We therefore write
\[
  Y_t:=Y_t^+=Y_t^-,
  \qquad t\in[0,T].
\]

\medskip
\noindent
\textit{Step 2: an explicit decomposition for $\Lambda^{\tau,+}$.}

We first show that for every integer $n\ge1$,
\begin{equation}\label{eq:prop2-main-identity}
  \hat\mu_\tau(n)+\sum_{s=1}^{n-1}\psi_{n-s}^\tau\hat\mu_\tau(s)
  =
  \mu_\tau+\xi_0\mu_\tau\frac1{1-a_\tau}
  +\mu_\tau(1-\xi_0)\sum_{s=1}^{n-1}\psi_{n-s}^\tau.
\end{equation}
Indeed, Definition~\ref{def:params} gives
\[
  \hat\mu_\tau(n)
  =
  \mu_\tau
  +
  \xi_0\mu_\tau
  \left(
    \frac{1}{1-a_\tau}\left(1-\sum_{r=1}^{n-1}\varphi_r^\tau\right)
    -
    \sum_{r=1}^{n-1}\varphi_r^\tau
  \right).
\]
Hence,
\begin{align*}
&\hat\mu_\tau(n)+\sum_{s=1}^{n-1}\psi_{n-s}^\tau\hat\mu_\tau(s)
  \\
  &=
  \mu_\tau
  +
  \xi_0\mu_\tau
  \left(
    \frac{1}{1-a_\tau}\left(1-\sum_{r=1}^{n-1}\varphi_r^\tau\right)
    -
    \sum_{r=1}^{n-1}\varphi_r^\tau
  \right)
  \\
  &\quad
  +
  \mu_\tau\sum_{s=1}^{n-1}\psi_{n-s}^\tau
  +
  \xi_0\mu_\tau
  \sum_{s=1}^{n-1}\psi_{n-s}^\tau
  \left(
    \frac{1}{1-a_\tau}\left(1-\sum_{r=1}^{s-1}\varphi_r^\tau\right)
    -
    \sum_{r=1}^{s-1}\varphi_r^\tau
  \right).
\end{align*}
Since
\[
  1-\sum_{r=1}^{s-1}\varphi_r^\tau
  =1-\sum_{r\ge1}\varphi_r^\tau+\sum_{r=s}^{\infty}\varphi_r^\tau
  =(1-a_\tau)+\sum_{r=s}^{\infty}\varphi_r^\tau,
\]
we obtain
\[
  \frac{1}{1-a_\tau}\left(1-\sum_{r=1}^{s-1}\varphi_r^\tau\right)
  -\sum_{r=1}^{s-1}\varphi_r^\tau
  =
  1+\left(1+\frac{1}{1-a_\tau}\right)\sum_{r=s}^{\infty}\varphi_r^\tau
  -\sum_{r\ge1}\varphi_r^\tau.
\]
Since $\sum_{r\ge1}\varphi_r^\tau=a_\tau$, this becomes
\[
  \frac{1}{1-a_\tau}\left(1-\sum_{r=1}^{s-1}\varphi_r^\tau\right)
  -\sum_{r=1}^{s-1}\varphi_r^\tau
  =
  (1-a_\tau)+\frac{1}{1-a_\tau}\sum_{r=s}^{\infty}\varphi_r^\tau+\sum_{r=s}^{\infty}\varphi_r^\tau.
\]
To compute the convolution term explicitly, we rewrite $\hat\mu_\tau(s)$ in terms of the tail
$\sum_{r=s}^\infty \varphi_r^\tau$ and then use the renewal identity
\[
  \psi^\tau=\varphi^\tau+\psi^\tau*\varphi^\tau,
\]
which implies, for every $n\ge1$,
\[
  \sum_{s=1}^{n-1}\psi_{n-s}^\tau\varphi_s^\tau
  =
  \psi_n^\tau-\varphi_n^\tau.
\]
Using this identity inside the convolution term above and rearranging yields
\[
  \sum_{s=1}^{n-1}\psi_{n-s}^\tau\hat\mu_\tau(s)
  =
  \xi_0\mu_\tau\left(\frac1{1-a_\tau}+1\right)\sum_{s=1}^{n-1}\varphi_{n-s}^\tau
  +
  \mu_\tau(1-\xi_0)\sum_{s=1}^{n-1}\psi_{n-s}^\tau.
\]
Substituting this back into the expression for
$\hat\mu_\tau(n)+\sum_{s=1}^{n-1}\psi_{n-s}^\tau\hat\mu_\tau(s)$ gives
\eqref{eq:prop2-main-identity}.

Now recall from \eqref{eq:lambda-decomp} that
\[
  \lambda_n^{\tau,+}
  =
  \hat\mu_\tau(n)
  +
  \sum_{s=1}^{n-1}\psi_{n-s}^\tau\hat\mu_\tau(s)
  +
  \frac1{1+\beta}\sum_{s=1}^{n-1}\psi_{n-s}^\tau
  \left(
    \Delta M_s^{\tau,+}
    +
    \beta \Delta M_s^{\tau,-}
  \right),
\]
where
\[
  \Delta M_s^{\tau,\pm}:=M_s^{\tau,\pm}-M_{s-1}^{\tau,\pm}.
\]
Using \eqref{eq:prop2-main-identity}, we obtain
\begin{align}
  \lambda_n^{\tau,+}
  &=
  \mu_\tau
  +
  \mu_\tau\sum_{s=1}^{n-1}\psi_{n-s}^\tau
  +
  \xi_0\mu_\tau
  \left(
    \frac1{1-a_\tau}
    -
    \sum_{s=1}^{n-1}\psi_{n-s}^\tau
  \right)
  \nonumber\\
  &\quad
  +
  \frac1{1+\beta}\sum_{s=1}^{n-1}\psi_{n-s}^\tau
  \left(
    \Delta M_s^{\tau,+}
    +
    \beta \Delta M_s^{\tau,-}
  \right).
  \label{eq:lambda-explicit-prop2}
\end{align}

We now sum \eqref{eq:lambda-explicit-prop2} over $n=1,\dots,m$.
For the deterministic term, Tonelli's theorem gives
\[
  \sum_{n=1}^{m}\sum_{s=1}^{n-1}\psi_{n-s}^\tau
  =
  \sum_{s=1}^{m-1}\sum_{j=1}^{m-s}\psi_j^\tau
  =
  \sum_{j=1}^{m-1}(m-j)\psi_j^\tau
  =
  \sum_{s=1}^{m-1}s\psi_{m-s}^\tau.
\]
For the martingale term, define
\[
  H_s^\tau:=M_s^{\tau,+}+\beta M_s^{\tau,-},
  \qquad
  \Delta H_s^\tau:=H_s^\tau-H_{s-1}^\tau.
\]
Then
\begin{align*}
  \sum_{n=1}^{m}\sum_{s=1}^{n-1}\psi_{n-s}^\tau \Delta H_s^\tau
  &=
  \sum_{s=1}^{m-1}\left(\sum_{j=1}^{m-s}\psi_j^\tau\right)\Delta H_s^\tau.
\end{align*}
Set
\[
  A_s^\tau:=\sum_{j=1}^{m-s}\psi_j^\tau,
  \qquad 1\le s\le m-1.
\]
By discrete summation by parts,
\[
  \sum_{s=1}^{m-1}A_s^\tau \Delta H_s^\tau
  =
  A_{m-1}^\tau H_{m-1}^\tau
  +
  \sum_{s=1}^{m-2}(A_s^\tau-A_{s+1}^\tau)H_s^\tau.
\]
Since
\[
  A_s^\tau-A_{s+1}^\tau=\psi_{m-s}^\tau,
  \qquad
  A_{m-1}^\tau=\psi_1^\tau,
\]
this yields
\[
  \sum_{n=1}^{m}\sum_{s=1}^{n-1}\psi_{n-s}^\tau \Delta H_s^\tau
  =
  \sum_{s=1}^{m-1}\psi_{m-s}^\tau H_s^\tau.
\]
Consequently,
\begin{align}
  \sum_{n=1}^{m}\lambda_n^{\tau,+}
  &=
  \mu_\tau m
  +
  \mu_\tau\sum_{s=1}^{m-1}s\psi_{m-s}^\tau
  +
  \xi_0\mu_\tau
  \left(
    \frac{m}{1-a_\tau}
    -
    \sum_{s=1}^{m-1}s\psi_{m-s}^\tau
  \right)
  \nonumber\\
  &\quad
  +
  \frac1{1+\beta}
  \sum_{s=1}^{m-1}\psi_{m-s}^\tau
  \left(
    M_s^{\tau,+}+\beta M_s^{\tau,-}
  \right).
  \label{eq:sum-lambda-prop2}
\end{align}

Taking $m=\lfloor \tau t\rfloor$ and multiplying by
$\frac{1-a_\tau}{\tau^\alpha\mu}$, we obtain
\begin{equation}\label{eq:Lambda-decomp-prop2}
  \Lambda_t^{\tau,+}=\Lambda_t^{\tau,-}=T_1^\tau(t)+T_2^\tau(t)+T_3^\tau(t),
\end{equation}
where
\begin{align*}
  T_1^\tau(t)
  &:=
  (1-a_\tau)\frac{\lfloor \tau t\rfloor}{\tau},
  \\
  T_2^\tau(t)
  &:=
  \frac{1-a_\tau}{\tau}
  \sum_{s=1}^{\lfloor \tau t\rfloor-1}s\psi_{\lfloor \tau t\rfloor-s}^\tau
  +
  \xi_0
  \left(
    \frac{\lfloor \tau t\rfloor}{\tau}
    -
    \frac{1-a_\tau}{\tau}
    \sum_{s=1}^{\lfloor \tau t\rfloor-1}s\psi_{\lfloor \tau t\rfloor-s}^\tau
  \right),
  \\
  T_3^\tau(t)
  &:=
  \frac{1-a_\tau}{\tau^\alpha\mu(1+\beta)}
  \sum_{s=1}^{\lfloor \tau t\rfloor-1}
  \psi_{\lfloor \tau t\rfloor-s}^\tau
  \left(
    M_s^{\tau,+}+\beta M_s^{\tau,-}
  \right).
\end{align*}
Using
\[
  Z_{s/\tau}^{\tau,\pm}
  =
  \sqrt{\frac{1-a_\tau}{\tau^\alpha\mu}}\,M_s^{\tau,\pm},
\]
we can rewrite $T_3^\tau$ as
\begin{equation}\label{eq:T3-rewrite-prop2}
  T_3^\tau(t)
  =
  \frac{\tau(1-a_\tau)}{\sqrt{\gamma\mu(1+\beta)^2}}
  \sum_{s=1}^{\lfloor \tau t\rfloor-1}
  \frac1{\tau}\psi_{\lfloor \tau t\rfloor-s}^\tau
  \left(
    Z_{s/\tau}^{\tau,+}+\beta Z_{s/\tau}^{\tau,-}
  \right).
\end{equation}

\medskip
\noindent
\textit{Step 3: passage to the limit in the decomposition.}

Recall the rescaled renewal density
\[
  \mathsf S^\tau(u):=
  \frac{1-a_\tau}{a_\tau}\tau\psi_{\lfloor \tau u\rfloor}^\tau,
  \qquad u\ge0,
\]
and Lemma~\ref{lem:kernel_convergence}, which states that
\[
  \mathsf S^\tau(u)\,du \Rightarrow f_{\alpha,\gamma}(u)\,du
\]
weakly on $[0,\infty)$, where
$f_{\alpha,\gamma}(u):=\gamma u^{\alpha-1}E_{\alpha,\alpha}(-\gamma u^\alpha)$
and 
$F_{\alpha,\gamma}(t):=\int_0^t f_{\alpha,\gamma}(u)\,du$.

Fix $t\in[0,T]$. We identify the limits of the three terms in
\eqref{eq:Lambda-decomp-prop2}.

First,
\[
  T_1^\tau(t)=(1-a_\tau)\frac{\lfloor \tau t\rfloor}{\tau}\to0,
\]
since $1-a_\tau=\gamma\tau^{-\alpha}\to0$.

Next, writing $j=\lfloor \tau t\rfloor-s$, we have
\begin{align*}
  \frac{1-a_\tau}{\tau}
  \sum_{s=1}^{\lfloor \tau t\rfloor-1}s\psi_{\lfloor \tau t\rfloor-s}^\tau
  =
  \frac{1-a_\tau}{\tau}
  \sum_{j=1}^{\lfloor \tau t\rfloor-1}
  (\lfloor \tau t\rfloor-j)\psi_j^\tau
 =
  \frac{a_\tau}{\tau}
  \sum_{j=1}^{\lfloor \tau t\rfloor-1}
  \mathsf S^\tau(j/\tau)
  \left(\frac{\lfloor \tau t\rfloor}{\tau}-\frac{j}{\tau}\right).
\end{align*}
Since the function
$u\mapsto (t-u)\mathbf 1_{[0,t]}(u)$
is bounded and continuous on $[0,\infty)$, the weak convergence of
$\mathsf S^\tau(u)\,du$ implies
\[
  \frac{1-a_\tau}{\tau}
  \sum_{s=1}^{\lfloor \tau t\rfloor-1}s\psi_{\lfloor \tau t\rfloor-s}^\tau
  \to
  \int_0^t (t-u)f_{\alpha,\gamma}(u)\,du
  =
  \int_0^t F_{\alpha,\gamma}(t-s)\,ds.
\]
Also,
$\frac{\lfloor \tau t\rfloor}{\tau}\to t$.
Hence,
\begin{equation}\label{eq:T2-limit-prop2}
  T_2^\tau(t)
  \to
  \int_0^t F_{\alpha,\gamma}(t-s)\,ds
  +
  \xi_0\left(
    t-\int_0^t F_{\alpha,\gamma}(t-s)\,ds
  \right).
\end{equation}

We now treat $T_3^\tau(t)$. Using \eqref{eq:T3-rewrite-prop2}, write
\begin{align*}
  T_3^\tau(t)
  &=
  \frac{a_\tau}{\sqrt{\gamma\mu(1+\beta)^2}}
  \sum_{s=1}^{\lfloor \tau t\rfloor-1}
  \frac1{\tau}
  \mathsf S^\tau\!\left(\frac{\lfloor \tau t\rfloor-s}{\tau}\right)
  \left(
    Z_{s/\tau}^{\tau,+}+\beta Z_{s/\tau}^{\tau,-}
  \right).
\end{align*}
Set
\[
  G^\tau(s):=Z_s^{\tau,+}+\beta Z_s^{\tau,-},
  \qquad
  G(s):=Z_s^++\beta Z_s^-.
\]
Then
\begin{align*}
  &T_3^\tau(t)
  -
  \frac1{\sqrt{\gamma\mu(1+\beta)^2}}
  \int_0^t f_{\alpha,\gamma}(t-s)G(s)\,ds
  \\
  &=
  \frac{a_\tau}{\sqrt{\gamma\mu(1+\beta)^2}}
  \sum_{s=1}^{\lfloor \tau t\rfloor-1}
  \frac1{\tau}
  \mathsf S^\tau\!\left(\frac{\lfloor \tau t\rfloor-s}{\tau}\right)
  \left(G^\tau(s/\tau)-G(s/\tau)\right)
  \\
  &\quad
  +
  \frac{1}{\sqrt{\gamma\mu(1+\beta)^2}}
  \left[
    a_\tau
    \sum_{s=1}^{\lfloor \tau t\rfloor-1}
    \frac1{\tau}
    \mathsf S^\tau\!\left(\frac{\lfloor \tau t\rfloor-s}{\tau}\right)
    G(s/\tau)
    -
    \int_0^t f_{\alpha,\gamma}(t-s)G(s)\,ds
  \right]
  \\
  &=:I_1^\tau(t)+I_2^\tau(t).
\end{align*}
For the first term, since $a_\tau\le1$ and
\[
  \frac1{\tau}\sum_{s=1}^{\lfloor \tau t\rfloor-1}
  \mathsf S^\tau\!\left(\frac{\lfloor \tau t\rfloor-s}{\tau}\right)
  \le
  \int_0^\infty \mathsf S^\tau(u)\,du
  =1,
\]
we get
\[
  |I_1^\tau(t)|
  \le
  \frac{1}{\sqrt{\gamma\mu(1+\beta)^2}}
  \sup_{0\le r\le T}|G^\tau(r)-G(r)|
  \to0
\]
almost surely by \eqref{eq:prop2-skorokhod}.

For the second term, the function
$u\mapsto G(t-u)\mathbf 1_{[0,t]}(u)$
is bounded and continuous on $[0,\infty)$ because $G$ is continuous. Hence the weak
convergence of $\mathsf S^\tau(u)\,du$ implies that
$I_2^\tau(t)\to0$.
Therefore,
\begin{equation}\label{eq:T3-limit-prop2}
  T_3^\tau(t)
  \to
  \frac1{\sqrt{\gamma\mu(1+\beta)^2}}
  \int_0^t
  f_{\alpha,\gamma}(t-s)\left(Z_s^++\beta Z_s^-\right)\,ds.
\end{equation}

Combining \eqref{eq:Lambda-decomp-prop2}, \eqref{eq:T2-limit-prop2},
\eqref{eq:T3-limit-prop2}, and Lemma~\ref{lem:tightness}, which yields
\[
  \sup_{t\in[0,T]}|\Lambda_t^{\tau,+}-Y_t^{\tau,+}|\to0
  \qquad\text{in probability},
\]
we conclude that for every fixed $t\in[0,T]$,
\begin{equation}\label{eq:Y-limit-eqn}
\begin{aligned}
  Y_t
  &=
  \int_0^t F_{\alpha,\gamma}(t-s)\,ds
  +
  \xi_0\left(
    t-\int_0^t F_{\alpha,\gamma}(t-s)\,ds
  \right)
  \\
  &\qquad
  +
  \frac1{\sqrt{\gamma\mu(1+\beta)^2}}
  \int_0^t
  f_{\alpha,\gamma}(t-s)\left(Z_s^++\beta Z_s^-\right)\,ds.
\end{aligned}
\end{equation}

\medskip
\noindent
\textit{Step 4: $Y$ is absolutely continuous and its density.}

Define
\[
  c_{\beta,\gamma,\mu}:=\frac1{\sqrt{\gamma\mu(1+\beta)^2}},
  \qquad
  H_t:=Z_t^++\beta Z_t^-.
\]
Then \eqref{eq:Y-limit-eqn} becomes
\[
  Y_t
  =
  \int_0^t F_{\alpha,\gamma}(t-s)\,ds
  +
  \xi_0\left(
    t-\int_0^t F_{\alpha,\gamma}(t-s)\,ds
  \right)
  +
  c_{\beta,\gamma,\mu}\int_0^t f_{\alpha,\gamma}(t-s)H_s\,ds.
\]
Since $H$ is a continuous semimartingale and $F_{\alpha,\gamma}(0)=0$,
integration by parts yields
\[
  \int_0^t f_{\alpha,\gamma}(t-s)H_s\,ds
  =
  \int_0^t F_{\alpha,\gamma}(t-s)\,dH_s.
\]
Now $f_{\alpha,\gamma}\in L^2([0,T])$ because $\alpha>1/2$, and by
Proposition~\ref{proposition1},
\[
  [H]_t=[Z^++\beta Z^-]_t=(1+\beta^2)Y_t,
\]
which is finite on $[0,T]$. Hence,
\[
  \int_0^T\int_0^t f_{\alpha,\gamma}(t-s)^2\,d[H]_s\,dt<\infty
  \qquad\text{a.s.},
\]
so that stochastic Fubini's theorem applies. Therefore,
\begin{align*}
  \int_0^t F_{\alpha,\gamma}(t-s)\,dH_s
  =
  \int_0^t\int_s^t f_{\alpha,\gamma}(r-s)\,dr\,dH_s
 =
  \int_0^t\int_0^r f_{\alpha,\gamma}(r-s)\,dH_s\,dr.
\end{align*}
Substituting this into the previous display, we obtain
\[
  Y_t=\int_0^t v_s\,ds,
\]
where
\begin{equation}\label{eq:v-density-pre}
  v_t
  =
  \xi_0\left(1-F_{\alpha,\gamma}(t)\right)
  +
  F_{\alpha,\gamma}(t)
  +
  c_{\beta,\gamma,\mu}
  \int_0^t
  f_{\alpha,\gamma}(t-s)\,dH_s.
\end{equation}

\medskip
\noindent
\textit{Step 5: Brownian representation of $Z^+$ and $Z^-$.}

By Proposition~\ref{proposition1}, every limit point satisfies
\[
  [Z^+,Z^+]_t=Y_t,\qquad [Z^-,Z^-]_t=Y_t,\qquad [Z^+,Z^-]_t=0.
\]
Since $Y_t=\int_0^t v_s\,ds$ is absolutely continuous, the quadratic variations
of $Z^+$ and $Z^-$ are absolutely continuous. Hence, by the representation theorem
for continuous local martingales with absolutely continuous quadratic variation,
there exist standard Brownian motions $B^1$ and $B^2$ such that
\[
  Z_t^+=\int_0^t\sqrt{v_s}\,dB_s^1,
  \qquad
  Z_t^-=\int_0^t\sqrt{v_s}\,dB_s^2.
\]
Moreover, since $[Z^+,Z^-]\equiv0$, the two-dimensional Brownian motion
$(B^1,B^2)$ has covariance matrix equal to the identity, and therefore
$B^1$ and $B^2$ are independent.

Thus
\[
  dH_s
  =
  d(Z_s^++\beta Z_s^-)
  =
  \sqrt{v_s}\,\left(dB_s^1+\beta\,dB_s^2\right).
\]
Define
\[
  W_t^V:=\frac{B_t^1+\beta B_t^2}{\sqrt{1+\beta^2}}.
\]
Then $W^V$ is a standard Brownian motion, and
\[
  dH_s=\sqrt{1+\beta^2}\,\sqrt{v_s}\,dW_s^V.
\]
Substituting this into \eqref{eq:v-density-pre}, we obtain
\[
  v_t
  =
  \xi_0\left(1-F_{\alpha,\gamma}(t)\right)
  +
  F_{\alpha,\gamma}(t)
  +
  \sqrt{\frac{1+\beta^2}{\gamma\mu(1+\beta)^2}}
  \int_0^t
  f_{\alpha,\gamma}(t-s)\sqrt{v_s}\,dW_s^V.
\]

Therefore
\[
  Y_t=\int_0^t v_s\,ds,
  \qquad
  Z_t^+=\int_0^t\sqrt{v_s}\,dB_s^1,
  \qquad
  Z_t^-=\int_0^t\sqrt{v_s}\,dB_s^2,
\]
with
\[
  v_t
  =
  \xi_0\left(1-F_{\alpha,\gamma}(t)\right)
  +
  F_{\alpha,\gamma}(t)
  +
  \sqrt{\frac{1+\beta^2}{\gamma\mu(1+\beta)^2}}
  \int_0^t
  f_{\alpha,\gamma}(t-s)\sqrt{v_s}\,dW_s^V.
\]
This proves the claimed representation. The H\"older regularity of $v$
follows from Proposition~6.2 of \cite{el2019characteristic}. This completes
the proof of Proposition~\ref{prop2}.
%%%%%%%%%%%%%%%%%%%%%%%%%%%%%%%%%%%%%%%%

\subsection{Proof of Theorem~\ref{theorem1}}\label{app:inar-moment}

From Proposition~\ref{proposition1}, we have obtained $C$-tightness of $(\boldsymbol\Lambda^{\tau},\mathbf Y^{\tau},\mathbf Z^{\tau})_{\tau\ge1}$, so every subsequence admits a further subsequence converging in law to some limit point $(\boldsymbol\Lambda,\mathbf Y,\mathbf Z)$. The same proposition also identifies the limiting quadratic covariation:
\[
  [\mathbf Z,\mathbf Z]=\mathrm{diag}(\mathbf Y),
  \qquad
  [Z^+,Z^-]_t=0.
\]

Proposition~\ref{prop2} then shows that every such limit point admits a Brownian representation with independent components and that its rate process $v$, defined by $Y_t=\int_0^t v_s\,ds$, satisfies the stochastic Volterra equation \eqref{finalequation1}. Rewriting that equation in the rough Heston form \eqref{eq:vt} and invoking the well-posedness and uniqueness results for the rough Heston variance equation from \cite{el2019characteristic} identifies the limit uniquely.

Hence all subsequential limit points coincide, and the whole sequence converges in law to the unique limit described in Theorem~\ref{theorem1}. This completes the proof of Theorem~\ref{theorem1}.

%%%%%%%%%%%%%%%%%%%%%%%%%%%%%%%%%

\subsection{Proof of Theorem~\ref{maintheorem}}\label{sec:proofoftheorem2}
The proof follows the logic of Corollary 2.1 of \cite{el2019characteristic}, leveraging the convergence result established in our Theorem \ref{theorem1} and applying the continuous mapping theorem. We outline the key steps.

First, recall the definition of the properly rescaled microscopic price process $P_t^\tau$:
\begin{align*}
    {P_t^{\tau}}
    = {\sqrt{\frac{\theta}{2}}\sqrt{\frac{1-a_{\tau}}{\tau^\alpha\mu}}\left( N_{\lfloor t\tau\rfloor}^{\tau,+}-N_{\lfloor t\tau\rfloor}^{\tau,-} \right)-\frac\theta2\frac{1-a_{\tau}}{\tau^\alpha\mu}N_{\lfloor t\tau\rfloor}^{\tau,+}} = {\sqrt{\frac{\theta}{2}}\left(Z_t^{\tau,+} - Z_t^{\tau,-}\right) - \frac{\theta}{2} Y_t^{\tau,+}.}
\end{align*}
The mapping from the process triplet $(Z_t^{\tau,+}, Z_t^{\tau,-}, Y_t^{\tau,+})$ to $P_t^\tau$ is a continuous linear combination. Theorem \ref{theorem1} establishes the joint convergence in law of these processes to their limits $(Z_t^+, Z_t^-, Y_t^+)$ in the Skorokhod topology. By the continuous mapping theorem, $P_t^\tau$ converges in law to a process $P_t$ given by the same linear combination of the limits:
\begin{align*}
 P_t = {\sqrt{\frac{\theta}{2}}(Z_t^+ - Z_t^-) - \frac{\theta}{2} Y_t^+} = \sqrt{\frac{\theta}{2}} \left( \int_0^t \sqrt{v_s} dB_s^1 - \int_0^t \sqrt{v_s} dB_s^2 \right) - \frac{\theta}{2} \int_0^t v_s ds,
\end{align*}
where we have substituted the integral representations of the limit processes from Theorem \ref{theorem1}.

Now define the variance process $V_t:=\theta v_t$ and the Brownian motion
\[
  W_t^X:=\frac{1}{\sqrt{2}}(B_t^1-B_t^2).
\]
Then
\[
  \sqrt{\frac{\theta}{2}}\left(dB_t^1-dB_t^2\right)=\sqrt{\theta}\,dW_t^X,
\]
so that
\[
  \sqrt{\frac{\theta}{2}}\int_0^t \sqrt{v_s}\,(dB_s^1-dB_s^2)
  =\sqrt{\theta}\int_0^t \sqrt{v_s}\,dW_s^X
  =\int_0^t \sqrt{\theta v_s}\,dW_s^X
  =\int_0^t \sqrt{V_s}\,dW_s^X.
\]
Hence,
\[
  P_t = \int_0^t \sqrt{V_s} dW_s^X - \frac{1}{2} \int_0^t V_s ds.
\]
This is precisely the log-price process in the rough Heston model under the risk-neutral measure.

To obtain the SDE for the variance process $V_t$, we simply multiply the SDE for $v_t$ from Theorem \ref{theorem1} by the constant $\theta$. This yields equation \eqref{eq:V}.

Finally, we compute the correlation between the Brownian motions $W^X$ and $W^V$. Using the properties of quadratic covariation and the definitions of $W_t^X$ and $W_t^V = (B_t^1 + \beta B_t^2) / \sqrt{1+\beta^2}$, we have:
\begin{align*}
  &d\left\langle W^X, W^V \right\rangle_t 
  \\
  &= d\left\langle \frac{1}{\sqrt{2}}\left(B^1 - B^2\right), \frac{1}{\sqrt{1+\beta^2}}\left(B^1 + \beta B^2\right) \right\rangle_t \\
  &= \frac{1}{\sqrt{2(1+\beta^2)}} \left( d\left\langle B^1, B^1 \right\rangle_t + \beta d\left\langle B^1, B^2 \right\rangle_t - d\left\langle B^2, B^1 \right\rangle_t - \beta d\left\langle B^2, B^2 \right\rangle_t \right).
\end{align*}
Since $B^1$ and $B^2$ are independent standard Brownian motions, we have $d\langle B^1, B^1 \rangle_t = d\langle B^2, B^2 \rangle_t = dt$ and $d\langle B^1, B^2 \rangle_t = 0$. This simplifies to:
\[
    d\left\langle W^X,W^V \right\rangle_t=\frac{1-\beta}{\sqrt{2(1+\beta^2)}}dt.
\]
This completes the proof.

%%%%%%%%%%%%%%%%%%%%%%%%%%%%%%%%%%%%%%%%%%%%%%%%%%%%%%%%%%%%

\section{Technical Results for the Weak-Error Analysis}
\label{supp:weak-error-proofs}
%%%%%%%%%%%%%%%%%%%%%%%%%%%%%%%%%%%%%

For $\alpha\in(0,1)$ and $f\in C^\alpha([0,T])$, we define
\begin{equation}\label{eq:holder-seminorm}
  [f]_{C^\alpha([0,T])}
  :=
  \sup_{\substack{x,y\in[0,T]\\x\neq y}}
  \frac{|f(x)-f(y)|}{|x-y|^\alpha},
\end{equation}
and
\begin{equation}\label{eq:holder-norm}
  \|f\|_{C^\alpha([0,T])}
  :=
  \|f\|_\infty+[f]_{C^\alpha([0,T])}.
\end{equation}

\subsection{Proof of Proposition~\ref{prop:baseline-riemann}}\label{app:baseline-riemann}

\begin{proof}
We use the explicit form of $\hat\mu_\tau(m)$ from Definition~\ref{def:params}. Write
\[
  \overline\varphi(n):=\sum_{k=n+1}^{\infty}\varphi_k,\qquad n\ge0.
\]
For the kernel chosen in Definition~\ref{def:params}, a direct telescoping computation gives
\begin{equation}\label{eq:tail-phi-explicit-final}
  \overline\varphi(0)=1,
  \qquad
  \overline\varphi(n)=\frac{1}{\Gamma(1-\alpha)}n^{-\alpha},
  \qquad n\ge1.
\end{equation}
Moreover,
\[
  1-\sum_{s=1}^{m-1}\varphi_s^\tau
  =1-a_\tau\sum_{s=1}^{m-1}\varphi_s
  =\varepsilon_\tau+a_\tau\overline\varphi(m-1).
\]
Substituting this into the definition of $\hat\mu_\tau(m)$ yields
\begin{align}
  \hat\mu_\tau(m)
  &=\mu_\tau+\xi_0\mu_\tau\left(
      \frac{\varepsilon_\tau+a_\tau\overline\varphi(m-1)}{\varepsilon_\tau}
      -a_\tau\sum_{s=1}^{m-1}\varphi_s
    \right)
  \nonumber\\
  &=\mu_\tau
    +\xi_0\mu_\tau\varepsilon_\tau
    +\xi_0\mu_\tau a_\tau\left(\varepsilon_\tau^{-1}+1\right)\overline\varphi(m-1).
  \label{eq:muhat-decomp-final}
\end{align}
Multiplying by $\varepsilon_\tau/\mu$ and using $\mu_\tau=\mu\tau^{\alpha-1}$ and $\varepsilon_\tau=\gamma\tau^{-\alpha}$, we obtain
\begin{equation}\label{eq:muhat-scaled-decomp-final}
  \frac{\varepsilon_\tau}{\mu}\hat\mu_\tau(m)
  = \frac{\gamma}{\tau}
    + \frac{\xi_0\gamma^2}{\tau^{\alpha+1}}
    + \xi_0 a_\tau\tau^{\alpha-1}\overline\varphi(m-1)
    + \frac{\xi_0\gamma a_\tau}{\tau}\overline\varphi(m-1).
\end{equation}
Hence,
\begin{align}
  &\frac{\varepsilon_\tau}{\mu}
    \sum_{m=1}^{N_\tau}\hat\mu_\tau(m)
    f\!\left(T-\frac{m}{\tau}\right)
  \nonumber\\
  &=
    \frac{\gamma}{\tau}
    \sum_{m=1}^{N_\tau}f\!\left(T-\frac{m}{\tau}\right)
    +
    \xi_0 a_\tau\tau^{\alpha-1}
    \sum_{m=1}^{N_\tau}\overline\varphi(m-1)
    f\!\left(T-\frac{m}{\tau}\right)
    +R_\tau(f),
  \label{eq:baseline-split-final}
\end{align}
where, by \eqref{eq:tail-phi-explicit-final},
\begin{align}
  |R_\tau(f)|
  &\le C_T\|f\|_{\infty}
  \left(
    \tau^{-\alpha}
    +
    \frac1\tau\sum_{m=2}^{N_\tau}(m-1)^{-\alpha}
  \right)
  \nonumber\\
  &\le C_T\|f\|_{\infty}
  \left(
    \tau^{-\alpha}
    +\tau^{-1}N_\tau^{1-\alpha}
  \right)
  \le C_T\|f\|_{\infty}\tau^{-\alpha}.
  \label{eq:Rtau-bound-final}
\end{align}

We now treat the two main terms in \eqref{eq:baseline-split-final}.

\medskip
\noindent
\emph{Step 1: regular Riemann sum.}
Set
\[
x_m:=T-\frac{m}{\tau},\qquad m=1,\dots,N_\tau.
\]
Then
\[
\frac{1}{\tau}\sum_{m=1}^{N_\tau}f\!\left(T-\frac{m}{\tau}\right)
=
\sum_{m=1}^{N_\tau}\int_{x_m}^{x_m+1/\tau} f(x_m)\,ds.
\]
Therefore,
\begin{align*}
&\left|
\frac{1}{\tau}\sum_{m=1}^{N_\tau}f\!\left(T-\frac{m}{\tau}\right)
-\int_0^T f(s)\,ds
\right|
\\
&\le
\sum_{m=1}^{N_\tau}
\int_{x_m}^{x_m+1/\tau}|f(x_m)-f(s)|\,ds
+
\left|\int_{N_\tau/\tau}^{T}f(s)\,ds\right|.
\end{align*}
Since $f\in C^\alpha([0,T])$, by \eqref{eq:holder-seminorm}, for
$s\in[x_m,x_m+1/\tau]$,
\[
  |f(x_m)-f(s)|
  \le [f]_{C^\alpha([0,T])}|s-x_m|^\alpha
  \le [f]_{C^\alpha([0,T])}\tau^{-\alpha},
\]
and hence,
\[
\sum_{m=1}^{N_\tau}
\int_{x_m}^{x_m+1/\tau}|f(x_m)-f(s)|\,ds
\le
N_\tau\cdot \tau^{-1-\alpha}[f]_{C^\alpha([0,T])}
\le C_T\|f\|_{C^\alpha([0,T])}\tau^{-\alpha}.
\]
Also, since $0\le T-N_\tau/\tau\le 1/\tau$,
\[
\left|\int_{N_\tau/\tau}^{T}f(s)\,ds\right|
\le
\frac{1}{\tau}\|f\|_\infty
\le
C_T\|f\|_{C^\alpha([0,T])}\tau^{-1}
\le
C_T\|f\|_{C^\alpha([0,T])}\tau^{-\alpha},
\]
because $\alpha\in(0,1)$. Thus
\begin{equation}\label{eq:regular-riemann-final}
  \left|
    \frac{\gamma}{\tau}\sum_{m=1}^{N_\tau}f\!\left(T-\frac{m}{\tau}\right)
    -\gamma\int_0^T f(s)\,ds
  \right|
  \le C_T\|f\|_{C^{\alpha}([0,T])}\tau^{-\alpha}.
\end{equation}

\medskip
\noindent
\emph{Step 2: weakly singular Riemann sum.}
Set
\[
  g(u):=f(T-u),\qquad u\in[0,T].
\]
Then $g\in C^{\alpha}([0,T])$ and
\[
  \|g\|_{C^{\alpha}([0,T])}\le C\|f\|_{C^{\alpha}([0,T])}.
\]
By \eqref{eq:tail-phi-explicit-final}, the term $m=1$ must be treated separately, since
\[
  \overline\varphi(0)=1,
  \qquad
  \overline\varphi(j)=\frac{1}{\Gamma(1-\alpha)}j^{-\alpha},
  \qquad j\ge1.
\]
Hence,
\begin{equation}\label{eq:singular-sum-decomp-final}
\begin{aligned}
  \tau^{\alpha-1}\sum_{m=1}^{N_\tau}\overline\varphi(m-1)
  f\!\left(T-\frac{m}{\tau}\right)
  =
  \tau^{\alpha-1}g\!\left(\frac1\tau\right)
  +
  \frac{1}{\Gamma(1-\alpha)}\frac1\tau
  \sum_{j=1}^{N_\tau-1}\left(\frac{j}{\tau}\right)^{-\alpha}
  g\!\left(\frac{j+1}{\tau}\right).
\end{aligned}
\end{equation}
We compare this with
\[
  I^{1-\alpha}f(T)
  =
  \frac{1}{\Gamma(1-\alpha)}\int_0^T u^{-\alpha}g(u)\,du.
\]

We split the error into the first cell and the remaining cells.

\medskip
\noindent
\emph{First cell.}
Write
\begin{align*}
  \Delta_{0,\tau}
  &:=
  \tau^{\alpha-1}g\!\left(\frac1\tau\right)
  -
  \frac{1}{\Gamma(1-\alpha)}
  \int_0^{1/\tau}u^{-\alpha}g(u)\,du.
\end{align*}
Using
\[
g\!\left(\frac1\tau\right)=g(0)+\mathcal O\!\left(\tau^{-\alpha}\|g\|_{C^\alpha([0,T])}\right),
\]
we obtain
\begin{align*}
  \Delta_{0,\tau}
  &=
  g(0)\left[
    \tau^{\alpha-1}
    -
    \frac{1}{\Gamma(1-\alpha)}
    \int_0^{1/\tau}u^{-\alpha}\,du
  \right]
  +\mathcal O\left(\tau^{-1}\|g\|_{C^\alpha([0,T])}\right).
\end{align*}
Since
\[
\int_0^{1/\tau}u^{-\alpha}\,du
  =
  \frac{1}{1-\alpha}\tau^{-(1-\alpha)}
\]
and
$\Gamma(2-\alpha)=(1-\alpha)\Gamma(1-\alpha)$,
we can rewrite the bracket as
\[
\tau^{-(1-\alpha)}
  \left(
    1-\frac{1}{\Gamma(2-\alpha)}
  \right).
\]
Therefore,
\begin{equation}\label{eq:first-cell-mainterm}
  \Delta_{0,\tau}
  =
  \left(
    1-\frac{1}{\Gamma(2-\alpha)}
  \right)\tau^{-(1-\alpha)}g(0)
  +
  \mathcal O\left(\tau^{-1}\|g\|_{C^\alpha([0,T])}\right).
\end{equation}
Since $\alpha\in(1/2,1)$, we have $\tau^{-1}\le\tau^{-\alpha}$ for large $\tau$. Moreover,
$1-\frac{1}{\Gamma(2-\alpha)}
  \to0$
as $\alpha\uparrow1^-$.

\medskip
\noindent
\emph{Remaining cells.}
For $1\le j\le N_\tau-1$, set
\[
  E_j:=
  \left|
    \frac1\tau\left(\frac{j}{\tau}\right)^{-\alpha}g\!\left(\frac{j+1}{\tau}\right)
    -
    \int_{j/\tau}^{(j+1)/\tau}u^{-\alpha}g(u)\,du
  \right|.
\]
Then
\begin{align*}
  E_j
  &\le
  \left|
    \frac1\tau\left(\frac{j}{\tau}\right)^{-\alpha}g\!\left(\frac{j+1}{\tau}\right)
    -
    g\!\left(\frac{j+1}{\tau}\right)
    \int_{j/\tau}^{(j+1)/\tau}u^{-\alpha}\,du
  \right|
  \\
  &\quad+
  \int_{j/\tau}^{(j+1)/\tau}
  u^{-\alpha}
  \left|g\!\left(\frac{j+1}{\tau}\right)-g(u)\right|\,du.
\end{align*}
Using the mean-value theorem for $u^{-\alpha}$ and the $\alpha$-H\"older continuity of $g$,
\[
  E_j
  \le
  C_T\|g\|_{\infty}\tau^{\alpha-2}j^{-\alpha-1}
  +
  C_T\|g\|_{C^\alpha([0,T])}\tau^{-1}j^{-\alpha}.
\]
Summing over $1\le j\le N_\tau-1$, we obtain
\begin{align}
  \sum_{j=1}^{N_\tau-1}E_j
  &\le
  C_T\|f\|_{C^\alpha([0,T])}
  \left(
    \tau^{\alpha-2}\sum_{j=1}^{N_\tau-1}j^{-\alpha-1}
    +
    \tau^{-1}\sum_{j=1}^{N_\tau-1}j^{-\alpha}
  \right)
  \nonumber\\
  &\le
  C_T\|f\|_{C^\alpha([0,T])}
  \left(
    \tau^{\alpha-2}
    +\tau^{-1}N_\tau^{1-\alpha}
  \right)
  \le
  C_T\|f\|_{C^\alpha([0,T])}\tau^{-\alpha}.
  \label{eq:Ej-sum-bound}
\end{align}
The terminal interval $[N_\tau/\tau,T]$ has length at most $\tau^{-1}$, so
\[
  \int_{N_\tau/\tau}^T u^{-\alpha}|g(u)|\,du
  \le C_T\|f\|_\infty \tau^{-1}
  \le C_T\|f\|_{C^\alpha([0,T])}\tau^{-\alpha}.
\]

Combining \eqref{eq:first-cell-mainterm}, \eqref{eq:Ej-sum-bound}, and the tail estimate gives
\begin{equation}\label{eq:singular-riemann-refined}
\begin{aligned}
  &\left|
    \tau^{\alpha-1}\sum_{m=1}^{N_\tau}\overline\varphi(m-1)
    f\!\left(T-\frac{m}{\tau}\right)
    -
    I^{1-\alpha}f(T)
  \right|
  \\
  &\le
  C_T\|f\|_{C^{\alpha}([0,T])}\tau^{-\alpha}
  +
  C_T\left|1-\frac{1}{\Gamma(2-\alpha)}\right|
  \|f\|_{C^{\alpha}([0,T])}\tau^{-(1-\alpha)}.
\end{aligned}
\end{equation}

It remains to reinsert the prefactor $a_\tau$. First,
\begin{align*}
\tau^{\alpha-1}\sum_{m=1}^{N_\tau}\overline\varphi(m-1)
&=
\tau^{\alpha-1}\left(
1+\frac{1}{\Gamma(1-\alpha)}\sum_{j=1}^{N_\tau-1}j^{-\alpha}
\right)
\\
&\le
C_T\left(\tau^{\alpha-1}+\tau^{\alpha-1}N_\tau^{1-\alpha}\right)
\le C_T.
\end{align*}
Therefore,
\begin{align*}
  &\left|
    \xi_0 a_\tau\tau^{\alpha-1}
    \sum_{m=1}^{N_\tau}\overline\varphi(m-1)
    f\!\left(T-\frac{m}{\tau}\right)
    -
    \xi_0\tau^{\alpha-1}
    \sum_{m=1}^{N_\tau}\overline\varphi(m-1)
    f\!\left(T-\frac{m}{\tau}\right)
  \right|
  \\
  &\le
  |\xi_0|\,|a_\tau-1|\,\tau^{\alpha-1}
  \sum_{m=1}^{N_\tau}\overline\varphi(m-1)
  \left|f\!\left(T-\frac{m}{\tau}\right)\right|
\le  C_T\|f\|_{\infty}\tau^{-\alpha}.
\end{align*}
Together with \eqref{eq:singular-riemann-refined}, this yields
\begin{equation}\label{eq:singular-riemann-final}
\begin{aligned}
  &\left|
    \xi_0 a_\tau\tau^{\alpha-1}
    \sum_{m=1}^{N_\tau}\overline\varphi(m-1)
    f\!\left(T-\frac{m}{\tau}\right)
    -
    \xi_0 I^{1-\alpha}f(T)
  \right|
  \\
  &\le
  C_T\|f\|_{C^{\alpha}([0,T])}\tau^{-\alpha}
  +
  C_T\left|1-\frac{1}{\Gamma(2-\alpha)}\right|
  \|f\|_{C^{\alpha}([0,T])}\tau^{-(1-\alpha)}.
\end{aligned}
\end{equation}

Finally, combining \eqref{eq:baseline-split-final}, \eqref{eq:Rtau-bound-final},
\eqref{eq:regular-riemann-final}, and \eqref{eq:singular-riemann-final}, we obtain
\begin{align*}
  &\left|
    \frac{\varepsilon_\tau}{\mu}
    \sum_{m=1}^{N_\tau}\hat\mu_\tau(m)
    f\!\left(T-\frac{m}{\tau}\right)
    -
    \left[
      \gamma\int_0^T f(s)\,ds
      +
      \xi_0 I^{1-\alpha}f(T)
    \right]
  \right|\\
  \le&
  C_T\|f\|_{C^{\alpha}([0,T])}\tau^{-\alpha}
  +
  C_{\mathrm{base}}(\alpha)\|f\|_{C^{\alpha}([0,T])}\tau^{-(1-\alpha)},
\end{align*}
with
\[
  C_{\mathrm{base}}(\alpha)
  :=
  C_T\left|1-\frac{1}{\Gamma(2-\alpha)}\right|.
\]
In particular,
$C_{\mathrm{base}}(\alpha)\to0$
as $\alpha\uparrow1^-$.
This proves \eqref{eq:baseline-riemann}.
\end{proof}
%%%%%%%%%%%%%%%%%%%%%%%%%%%%%%%%%%%%%%%%%%%%%%%%%%%%%%%%
\subsection{Proof of Lemma~\ref{lem:phi-expansion}}\label{app:phi-expansion}

\begin{proof}
Write
\[
  a_n:=\frac{1}{\Gamma(1-\alpha)}\frac{1}{n^\alpha},
  \qquad n\ge1.
\]
From Definition~\ref{def:params}, for $n\ge2$,
$\varphi_n=a_{n-1}-a_n$
and $\varphi_1=1-a_1$. Hence,
\begin{align*}
\sum_{n\ge2}\varphi_n e^{-sn}
&=
\sum_{n\ge2}(a_{n-1}-a_n)e^{-sn}\\
&=
e^{-s}\sum_{m\ge1}a_m e^{-sm}-\sum_{n\ge2}a_n e^{-sn}
=
-\left(1-e^{-s}\right)\sum_{m\ge1}a_m e^{-sm}+a_1 e^{-s}.
\end{align*}
Adding $\varphi_1 e^{-s}=(1-a_1)e^{-s}$ gives the exact identity
\[
\hat\varphi(s)
=
e^{-s}-\left(1-e^{-s}\right)A(s),
\qquad
A(s):=\sum_{m\ge1}a_m e^{-sm}
=
\frac{1}{\Gamma(1-\alpha)}\mathrm{Li}_\alpha(e^{-s}),
\]
where $\mathrm{Li}_\alpha$ is the polylogarithm. For $0<\alpha<1$, the classical expansion
\cite{DLMF} yields
\[
\mathrm{Li}_\alpha(e^{-s})
=
\Gamma(1-\alpha)s^{\alpha-1}+\zeta(\alpha)+\mathcal O(s),
\qquad s\downarrow0.
\]
Therefore
\[
A(s)
=
s^{\alpha-1}+\frac{\zeta(\alpha)}{\Gamma(1-\alpha)}+\mathcal O(s),
\]
and, since
\[
1-e^{-s}=s-\frac{s^2}{2}+\mathcal O(s^3),
\qquad
e^{-s}=1-s+\frac{s^2}{2}+\mathcal O(s^3),
\]
we obtain
\begin{align*}
  1-\hat\varphi(s)
  =
  1-e^{-s}+\left(1-e^{-s}\right)A(s)
  =
  s^\alpha
  +\left(1+\frac{\zeta(\alpha)}{\Gamma(1-\alpha)}\right)s
  +\mathcal O(s^{1+\alpha})+\mathcal O(s^2).
\end{align*}
Because $1+\alpha<2$ for $\alpha\in(0,1)$, the $\mathcal O(s^2)$ term is absorbed into $\mathcal O(s^{1+\alpha})$. This proves \eqref{eq:phi-expansion-main}, with $c_1(\alpha)$ given by \eqref{eq:c1-alpha-main}.

Finally, as $\alpha\uparrow1^-$, we have
$\zeta(\alpha)\sim -\frac{1}{1-\alpha}$,
$\Gamma(1-\alpha)\sim \frac{1}{1-\alpha}$,
and therefore
$\frac{\zeta(\alpha)}{\Gamma(1-\alpha)}\to-1$.
This proves the limit statement in Lemma~\ref{lem:phi-expansion}.
\end{proof}

%%%%%%%%%%%%%%%%%%%%%%%%%%%%%%%%%%%%%





\subsection{Proof of Proposition~\ref{prop:exact-transform-weak}}\label{app:weak-euro-exact}
We use the branching--immigration representation of the Poisson INAR($\infty$) process.

Fix $z\in\mathcal D_\tau(T)$. Consider first a single ``$+$''-immigrant born at the current time. Its direct contribution to $zP_T^\tau$ is $\Theta_+^\tau(z)$, hence the factor $e^{\Theta_+^\tau(z)}$.

Now fix a lag $k\in\{1,\dots,n\}$. Because $q_k^{\tau,+}$ is the lag-$k$ coefficient associated with a ``$+$''-parent, the numbers of ``$+$''-children and ``$-$''-children produced at lag $k$ are independent Poisson random variables, both with mean $q_k^{\tau,+}$. Each such child starts an independent subcluster with remaining horizon $n-k$. Therefore, by the Poisson generating-function identity
$\E\left[s^{\mathrm{Poisson}(\lambda)}\right]=e^{\lambda(s-1)}$,
the lag-$k$ continuation factor equals
\[
  \exp\!\left(q_k^{\tau,+}(\mathcal M_{n-k}^{\tau,+}(z)-1)\right)\,
  \exp\!\left(q_k^{\tau,+}(\mathcal M_{n-k}^{\tau,-}(z)-1)\right)
  =
  \exp\!\left(q_k^{\tau,+}G_{n-k}^\tau(z)\right).
\]
Multiplying these independent continuation factors over $k=1,\dots,n$, and including the ancestor contribution $e^{\Theta_+^\tau(z)}$, we obtain
\[
  \mathcal M_n^{\tau,+}(z)
  =
  e^{\Theta_+^\tau(z)}
  \prod_{k=1}^n
  \exp\!\left(q_k^{\tau,+}G_{n-k}^\tau(z)\right)
  =
  \exp\!\left(
    \Theta_+^\tau(z)
    +\sum_{k=1}^n q_k^{\tau,+}G_{n-k}^\tau(z)
  \right),
\]
which is exactly \eqref{eq:Mplus-exact}.

The proof of \eqref{eq:Mminus-exact} is identical. For a single ``$-$''-immigrant, the lag-$k$ numbers of ``$+$''- and ``$-$''-children are independent Poisson random variables with common mean $q_k^{\tau,-}$, so the lag-$k$ continuation factor is
\[
  \exp\!\left(q_k^{\tau,-}(\mathcal M_{n-k}^{\tau,+}(z)-1)\right)\,
  \exp\!\left(q_k^{\tau,-}(\mathcal M_{n-k}^{\tau,-}(z)-1)\right)
  =
  \exp\!\left(q_k^{\tau,-}G_{n-k}^\tau(z)\right),
\]
which gives \eqref{eq:Mminus-exact}.

Now define $G_n^\tau(z):=\mathcal M_n^{\tau,+}(z)+\mathcal M_n^{\tau,-}(z)-2$. Summing \eqref{eq:Mplus-exact} and \eqref{eq:Mminus-exact} and subtracting $2$ yields the closed recursion \eqref{eq:G-rec-exact}.

Finally, at each time $m\in\{1,\dots,N_\tau\}$, the numbers of ``$+$''- and ``$-$''-immigrants are independent Poisson random variables with common mean $\hat\mu_\tau(m)$, and clusters generated by distinct immigrants are independent. Hence, again by the Poisson generating-function identity,
\[
  \Phi^\tau(z,T)
  =
  \prod_{m=1}^{N_\tau}
  \exp\!\left(
    \hat\mu_\tau(m)\left(\mathcal M_{N_\tau-m}^{\tau,+}(z)-1\right)
  \right)
  \exp\!\left(
    \hat\mu_\tau(m)\left(\mathcal M_{N_\tau-m}^{\tau,-}(z)-1\right)
  \right).
\]
Combining the two exponential factors and using the definition of $G_n^\tau(z)$, we obtain
\[
  \Phi^\tau(z,T)
  =
  \exp\!\left(
    \sum_{m=1}^{N_\tau}\hat\mu_\tau(m)\,
    \left(\mathcal M_{N_\tau-m}^{\tau,+}(z)+\mathcal M_{N_\tau-m}^{\tau,-}(z)-2\right)
  \right),
\]
which is exactly \eqref{eq:Phi-exact}. This completes the proof.

%%%%%%%%%%%%%%%%%%%%%%%%%%%%%%%%%%%%%%%%%%%%%%%%

\subsection{Uniform strip smallness}\label{app:weak-euro-second-order}

\begin{proposition}[Uniform strip smallness]
\label{prop:smallness-new33}
For every fixed $\eta\in\mathbb R$ and $R>0$, there exist $\tau_0\ge1$ and a
constant $C_{\eta,R,T}>0$ such that for all $\tau\ge\tau_0$,
with $\varepsilon_\tau:=1-a_\tau=\gamma\tau^{-\alpha}$,
\begin{equation}\label{eq:G-small-new33}
  \sup_{0\le n\le N_\tau}\sup_{z\in\Gamma_{\eta,R}}
  |G_n^\tau(z)|
  \le C_{\eta,R,T}\,\varepsilon_\tau,
\end{equation}
and
\begin{equation}\label{eq:AB-small-new33}
  \sup_{0\le n\le N_\tau}\sup_{z\in\Gamma_{\eta,R}}
  \left(|A_n^\tau(z)|+|B_n^\tau(z)|\right)
  \le C_{\eta,R,T}\,\varepsilon_\tau.
\end{equation}
In particular, after possibly enlarging the constant $C_{\eta,R,T}$,
\begin{equation}\label{eq:SD-small-new33}
  \sup_{0\le n\le N_\tau}\sup_{z\in\Gamma_{\eta,R}}
  \left(|S_n^\tau(z)|+|D_n^\tau(z)|\right)
  \le C_{\eta,R,T}\,\varepsilon_\tau.
\end{equation}
\end{proposition}

\begin{proof}
We proceed by induction on $n$.

\medskip
\noindent
\emph{Step 1: the case $n=0$.}
By definition,
\[
  A_0^\tau(z)=\Theta_+^\tau(z),\qquad B_0^\tau(z)=\Theta_-^\tau(z),
\]
hence
\[
  |A_0^\tau(z)|+|B_0^\tau(z)|
  \le |z|\,(2c_\tau+d_\tau)
  \le C\varepsilon_\tau
\]
uniformly on $\Gamma_{\eta,R}$, because $c_\tau\asymp \varepsilon_\tau$ and
$d_\tau\asymp \varepsilon_\tau^2$.

Since $A_0^\tau(z),B_0^\tau(z)\to0$ uniformly on $\Gamma_{\eta,R}$, the Taylor expansion gives
\[
  e^{A_0^\tau(z)}+e^{B_0^\tau(z)}-2
  =A_0^\tau(z)+B_0^\tau(z)
   +\frac12\left((A_0^\tau(z))^2+(B_0^\tau(z))^2\right)
   +\mathcal O(\varepsilon_\tau^3),
\]
uniformly on $\Gamma_{\eta,R}$. Therefore,
\[
  \sup_{z\in\Gamma_{\eta,R}}|G_0^\tau(z)|\le C\varepsilon_\tau.
\]
This proves \eqref{eq:G-small-new33} and \eqref{eq:AB-small-new33} at $n=0$.

\medskip
\noindent
\emph{Step 2: an analytic estimate for the nonlinear map.}
Define
\[
  F(a,b):=e^a+e^b-2.
\]
Since $F$ is analytic near $(0,0)$ and
\[
  F(a,b)=a+b+\frac12(a^2+b^2)+O(|a|^3+|b|^3),
\]
there exist constants $\delta>0$ and $C_\delta>0$ such that whenever
$|a|+|b|\le \delta$,
\begin{equation}\label{eq:F-est-new33}
  |F(a,b)|\le |a+b|+C_\delta\left(|a|^2+|b|^2\right).
\end{equation}

\medskip
\noindent
\emph{Step 3: induction step.}
Assume that for some $n\ge1$,
\begin{equation}\label{eq:ind-hyp-G-new33}
  \sup_{0\le j\le n-1}\sup_{z\in\Gamma_{\eta,R}}|G_j^\tau(z)|
  \le M\varepsilon_\tau
\end{equation}
for some constant $M>0$ to be fixed later.

Then from \eqref{eq:AB-def-new33},
\begin{align*}
  |A_n^\tau(z)|
  \le |\Theta_+^\tau(z)|+\sum_{k=1}^n q_k^{\tau,+}|G_{n-k}^\tau(z)|
  \le M_{\eta,R}(c_\tau+d_\tau)+\frac{a_\tau}{1+\beta}M\varepsilon_\tau
   \le C\varepsilon_\tau,
\end{align*}
and similarly
\begin{align*}
  |B_n^\tau(z)|
  \le |\Theta_-^\tau(z)|+\sum_{k=1}^n q_k^{\tau,-}|G_{n-k}^\tau(z)|
 \le M_{\eta,R}c_\tau+\frac{\beta a_\tau}{1+\beta}M\varepsilon_\tau
   \le C\varepsilon_\tau.
\end{align*}
Hence
\[
  \sup_{z\in\Gamma_{\eta,R}}\left(|A_n^\tau(z)|+|B_n^\tau(z)|\right)\le C\varepsilon_\tau.
\]
After enlarging $\tau_0$ if necessary, we may suppose that
\[
  |A_n^\tau(z)|+|B_n^\tau(z)|\le \delta
\]
for all $\tau\ge\tau_0$, all $0\le n\le N_\tau$, and all $z\in\Gamma_{\eta,R}$, so that
\eqref{eq:F-est-new33} applies.

Now use \eqref{eq:F-est-new33} and \eqref{eq:G-from-AB-new33}:
\[
  |G_n^\tau(z)|
  \le |A_n^\tau(z)+B_n^\tau(z)|
      +C\left(|A_n^\tau(z)|^2+|B_n^\tau(z)|^2\right).
\]
Since
\[
  A_n^\tau(z)+B_n^\tau(z)=S_n^\tau(z)
  =-zd_\tau+\sum_{k=1}^n\varphi_k^\tau\,G_{n-k}^\tau(z),
\]
we obtain
\begin{align*}
  |A_n^\tau(z)+B_n^\tau(z)|
  \le M_{\eta,R}d_\tau
      +\sum_{k=1}^n\varphi_k^\tau\,|G_{n-k}^\tau(z)|
      \le C\varepsilon_\tau^2+a_\tau M\varepsilon_\tau.
\end{align*}
Moreover, because $A_n^\tau(z),B_n^\tau(z)=\mathcal O(\varepsilon_\tau)$ uniformly,
\[
  |A_n^\tau(z)|^2+|B_n^\tau(z)|^2\le C\varepsilon_\tau^2.
\]
Hence,
\[
  |G_n^\tau(z)|\le a_\tau M\varepsilon_\tau + C\varepsilon_\tau^2
  = \left((1-\varepsilon_\tau)M + C\varepsilon_\tau\right)\varepsilon_\tau.
\]
Choose $M\ge C$. Then
$(1-\varepsilon_\tau)M + C\varepsilon_\tau \le M$,
so that
\[
  |G_n^\tau(z)|\le M\varepsilon_\tau.
\]
This closes the induction and proves \eqref{eq:G-small-new33}.

Finally, \eqref{eq:AB-small-new33} follows from the bounds on $A_n^\tau,B_n^\tau$ obtained
above. The bound \eqref{eq:SD-small-new33} then follows, after possibly enlarging the
constant, from
\[
  S_n^\tau=A_n^\tau+B_n^\tau,\qquad D_n^\tau=A_n^\tau-B_n^\tau.
\]
\end{proof}

%%%%%%%%%%%%%%%%%%%%%%%%%%%%%%%%%%%%%%%%%%%%%%%%

\subsection{Proof of Proposition~\ref{prop:reduced-recursion-new33}}\label{app:proof-reduced-recursion-new33}
Set
\[
  U_n^\tau(z):=\sum_{k=1}^n \varphi_k^\tau\,G_{n-k}^\tau(z).
\]
Then
\[
  A_n^\tau(z)=z(c_\tau-d_\tau)+\frac{1}{1+\beta}U_n^\tau(z),
\qquad
  B_n^\tau(z)=-zc_\tau+\frac{\beta}{1+\beta}U_n^\tau(z).
\]

By Proposition~\ref{prop:smallness-new33}, the quantities $A_n^\tau(z)$ and $B_n^\tau(z)$
are uniformly $\mathcal O(\varepsilon_\tau)$ on $\{0,\dots,N_\tau\}\times\Gamma_{\eta,R}$.
Hence, for $\tau$ large enough, they belong to a fixed complex neighborhood of the origin on
which the third-order Taylor formula for the exponential map is uniform:
\[
  e^x = 1+x+\frac{x^2}{2}+\rho(x),
  \qquad |\rho(x)|\le C|x|^3.
\]
Applying this with $x=A_n^\tau(z)$ and $x=B_n^\tau(z)$, and subtracting $2$, yields
\[
  G_n^\tau(z)
  =A_n^\tau(z)+B_n^\tau(z)
   +\frac12\left((A_n^\tau(z))^2+(B_n^\tau(z))^2\right)
   +r_n^\tau(z),
\]
where
\[
  |r_n^\tau(z)|\le C\left(|A_n^\tau(z)|^3+|B_n^\tau(z)|^3\right)\le C\varepsilon_\tau^3
\]
uniformly for $0\le n\le N_\tau$ and $z\in\Gamma_{\eta,R}$.

The linear term is
\[
  A_n^\tau(z)+B_n^\tau(z)=U_n^\tau(z)-zd_\tau.
\]
For the quadratic term, a direct computation gives
\begin{align*}
  \frac12\left((A_n^\tau(z))^2+(B_n^\tau(z))^2\right)
  &=
    \frac12\left((z(c_\tau-d_\tau))^2+z^2c_\tau^2\right) \\
  &\quad
    + \frac{z}{1+\beta}\left((c_\tau-d_\tau)-\beta c_\tau\right)U_n^\tau(z)
    + \frac{1+\beta^2}{2(1+\beta)^2}\left(U_n^\tau(z)\right)^2.
\end{align*}
Since $d_\tau=\mathcal O(\varepsilon_\tau^2)$ and $c_\tau=\mathcal O(\varepsilon_\tau)$, we have
\[
  \frac12\left((z(c_\tau-d_\tau))^2+z^2c_\tau^2\right)
  = c_\tau^2 z^2 + \mathcal O(\varepsilon_\tau^3),
\]
uniformly on $\Gamma_{\eta,R}$, and
\[
  \frac{z}{1+\beta}\left((c_\tau-d_\tau)-\beta c_\tau\right)
  = \frac{1-\beta}{1+\beta}\,c_\tau z + \mathcal O(\varepsilon_\tau^2)
\]
uniformly on $\Gamma_{\eta,R}$.
Moreover, by Proposition~\ref{prop:smallness-new33},
$U_n^\tau(z)=\mathcal O(\varepsilon_\tau)$
uniformly on $\{0,\dots,N_\tau\}\times\Gamma_{\eta,R}$, and thus the $\mathcal O(\varepsilon_\tau^2)$
coefficient error in the cross term contributes only $\mathcal O(\varepsilon_\tau^3)$.
Collecting the terms proves \eqref{eq:G-reduced-new33}--\eqref{eq:R-reduced-new33}.
%%%%%%%%%%%%%%%%%%%%%%%%%%%%%%%%%%%%%%%%%%%%%%%%%%

\subsection{Equivalent resolvent reformulation}\label{app:proof-barH}

The three lemmas in this subsection are algebraic rewrites. Lemma~\ref{lem:barH}
replaces the averaged state variable by the current state up to a lower-order
error. Lemma~\ref{lem:renewal-inversion} rewrites the discrete recursion in
resolvent form, and Lemma~\ref{lem:continuous-resolvent} gives the analogous
continuous resolvent form. Together they reduce the comparison to matching the
two resolvent kernels.

\begin{lemma}
\label{lem:barH}
With the convention that the empty sum is zero when $n=0$, uniformly for
$0\le n\le N_\tau$ and $z\in\Gamma_{\eta,R}$,
\begin{equation}\label{eq:H-barH-close}
  \left|
    H_n^\tau(z)-\sum_{k=1}^n\varphi_k H_{n-k}^\tau(z)
  \right|
  \le C_{\eta,R,T}\tau^{-\alpha}.
\end{equation}
Consequently,
\begin{equation}\label{eq:Htau-rec-local}
  H_n^\tau(z)
  =
  \sum_{k=1}^n\varphi_k^\tau H_{n-k}^\tau(z)
  +\tau^{-\alpha}F_z(H_n^\tau(z))
  +\widehat E_n^\tau(z),
\end{equation}
where
\begin{equation}\label{eq:Ehat-bound}
  \sup_{0\le n\le N_\tau}\sup_{z\in\Gamma_{\eta,R}}
  |\widehat E_n^\tau(z)|
  \le C_{\eta,R,T}\tau^{-2\alpha}.
\end{equation}
\end{lemma}

\begin{proof}

We first treat the case $n=0$. Since
\[
  H_0^\tau(z)=\frac{\mu}{\theta\varepsilon_\tau}G_0^\tau(z),
  \qquad
  G_0^\tau(z)=e^{\Theta_+^\tau(z)}+e^{\Theta_-^\tau(z)}-2,
\]
with $\Theta_+^\tau(z)=z(c_\tau-d_\tau)$ and $\Theta_-^\tau(z)=-zc_\tau$,
and since $c_\tau=\mathcal O(\varepsilon_\tau)$, $d_\tau=\mathcal O(\varepsilon_\tau^2)$,
the linear terms cancel and therefore
\[
  \sup_{z\in\Gamma_{\eta,R}}|G_0^\tau(z)|\le C_{\eta,R,T}\varepsilon_\tau^2.
\]
Hence,
\[
\sup_{z\in\Gamma_{\eta,R}}|H_0^\tau(z)|
  \le C_{\eta,R,T}\varepsilon_\tau
  \le C_{\eta,R,T}\tau^{-\alpha}.
\]
Thus \eqref{eq:H-barH-close} holds for $n=0$.

Now let $1\le n\le N_\tau$. By Proposition~\ref{prop:smallness-new33} and the rescaling
\eqref{eq:H-rescaled-def},
\[
  \sup_{0\le n\le N_\tau}\sup_{z\in\Gamma_{\eta,R}}
  |H_n^\tau(z)|\le C_{\eta,R,T}.
\]
Hence,
\[
  \sup_{0\le n\le N_\tau}\sup_{z\in\Gamma_{\eta,R}}
  \left|\sum_{k=1}^n\varphi_k H_{n-k}^\tau(z)\right|
  \le C_{\eta,R,T},
\]
since $\sum_{k\ge1}\varphi_k=1$.

Subtracting $\sum_{k=1}^n\varphi_kH_{n-k}^\tau(z)$ from both sides of
\eqref{eq:Htau-rec-main-again} gives
\begin{align*}
  &H_n^\tau(z)-\sum_{k=1}^n\varphi_k H_{n-k}^\tau(z)\\
  &=
  -(1-a_\tau)\sum_{k=1}^n\varphi_k H_{n-k}^\tau(z)
  +\tau^{-\alpha}
   F_z\!\left(\sum_{k=1}^n\varphi_k H_{n-k}^\tau(z)\right)
  +E_n^\tau(z).
\end{align*}
Using $1-a_\tau=\gamma\tau^{-\alpha}$, the boundedness of the convolution term,
the boundedness of $F_z$ on the relevant range, and \eqref{eq:En-bound-again},
we obtain \eqref{eq:H-barH-close}.

Since $F_z$ is quadratic in $x$ and the relevant arguments remain uniformly
bounded on $\Gamma_{\eta,R}$, $F_z$ is uniformly Lipschitz on the relevant
range. Hence, by \eqref{eq:H-barH-close},
\[
  \left|
    F_z\!\left(\sum_{k=1}^n\varphi_k H_{n-k}^\tau(z)\right)-F_z(H_n^\tau(z))
  \right|
  \le C_{\eta,R,T}
  \left|
    \sum_{k=1}^n\varphi_k H_{n-k}^\tau(z)-H_n^\tau(z)
  \right|
  \le C_{\eta,R,T}\tau^{-\alpha}.
\]
Multiplying by $\tau^{-\alpha}$ and using
\[
  a_\tau\sum_{k=1}^n\varphi_k H_{n-k}^\tau(z)
  =\sum_{k=1}^n\varphi_k^\tau H_{n-k}^\tau(z),
\]
we obtain \eqref{eq:Htau-rec-local} with
\[
  \widehat E_n^\tau(z)
  :=
  E_n^\tau(z)
  +\tau^{-\alpha}
  \left[
    F_z\!\left(\sum_{k=1}^n\varphi_k H_{n-k}^\tau(z)\right)-F_z(H_n^\tau(z))
  \right].
\]
The bound \eqref{eq:Ehat-bound} follows immediately.
\end{proof}

%%%%%%%%%%%%%%%%%%%%%%%%%%%%%%%%%%%%%%%%%%%%%%%%%%

% Appendix B.2 (previously a separate subsection; now part of the merged
% "Equivalent resolvent reformulation" subsection above)

\begin{lemma}[Renewal inversion]
\label{lem:renewal-inversion}
Let
\[
  \psi^\tau:=\sum_{m\ge1}(\varphi^\tau)^{*m}
  \qquad\text{with}\qquad
  \varphi^\tau=a_\tau\varphi
\]
be the renewal kernel associated with $\varphi^\tau$, and define
\begin{equation}\label{eq:zeta-discrete-def}
  \mathsf S_k^\tau
  :=
  \frac{\tau(1-a_\tau)}{a_\tau}\psi_k^\tau,
  \qquad k\ge1.
\end{equation}
Then \eqref{eq:Htau-rec-local} is equivalent to
\begin{equation}\label{eq:Htau-resolvent-discrete}
\begin{aligned}
  H_n^\tau(z)
  &=
  \tau^{-\alpha}F_z(H_n^\tau(z))
  +\widehat E_n^\tau(z)
\\
  &\quad
  +
  \frac{a_\tau}{1-a_\tau}\tau^{-\alpha-1}
  \sum_{j=1}^{n-1}\mathsf S_{n-j}^\tau F_z(H_j^\tau(z))
  +
  \sum_{j=1}^{n-1}\psi_{n-j}^\tau\widehat E_j^\tau(z),
\end{aligned}
\end{equation}
for every $1\le n\le N_\tau$.
\end{lemma}

\begin{proof}
Equation \eqref{eq:Htau-rec-local} is a discrete renewal equation of the form
\[
  H^\tau = U^\tau + \varphi^\tau * H^\tau,
\]
where
\[
  U_n^\tau(z):=\tau^{-\alpha}F_z(H_n^\tau(z))+\widehat E_n^\tau(z).
\]
Since $\|\varphi^\tau\|_1=a_\tau<1$, the associated renewal kernel
$\psi^\tau:=\sum_{m\ge1}(\varphi^\tau)^{*m}$
is well defined in $\ell^1$. Iterating the renewal equation and summing the
Neumann series yields
\[
  H^\tau = U^\tau + \psi^\tau * U^\tau,
\]
that is,
\[
  H_n^\tau(z)
  =
  U_n^\tau(z)
  +\sum_{j=1}^{n-1}\psi_{n-j}^\tau U_j^\tau(z),
  \qquad n\ge1.
\]
Substituting the definition of $U_j^\tau(z)$ yields
\[
  H_n^\tau(z)
  =
  \tau^{-\alpha}F_z(H_n^\tau(z))
  +\widehat E_n^\tau(z)
  +\sum_{j=1}^{n-1}\psi_{n-j}^\tau
   \left(
     \tau^{-\alpha}F_z(H_j^\tau(z))+\widehat E_j^\tau(z)
   \right).
\]
Using \eqref{eq:zeta-discrete-def},
\[
  \psi_{n-j}^\tau
  =
  \frac{a_\tau}{1-a_\tau}\frac1\tau \mathsf{S}_{n-j}^\tau,
\]
we obtain \eqref{eq:Htau-resolvent-discrete}.
\end{proof}

%%%%%%%%%%%%%%%%%%%%%%%%%%%%%%%%%%%%%%%%%%%%%%%%%%

% Appendix B.3 (previously a separate subsection; now part of the merged
% "Equivalent resolvent reformulation" subsection above)

\begin{lemma}[Continuous resolvent form]
\label{lem:continuous-resolvent}
For each $z$ in the admissible strip under consideration, the solution of the
continuous fractional Riccati equation admits the representation
\begin{equation}\label{eq:continuous-resolvent}
  h(t,z)
  =
  \frac1\gamma\int_0^t f_{\alpha,\gamma}(t-s)F_z(h(s,z))\,ds,
  \qquad 0\le t\le T.
\end{equation}
\end{lemma}

\begin{proof}
Starting from \eqref{eq:riccati-volterra-weak}, we rewrite it as
\[
  h(t,z)
  =
  \frac1{\Gamma(\alpha)}\int_0^t (t-s)^{\alpha-1}
  \left(-\gamma h(s,z)+F_z(h(s,z))\right)\,ds.
\]
Taking Laplace transforms gives
\[
  \widehat h(\lambda,z)
  =
  \lambda^{-\alpha}
  \left(
    -\gamma \widehat h(\lambda,z)
    +\mathcal L_t[F_z(h(\cdot,z))](\lambda)
  \right),
\]
and hence,
\[
  \widehat h(\lambda,z)
  =
  \frac{1}{\lambda^\alpha+\gamma}
  \mathcal L_t[F_z(h(\cdot,z))](\lambda).
\]
Since
\[
  \mathcal L\!\left\{\frac1\gamma f_{\alpha,\gamma}\right\}(\lambda)
  =
  \frac{1}{\lambda^\alpha+\gamma},
\]
the claim follows by Laplace inversion.
\end{proof}

%%%%%%%%%%%%%%%%%%%%%%%%%%%%%%%%%%%%%%%%%%%%%%%%%%

\subsection{Lemma~\ref{lem:cum-resolvent-renewal}}\label{app:proof-cum-resolvent-renewal}

\begin{lemma}[Discrete cumulative renewal equation]
\label{lem:cum-resolvent-renewal}
For all $n\ge1$,
\begin{equation}\label{eq:cum-resolvent-renewal}
  \mathcal G_n^\tau
  =
  a_\tau\sum_{k=1}^n\varphi_k\,\mathcal G_{n-k}^\tau
  +\tau^{-\alpha}\sum_{k=1}^n\varphi_k.
\end{equation}
Moreover,
\begin{equation}\label{eq:cum-resolvent-uniform}
  \sup_{n\ge0}\mathcal G_n^\tau\le\frac1\gamma.
\end{equation}
\end{lemma}

\begin{proof}
Let
\[
  \Psi_\tau(q):=\sum_{n\ge1}\psi_n^\tau q^n
  =\frac{a_\tau\Phi(q)}{1-a_\tau\Phi(q)},
  \qquad
  \Phi(q):=\sum_{m\ge1}\varphi_m q^m.
\]
Then
\[
  \sum_{n\ge1}\left(\sum_{k=1}^n\psi_k^\tau\right)q^n
  =\frac{\Psi_\tau(q)}{1-q}
  =\frac{a_\tau\Phi(q)}{(1-q)(1-a_\tau\Phi(q))}.
\]
Multiplying by $(1-a_\tau)/(\gamma a_\tau)$, we obtain
\[
  \sum_{n\ge1}\mathcal G_n^\tau q^n
  =
  \frac{1-a_\tau}{\gamma}\,
  \frac{\Phi(q)}{(1-q)(1-a_\tau\Phi(q))}.
\]
Hence,
\[
  (1-a_\tau\Phi(q))\sum_{n\ge1}\mathcal G_n^\tau q^n
  =
  \frac{1-a_\tau}{\gamma}\frac{\Phi(q)}{1-q}.
\]
Comparing coefficients of $q^n$ gives
\[
  \mathcal G_n^\tau
  =
  a_\tau\sum_{k=1}^n\varphi_k\,\mathcal G_{n-k}^\tau
  +\frac{1-a_\tau}{\gamma}\sum_{k=1}^n\varphi_k.
\]
Since $(1-a_\tau)/\gamma=\tau^{-\alpha}$, this proves
\eqref{eq:cum-resolvent-renewal}.

Because $\varphi_k\ge0$, $a_\tau\in(0,1)$, and
$0\le \sum_{k=1}^n\varphi_k\le 1$,
we obtain
\[
  \mathcal G_n^\tau
  \le a_\tau \max_{0\le j\le n-1}\mathcal G_j^\tau+\tau^{-\alpha}.
\]
Setting $M_n^\tau:=\max_{0\le j\le n}\mathcal G_j^\tau$, we get
\[
  M_n^\tau\le a_\tau M_{n-1}^\tau+\tau^{-\alpha}.
\]
Since $M_0^\tau=0$, iteration yields
\[
  M_n^\tau
  \le \tau^{-\alpha}\sum_{r=0}^{n-1}a_\tau^r
  \le \frac{\tau^{-\alpha}}{1-a_\tau}
  =\frac1\gamma.
\]
This proves \eqref{eq:cum-resolvent-uniform}.
\end{proof}

%%%%%%%%%%%%%%%%%%%%%%%%%%%%%%%%%%%%%%%%%%%%%%%%%%

\subsection{Lemma~\ref{lem:cum-resolvent-laplace}}\label{app:proof-cum-resolvent-laplace}

\begin{lemma}[Laplace representation]
\label{lem:cum-resolvent-laplace}
Define the piecewise-constant interpolation
\[
  \mathcal G^\tau(t):=\mathcal G_n^\tau,
  \qquad t\in[n/\tau,(n+1)/\tau).
\]
Then, for any $z>0$,
\begin{equation}\label{eq:cum-resolvent-laplace-simplified}
  \widehat{\mathcal G^\tau}(z)
  =
  \frac{1-a_\tau}{\gamma z}\,
  \frac{\hat\varphi(z/\tau)}{1-a_\tau\hat\varphi(z/\tau)}.
\end{equation}
\end{lemma}

\begin{proof}
Let
\[
  q:=e^{-z/\tau}\in(0,1).
\]
Since $\mathcal G^\tau(t)=\mathcal G_n^\tau$ on $[n/\tau,(n+1)/\tau)$, we have
\begin{align*}
  \widehat{\mathcal G^\tau}(z)
  =
  \sum_{n\ge0}\mathcal G_n^\tau
  \int_{n/\tau}^{(n+1)/\tau}e^{-zt}\,dt
  =
  \frac{1-q}{z}\sum_{n\ge0}\mathcal G_n^\tau q^n.
\end{align*}
Since $\mathcal G_0^\tau=0$, Lemma~\ref{lem:cum-resolvent-renewal} yields
\[
  \sum_{n\ge0}\mathcal G_n^\tau q^n
  =
  \sum_{n\ge1}\mathcal G_n^\tau q^n
  =
  \frac{1-a_\tau}{\gamma}\,
  \frac{\Phi(q)}{(1-q)(1-a_\tau\Phi(q))},
\]
where $\Phi(q):=\sum_{m\ge1}\varphi_m q^m$. Therefore,
\[
  \widehat{\mathcal G^\tau}(z)
  =
  \frac{1-a_\tau}{\gamma z}\,
  \frac{\Phi(q)}{1-a_\tau\Phi(q)}.
\]
Since $\Phi(e^{-z/\tau})=\hat\varphi(z/\tau)$, this is exactly
\eqref{eq:cum-resolvent-laplace-simplified}.
\end{proof}

%%%%%%%%%%%%%%%%%%%%%%%%%%%%%%%%%%%%%%%%%%%%%%%%%%

\subsection{Lemma~\ref{lem:cum-resolvent-laplace-domain-comparison}}\label{app:proof-cum-resolvent-laplace-domain-comparison}

\begin{lemma}[Laplace-domain consistency]
\label{lem:cum-resolvent-laplace-domain-comparison}
Let $K\subset(0,\infty)$ be compact. Then uniformly for $z\in K$,
\[
  \widehat{\mathcal G^\tau}(z)-\widehat{\mathcal G}(z)
  =
  \mathcal O(\tau^{-\alpha})
  +
  \mathcal O\!\left(|c_1(\alpha)|\,\tau^{-(1-\alpha)}\right),
\]
where
\[
  \widehat{\mathcal G}(z)=\frac{1}{z(\gamma+z^\alpha)}.
\]
\end{lemma}

\begin{proof}
Set $s:=z/\tau$. By Lemma~\ref{lem:phi-expansion},
\[
  \hat\varphi(s)
  =
  1-s^\alpha-c_1(\alpha)s+\mathcal O(s^{1+\alpha}),
\]
uniformly for $z\in K$ as $\tau\to\infty$, because $K$ is compact. Therefore
\begin{align*}
  1-a_\tau\hat\varphi(s)
  &=
  (1-a_\tau)+a_\tau(1-\hat\varphi(s))
  \\
  &=
  \gamma\tau^{-\alpha}
  +a_\tau z^\alpha\tau^{-\alpha}
  +a_\tau c_1(\alpha)z\tau^{-1}
  +\mathcal O(\tau^{-1-\alpha})
  \\
  &=
  \tau^{-\alpha}
  \left(
    \gamma+z^\alpha+c_1(\alpha)z\tau^{-(1-\alpha)}
    +\mathcal O(\tau^{-\alpha})
  \right),
\end{align*}
again uniformly for $z\in K$. Since $K\subset(0,\infty)$, the quantity
$\gamma+z^\alpha$ is bounded away from $0$ on $K$. Therefore, for $\tau$ large
enough,
\[
  \left|1-a_\tau\hat\varphi(z/\tau)\right|
  \ge c_K\tau^{-\alpha}
\]
for some $c_K>0$.
Using Lemma~\ref{lem:cum-resolvent-laplace},
\[
  \widehat{\mathcal G^\tau}(z)
  =
  \frac{1-a_\tau}{\gamma z}\,
  \frac{\hat\varphi(s)}{1-a_\tau\hat\varphi(s)}
  =
  \frac{\hat\varphi(s)}
       {z\left(
         \gamma+z^\alpha+c_1(\alpha)z\tau^{-(1-\alpha)}
         +\mathcal O(\tau^{-\alpha})
       \right)}.
\]
Since
\[
  \hat\varphi(s)=1+\mathcal O(\tau^{-\alpha})
\]
uniformly on $K$, and since
\[
  \frac{1}{A+B}
  =
  \frac{1}{A}+\mathcal O\!\left(\frac{B}{A^2}\right)
\]
for $A=\gamma+z^\alpha$ and
$B=c_1(\alpha)z\tau^{-(1-\alpha)}+\mathcal O(\tau^{-\alpha})$,
a first-order expansion of the reciprocal around $\gamma+z^\alpha$ yields
\[
  \widehat{\mathcal G^\tau}(z)
  =
  \frac{1}{z(\gamma+z^\alpha)}
  +
  \mathcal O(\tau^{-\alpha})
  +
  \mathcal O\!\left(|c_1(\alpha)|\,\tau^{-(1-\alpha)}\right),
\]
uniformly for $z\in K$.
\end{proof}

%%%%%%%%%%%%%%%%%%%%%%%%%%%%%%%%%%%%%%%%%%%%%%%%%%

\subsection{Proposition~\ref{prop:riccati-regularity}}\label{app:riccati-regularity}

\begin{proposition}[Bounded-strip regularity of the continuous fractional Riccati solution]
\label{prop:riccati-regularity}
Fix $\eta\in\R$ and $R>0$ such that
\[
  \Gamma_{\eta,R}:=\{\eta+\ii\xi:\ |\xi|\le R\}\subset \mathcal S_T.
\]
Under Assumption~\ref{ass:riccati-bound}, the map $(t,z)\mapsto h(t,z)$ is
jointly continuous on $[0,T]\times\Gamma_{\eta,R}$, each map
$t\mapsto h(t,z)$ belongs to $AC([0,T])\cap C^1((0,T])$, and there exists a
constant $C_{\eta,R,T}>0$ such that, uniformly for $z\in\Gamma_{\eta,R}$,
\begin{align}
  |h(t,z)| &\le C_{\eta,R,T}t^\alpha,
  &&0\le t\le T,\label{eq:h-small-t-prop}\\
  |h(t,z)-h(s,z)| &\le C_{\eta,R,T}|t-s|^\alpha,
  &&0\le s,t\le T,\label{eq:h-holder-t-prop}\\
  |\partial_t h(t,z)| &\le C_{\eta,R,T}t^{\alpha-1},
  &&0<t\le T.\label{eq:h-derivative-t-prop}
\end{align}
Moreover, the map $z\mapsto h(\cdot,z)$ is locally Lipschitz from
$\Gamma_{\eta,R}$ into $C([0,T])$.
\end{proposition}

\begin{proof}
Throughout the proof, $C_{\eta,R,T}>0$ denotes a generic constant that may
change from line to line but depends only on $(\eta,R,T)$ and the fixed model
parameters.

Let
\[
  F_z(x):=\frac12(z^2-z)+\gamma(\rho\nu z-1)x+\frac{(\gamma\nu)^2}{2}x^2,
  \qquad x\in\C.
\]
Then the Volterra form of the fractional Riccati equation reads
\begin{equation}\label{eq:riccati-volterra-proof-reg}
  h(t,z)
  =
  \frac{1}{\Gamma(\alpha)}
  \int_0^t (t-s)^{\alpha-1}F_z(h(s,z))\,ds,
  \qquad 0\le t\le T.
\end{equation}

Since $\Gamma_{\eta,R}$ is compact and
\[
  \sup_{z\in\Gamma_{\eta,R}}\sup_{0\le t\le T}|h(t,z)|
  \le M_{\eta,R,T},
\]
the polynomial $F_z(x)$ and its derivative in $x$ are uniformly bounded on the
set
\[
  \{(z,x): z\in\Gamma_{\eta,R},\ |x|\le M_{\eta,R,T}\}.
\]
Hence there exists a constant $C_{\eta,R,T}>0$ such that
\begin{equation}\label{eq:F-uniform-bound-proof-reg}
  |F_z(h(t,z))|\le C_{\eta,R,T},
  \qquad
  0\le t\le T,\ z\in\Gamma_{\eta,R},
\end{equation}
and
\begin{equation}\label{eq:F-lipschitz-proof-reg}
  |F_z(x)-F_z(y)|\le C_{\eta,R,T}|x-y|,
\end{equation}
for all $z\in\Gamma_{\eta,R}$ and all $x,y$ with $|x|,|y|\le M_{\eta,R,T}$.

We divide the proof into four steps.

\medskip
\noindent
\textit{Step 1: small-time bound.}
By \eqref{eq:riccati-volterra-proof-reg} and \eqref{eq:F-uniform-bound-proof-reg},
\begin{align*}
  |h(t,z)|
  &\le
  \frac{1}{\Gamma(\alpha)}
  \int_0^t (t-s)^{\alpha-1}|F_z(h(s,z))|\,ds\\
  &\le
  \frac{C_{\eta,R,T}}{\Gamma(\alpha)}
  \int_0^t (t-s)^{\alpha-1}\,ds
  =
  \frac{C_{\eta,R,T}}{\Gamma(\alpha+1)}\,t^\alpha.
\end{align*}
This proves \eqref{eq:h-small-t-prop}.

\medskip
\noindent
\textit{Step 2: $\alpha$-H\"older continuity in time.}
Fix $0\le s<t\le T$. Subtracting the two Volterra formulas gives
\begin{align*}
  h(t,z)-h(s,z)
  &=
  \frac{1}{\Gamma(\alpha)}
  \int_0^s
  \left((t-u)^{\alpha-1}-(s-u)^{\alpha-1}\right)F_z(h(u,z))\,du\\
  &\quad
  +
  \frac{1}{\Gamma(\alpha)}
  \int_s^t (t-u)^{\alpha-1}F_z(h(u,z))\,du.
\end{align*}
Taking absolute values and using \eqref{eq:F-uniform-bound-proof-reg},
\begin{align*}
  |h(t,z)-h(s,z)|
  &\le
  \frac{C_{\eta,R,T}}{\Gamma(\alpha)}
  \int_0^s
  \left|(t-u)^{\alpha-1}-(s-u)^{\alpha-1}\right|du
  \\
  &\qquad\qquad\qquad
  +
  \frac{C_{\eta,R,T}}{\Gamma(\alpha)}
  \int_s^t (t-u)^{\alpha-1}du.
\end{align*}
Since $\alpha-1\in(-1,0)$, the map $x\mapsto x^{\alpha-1}$ is decreasing on
$(0,\infty)$, and therefore
\begin{align*}
  \int_0^s\left((s-u)^{\alpha-1}-(t-u)^{\alpha-1}\right)\,du
  &=
  \frac{1}{\alpha}\left(s^\alpha-t^\alpha+(t-s)^\alpha\right)
  \le \frac{1}{\alpha}(t-s)^\alpha.
\end{align*}
Also,
$\int_s^t (t-u)^{\alpha-1}\,du=\frac1\alpha (t-s)^\alpha$.
Hence,
\[
  |h(t,z)-h(s,z)|
  \le C_{\eta,R,T}|t-s|^\alpha,
\]
uniformly in $z\in\Gamma_{\eta,R}$. This proves \eqref{eq:h-holder-t-prop}.

\medskip
\noindent
\textit{Step 3: continuity and Lipschitz dependence on $z$.}
Fix $z_1,z_2\in\Gamma_{\eta,R}$ and write
\[
  \delta h(t):=h(t,z_1)-h(t,z_2).
\]
Subtracting the two Volterra equations gives
\begin{align*}
  \delta h(t)
  &=
  \frac{1}{\Gamma(\alpha)}
  \int_0^t (t-s)^{\alpha-1}
  \left(F_{z_1}(h(s,z_1))-F_{z_2}(h(s,z_2))\right)\,ds.
\end{align*}
Split the integrand as
\begin{align*}
  F_{z_1}(h(s,z_1))-F_{z_2}(h(s,z_2))
  &=
  \left(F_{z_1}(h(s,z_1))-F_{z_1}(h(s,z_2))\right)
  \\
  &\quad+
  \left(F_{z_1}(h(s,z_2))-F_{z_2}(h(s,z_2))\right).
\end{align*}
By \eqref{eq:F-lipschitz-proof-reg},
\[
  \left|F_{z_1}(h(s,z_1))-F_{z_1}(h(s,z_2))\right|
  \le C_{\eta,R,T}|\delta h(s)|.
\]
Since $F_z(x)$ is polynomial in $(z,x)$ and $z\in\Gamma_{\eta,R}$, $|x|\le
M_{\eta,R,T}$, there exists $C_{\eta,R,T}>0$ such that
\[
  \left|F_{z_1}(h(s,z_2))-F_{z_2}(h(s,z_2))\right|
  \le C_{\eta,R,T}|z_1-z_2|.
\]
Therefore,
\[
  |\delta h(t)|
  \le
  C_{\eta,R,T}|z_1-z_2|t^\alpha
  +
  C_{\eta,R,T}
  \int_0^t (t-s)^{\alpha-1}|\delta h(s)|\,ds.
\]
Set
\[
u(t):=\sup_{0\le s\le t}|h(s,z_1)-h(s,z_2)|.
\]
Then the previous estimate implies
\[
u(t)\le C_{\eta,R,T}|z_1-z_2|
   +C_{\eta,R,T}\int_0^t (t-s)^{\alpha-1}u(s)\,ds,
   \qquad 0\le t\le T.
\]
By the weakly singular Gr\"onwall inequality, it follows that
\[
u(t)\le C_{\eta,R,T}|z_1-z_2|\,E_\alpha(C_{\eta,R,T}t^\alpha),
\qquad 0\le t\le T.
\]
Since $t\in[0,T]$, we obtain
\[
\sup_{0\le t\le T}|h(t,z_1)-h(t,z_2)|
\le C_{\eta,R,T}|z_1-z_2|.
\]
In particular, $z\mapsto h(\cdot,z)$ is locally Lipschitz into $C([0,T])$.
Combining this with \eqref{eq:h-holder-t-prop} gives joint continuity of
$(t,z)\mapsto h(t,z)$ on $[0,T]\times\Gamma_{\eta,R}$.

\medskip
\noindent
\textit{Step 4: absolute continuity and derivative bound.}
Define
\[
  g_z(t):=F_z(h(t,z)),\qquad 0\le t\le T.
\]
By \eqref{eq:h-holder-t-prop} and the local Lipschitz property
\eqref{eq:F-lipschitz-proof-reg},
\[
  |g_z(t)-g_z(s)|\le C_{\eta,R,T}|t-s|^\alpha,
  \qquad 0\le s,t\le T.
\]
Thus $g_z\in C^\alpha([0,T])$ uniformly in $z$, and
$h(\cdot,z)=I^\alpha g_z$.
We now prove that $h(\cdot,z)$ is differentiable on $(0,T]$.
Fix $z\in\Gamma_{\eta,R}$ and $t\in(0,T]$. Recall that
\[
  h(t,z)=\frac{1}{\Gamma(\alpha)}\int_0^t (t-s)^{\alpha-1}g_z(s)\,ds,
\]
where
\[
  |g_z(s)-g_z(r)|\le C_{\eta,R,T}|s-r|^\alpha,
  \qquad 0\le r,s\le T.
\]

Let $h_0$ be such that $|h_0|<\min\{t,T-t\}$, and let
$u\in(-h_0,h_0)\setminus\{0\}$.

\smallskip
\noindent
\emph{Case 1: $u>0$.}
In this case,
\begin{align*}
  \frac{h(t+u,z)-h(t,z)}{u}
  &=
  \frac{1}{\Gamma(\alpha)}
  \int_0^t
  \frac{(t+u-s)^{\alpha-1}-(t-s)^{\alpha-1}}{u}\,g_z(s)\,ds
  \\
  &\quad+
  \frac{1}{\Gamma(\alpha)}
  \frac{1}{u}\int_t^{t+u}(t+u-s)^{\alpha-1}g_z(s)\,ds.
\end{align*}
Adding and subtracting $g_z(t)$, we rewrite this as
\begin{align}
  &\frac{h(t+u,z)-h(t,z)}{u}\nonumber\\
  &=
  \frac{1}{\Gamma(\alpha)}
  \int_0^t
  \frac{(t+u-s)^{\alpha-1}-(t-s)^{\alpha-1}}{u}
  \left(g_z(s)-g_z(t)\right)\,ds
  \nonumber\\
  &\qquad+
  \frac{1}{\Gamma(\alpha)}
  \frac{1}{u}\int_t^{t+u}(t+u-s)^{\alpha-1}
  \left(g_z(s)-g_z(t)\right)\,ds
  \nonumber\\
  &\qquad\qquad+
  \frac{g_z(t)}{\Gamma(\alpha)}
  \left[
    \int_0^t
    \frac{(t+u-s)^{\alpha-1}-(t-s)^{\alpha-1}}{u}\,ds
    +
    \frac{1}{u}\int_t^{t+u}(t+u-s)^{\alpha-1}\,ds
  \right].
  \label{eq:dq-split-proof-pos}
\end{align}

For the bracketed term, a direct computation gives
\begin{align*}
  &\int_0^t
    \frac{(t+u-s)^{\alpha-1}-(t-s)^{\alpha-1}}{u}\,ds
    +
    \frac{1}{u}\int_t^{t+u}(t+u-s)^{\alpha-1}\,ds
  \\
  &=
  \frac{1}{u}\int_0^{t+u}r^{\alpha-1}\,dr
  -
  \frac{1}{u}\int_0^t r^{\alpha-1}\,dr
  =
  \frac{(t+u)^\alpha-t^\alpha}{\alpha u},
\end{align*}
and therefore
\begin{equation}\label{eq:bracket-limit-pos}
  \frac{g_z(t)}{\Gamma(\alpha)}
  \left[
    \int_0^t
    \frac{(t+u-s)^{\alpha-1}-(t-s)^{\alpha-1}}{u}\,ds
    +
    \frac{1}{u}\int_t^{t+u}(t+u-s)^{\alpha-1}\,ds
  \right]
  \to
  \frac{g_z(t)}{\Gamma(\alpha)}\,t^{\alpha-1}
\end{equation}
as $u\downarrow0$.

For the first integral in \eqref{eq:dq-split-proof-pos}, fix $0<u<t/2$. By the mean value theorem,
for each $s\in[0,t)$ there exists $\theta=\theta(s,u)\in(0,1)$ such that
\[
  \frac{(t+u-s)^{\alpha-1}-(t-s)^{\alpha-1}}{u}
  =
  (\alpha-1)(t-s+\theta u)^{\alpha-2}.
\]
Hence
\[
  \left|
    \frac{(t+u-s)^{\alpha-1}-(t-s)^{\alpha-1}}{u}
  \right|
  \le C_\alpha (t-s)^{\alpha-2},
\]
and thus, using the $\alpha$-H\"older continuity of $g_z$,
\[
  \left|
    \frac{(t+u-s)^{\alpha-1}-(t-s)^{\alpha-1}}{u}
    \left(g_z(s)-g_z(t)\right)
  \right|
  \le C_{\eta,R,T}(t-s)^{2\alpha-2}.
\]
Since $\alpha>1/2$, we have $2\alpha-2>-1$, so $(t-s)^{2\alpha-2}\in L^1(0,t)$.
Moreover, for each fixed $s\in[0,t)$,
\[
  \frac{(t+u-s)^{\alpha-1}-(t-s)^{\alpha-1}}{u}
  \to
  (\alpha-1)(t-s)^{\alpha-2}
\]
as $u\downarrow0$. Hence, by the dominated convergence theorem,
\begin{align}
  &\int_0^t
  \frac{(t+u-s)^{\alpha-1}-(t-s)^{\alpha-1}}{u}
  \left(g_z(s)-g_z(t)\right)\,ds
  \nonumber\\
  &\to
  (\alpha-1)\int_0^t (t-s)^{\alpha-2}\left(g_z(s)-g_z(t)\right)\,ds
  \qquad\text{as }u\downarrow0.
  \label{eq:first-integral-limit-pos}
\end{align}

For the second integral in \eqref{eq:dq-split-proof-pos}, again by the $\alpha$-H\"older continuity of $g_z$,
\begin{align}
  &\left|
    \frac{1}{u}\int_t^{t+u}(t+u-s)^{\alpha-1}
    \left(g_z(s)-g_z(t)\right)\,ds
  \right|
  \nonumber\\
  &\le
  \frac{C_{\eta,R,T}}{u}\int_t^{t+u}(t+u-s)^{\alpha-1}(s-t)^\alpha\,ds
  \nonumber\\
 &=
  \frac{C_{\eta,R,T}}{u}\int_0^u r^{\alpha-1}(u-r)^\alpha\,dr
 =
  C_{\eta,R,T}u^{2\alpha-1}\int_0^1 x^{\alpha-1}(1-x)^\alpha\,dx
  \to0
\label{eq:second-integral-limit-pos}
\end{align}
as $u\downarrow0$.

Combining the three limits above in \eqref{eq:bracket-limit-pos}, \eqref{eq:first-integral-limit-pos}, and \eqref{eq:second-integral-limit-pos},
we conclude that
\begin{equation}
  \lim_{u\downarrow0}\frac{h(t+u,z)-h(t,z)}{u}
  =
  \frac{g_z(t)}{\Gamma(\alpha)}t^{\alpha-1}
  +
  \frac{\alpha-1}{\Gamma(\alpha)}
  \int_0^t (t-s)^{\alpha-2}\left(g_z(s)-g_z(t)\right)\,ds .
  \label{eq:right-derivative-proof}
\end{equation}

\smallskip
\noindent
\emph{Case 2: $u<0$.}
In this case,
\begin{align*}
  \frac{h(t+u,z)-h(t,z)}{u}
  &=
  \frac{1}{\Gamma(\alpha)}
  \int_0^{t+u}
  \frac{(t+u-s)^{\alpha-1}-(t-s)^{\alpha-1}}{u}\,g_z(s)\,ds
  \\
  &\quad
  -
  \frac{1}{\Gamma(\alpha)}
  \frac{1}{u}\int_{t+u}^{t}(t-s)^{\alpha-1}g_z(s)\,ds.
\end{align*}
Adding and subtracting $g_z(t)$, we obtain
\begin{align}
  &\frac{h(t+u,z)-h(t,z)}{u}\nonumber\\
  &=
  \frac{1}{\Gamma(\alpha)}
  \int_0^{t+u}
  \frac{(t+u-s)^{\alpha-1}-(t-s)^{\alpha-1}}{u}
  \left(g_z(s)-g_z(t)\right)\,ds
  \nonumber\\
  &\qquad
  -
  \frac{1}{\Gamma(\alpha)}
  \frac{1}{u}\int_{t+u}^{t}(t-s)^{\alpha-1}
  \left(g_z(s)-g_z(t)\right)\,ds
  \nonumber\\
  &\qquad\qquad+
  \frac{g_z(t)}{\Gamma(\alpha)}
  \left[
    \int_0^{t+u}
    \frac{(t+u-s)^{\alpha-1}-(t-s)^{\alpha-1}}{u}\,ds
    -
    \frac{1}{u}\int_{t+u}^{t}(t-s)^{\alpha-1}\,ds
  \right].
  \label{eq:dq-split-proof-neg}
\end{align}

For the bracketed term, another direct computation gives
\begin{align*}
  &\int_0^{t+u}
  \frac{(t+u-s)^{\alpha-1}-(t-s)^{\alpha-1}}{u}\,ds
  -
  \frac{1}{u}\int_{t+u}^{t}(t-s)^{\alpha-1}\,ds
  \\
  &=
  \frac{1}{u}\int_0^{t+u}r^{\alpha-1}\,dr
  -
  \frac{1}{u}\int_0^t r^{\alpha-1}\,dr
  =
  \frac{(t+u)^\alpha-t^\alpha}{\alpha u},
\end{align*}
so again
\begin{equation}\label{eq:bracket-limit-neg}
  \frac{g_z(t)}{\Gamma(\alpha)}
  \left[
    \int_0^{t+u}
    \frac{(t+u-s)^{\alpha-1}-(t-s)^{\alpha-1}}{u}\,ds
    -
    \frac{1}{u}\int_{t+u}^{t}(t-s)^{\alpha-1}\,ds
  \right]
  \to
  \frac{g_z(t)}{\Gamma(\alpha)}\,t^{\alpha-1}
\end{equation}
as $u\uparrow0$.

For the first integral in \eqref{eq:dq-split-proof-neg}, the same dominated convergence argument as above applies, because for fixed $s\in[0,t)$,
\[
  \frac{(t+u-s)^{\alpha-1}-(t-s)^{\alpha-1}}{u}
  \to
  (\alpha-1)(t-s)^{\alpha-2},
\]
as $u\uparrow0$, and the same integrable domination by
$C_{\eta,R,T}(t-s)^{2\alpha-2}$ holds. Hence,
\begin{align}
  &\int_0^{t+u}
  \frac{(t+u-s)^{\alpha-1}-(t-s)^{\alpha-1}}{u}
  \left(g_z(s)-g_z(t)\right)\,ds
  \nonumber\\
  &\to
  (\alpha-1)\int_0^t (t-s)^{\alpha-2}\left(g_z(s)-g_z(t)\right)\,ds,
  \qquad\text{as }u\uparrow0.
  \label{eq:first-integral-limit-neg}
\end{align}

For the second integral in \eqref{eq:dq-split-proof-neg}, write $v:=-u>0$. Then
\begin{align}
  &\left|
    -\frac{1}{u}\int_{t+u}^{t}(t-s)^{\alpha-1}
    \left(g_z(s)-g_z(t)\right)\,ds
  \right|
  \nonumber\\
  &=
  \frac{1}{v}\left|
    \int_{t-v}^{t}(t-s)^{\alpha-1}
    \left(g_z(s)-g_z(t)\right)\,ds
  \right|
  \nonumber\\
  &\le
  \frac{C_{\eta,R,T}}{v}
  \int_{t-v}^{t}(t-s)^{\alpha-1}(t-s)^\alpha\,ds
 =
  C_{\eta,R,T}v^{2\alpha-1}
  \to0,
\label{eq:second-integral-limit-neg}
\end{align}
as $u\uparrow0$.

Combining the limits in \eqref{eq:bracket-limit-neg}, \eqref{eq:first-integral-limit-neg}, and \eqref{eq:second-integral-limit-neg},
we obtain
\begin{equation}
  \lim_{u\uparrow0}\frac{h(t+u,z)-h(t,z)}{u}
  =
  \frac{g_z(t)}{\Gamma(\alpha)}t^{\alpha-1}
  +
  \frac{\alpha-1}{\Gamma(\alpha)}
  \int_0^t (t-s)^{\alpha-2}\left(g_z(s)-g_z(t)\right)\,ds .
  \label{eq:left-derivative-proof}
\end{equation}

From \eqref{eq:right-derivative-proof} and \eqref{eq:left-derivative-proof}, the right and left derivatives agree. Therefore $h(\cdot,z)$ is differentiable on $(0,T]$, and
\begin{equation}\label{eq:h-derivative-formula-proof}
  \partial_t h(t,z)
  =
  \frac{g_z(t)}{\Gamma(\alpha)}t^{\alpha-1}
  +
  \frac{\alpha-1}{\Gamma(\alpha)}
  \int_0^t (t-s)^{\alpha-2}\left(g_z(s)-g_z(t)\right)\,ds .
\end{equation}

It remains to estimate the derivative. Since $|g_z(t)|\le C_{\eta,R,T}$,
\[
  \left|\frac{g_z(t)}{\Gamma(\alpha)}t^{\alpha-1}\right|
  \le C_{\eta,R,T}t^{\alpha-1}.
\]
For the integral term, using once more the $\alpha$-H\"older continuity of $g_z$,
\begin{align*}
  \left|
    \int_0^t (t-s)^{\alpha-2}\left(g_z(s)-g_z(t)\right)\,ds
  \right|
  \le
  C_{\eta,R,T}\int_0^t (t-s)^{2\alpha-2}\,ds
 =
  C_{\eta,R,T}\frac{t^{2\alpha-1}}{2\alpha-1}.
\end{align*}
Since $t\le T$, we have $t^{2\alpha-1}=t^{\alpha-1}t^\alpha\le T^\alpha t^{\alpha-1}$.
Therefore
\[
  |\partial_t h(t,z)|
  \le C_{\eta,R,T}t^{\alpha-1},
  \qquad 0<t\le T.
\]

Finally, since \eqref{eq:h-derivative-formula-proof} holds for every $t\in(0,T]$
and $|\partial_t h(t,z)|\le C_{\eta,R,T}\,t^{\alpha-1}$ with $t^{\alpha-1}\in L^1(0,T)$ (as $\alpha>0$), we have $\partial_t h(\cdot,z)\in L^1(0,T)$. Hence the representation
\[
  h(t,z)=h(0,z)+\int_0^t \partial_s h(s,z)\,ds,\qquad 0\le t\le T,
\]
which follows by integrating \eqref{eq:h-derivative-formula-proof} from $0$ to $t$ and using the dominated convergence theorem as the lower limit approaches~$0$, implies that $t\mapsto h(t,z)$ is absolutely continuous on $[0,T]$.
\end{proof}

%%%%%%%%%%%%%%%%%%%%%%%%%%%%%%%%%%%%%%%%%%%%



\subsection{Lemma~\ref{lem:h-time-regularity}}\label{app:proof-h-time-regularity}

\begin{lemma}[Regularity of the Riccati nonlinearity]
\label{lem:h-time-regularity}
Assume the hypotheses of Proposition~\ref{prop:riccati-regularity}. Fix $\eta\in\R$ and $R>0$,
and define
\[
  g_z(t)
  :=
  F_z(h(t,z))
  =
  \frac12(z^2-z)+\gamma\rho\nu z\,h(t,z)+\frac{(\gamma\nu)^2}{2}(h(t,z))^2.
\]
Then there exists a constant $C_{\eta,R,T}>0$ such that, uniformly for
$z\in\Gamma_{\eta,R}$,
\[
  \sup_{0\le t\le T}|g_z(t)|\le C_{\eta,R,T},
\]
$g_z$ is absolutely continuous on $[0,T]$, and
\begin{align}
  |g_z(t)-g_z(0)| &\le C_{\eta,R,T}t^\alpha,
  &&0\le t\le T,\label{eq:g-small-t-new}\\
  |g_z(t)-g_z(s)| &\le C_{\eta,R,T}|t-s|^\alpha,
  &&0\le s,t\le T,\label{eq:g-holder-t-new}\\
  |\partial_t g_z(t)| &\le C_{\eta,R,T}t^{\alpha-1},
  &&0<t\le T.\label{eq:g-derivative-t-new}
\end{align}
\end{lemma}

\begin{proof}
The bounds on $h$ stated in Proposition~\ref{prop:riccati-regularity} imply that
\[
  \sup_{0\le t\le T}\sup_{z\in\Gamma_{\eta,R}}|h(t,z)|
  \le C_{\eta,R,T},
\]
and therefore the polynomial $F_z(x)$ and its derivative in $x$ are uniformly
bounded on the set
\[
  \left\{(z,x):\ z\in\Gamma_{\eta,R},\ |x|\le C_{\eta,R,T}\right\}.
\]
Hence
\[
  \sup_{0\le t\le T}\sup_{z\in\Gamma_{\eta,R}}|g_z(t)|
  \le C_{\eta,R,T}.
\]

The estimate \eqref{eq:g-small-t-new} follows from
\[
  g_z(t)-g_z(0)
  =
  \gamma\rho\nu z\,h(t,z)+\frac{(\gamma\nu)^2}{2}\left(h(t,z)\right)^2
\]
and Proposition~\ref{prop:riccati-regularity}, which gives
$|h(t,z)|\le C_{\eta,R,T}t^\alpha$ uniformly on $\Gamma_{\eta,R}$.

Likewise, a direct application of the mean value theorem to the quadratic
polynomial $F_z$ together with the H\"older estimate for $h$ from
Proposition~\ref{prop:riccati-regularity} yields \eqref{eq:g-holder-t-new}.

Because $h(\cdot,z)$ is absolutely continuous on $[0,T]$ by
Proposition~\ref{prop:riccati-regularity}, the chain rule gives
\[
  \partial_t g_z(t)
  =
  \left(\gamma\rho\nu z+(\gamma\nu)^2 h(t,z)\right)\partial_t h(t,z),
  \qquad 0<t\le T.
\]
Using again the uniform bound on $h$ and the derivative estimate
\eqref{eq:h-derivative-t-prop} from
Proposition~\ref{prop:riccati-regularity}, we obtain
\eqref{eq:g-derivative-t-new}. Since
$|\partial_t g_z(t)|\le C_{\eta,R,T}t^{\alpha-1}$ and
$t^{\alpha-1}\in L^1(0,T)$, each $g_z$ is absolutely continuous on $[0,T]$.
\end{proof}

%%%%%%%%%%%%%%%%%%%%%%%%%%%%%%%%%%%%%%%%%%%%%%%%%%

\subsection{Lemma~\ref{lem:quadrature-c1}}\label{app:proof-quadrature-c1}

\begin{lemma}[Weakly singular quadrature constant]
\label{lem:quadrature-c1}
For $\alpha\in(0,1)$,
\begin{equation}\label{eq:quadrature-c1}
\begin{aligned}
  &\int_0^1\left(1-\frac{v^{-\alpha}}{\Gamma(1-\alpha)}\right)\,dv
  +
  \frac1{\Gamma(1-\alpha)}
  \sum_{j=1}^\infty\int_j^{j+1}(j^{-\alpha}-v^{-\alpha})\,dv\\
  &=
  1+\frac{\zeta(\alpha)}{\Gamma(1-\alpha)}
  =
  c_1(\alpha).
\end{aligned}
\end{equation}
Moreover,
\begin{equation}\label{eq:quadrature-c1-tail}
  \frac1{\Gamma(1-\alpha)}
  \sum_{j=n}^\infty\int_j^{j+1}(j^{-\alpha}-v^{-\alpha})\,dv
  \le C_\alpha n^{-\alpha},
  \qquad n\ge1.
\end{equation}
\end{lemma}

\begin{proof}
For $j\ge1$,
\[
  \int_j^{j+1}(j^{-\alpha}-v^{-\alpha})\,dv
  =
  j^{-\alpha}
  -
  \frac{(j+1)^{1-\alpha}-j^{1-\alpha}}{1-\alpha}.
\]
Therefore,
\begin{align*}
  \sum_{j=1}^{m}\int_j^{j+1}(j^{-\alpha}-v^{-\alpha})\,dv
  =
  \sum_{j=1}^{m}j^{-\alpha}
  -
  \frac{(m+1)^{1-\alpha}-1}{1-\alpha}.
\end{align*}
Hence,
\begin{align*}
  &\int_0^1\left(1-\frac{v^{-\alpha}}{\Gamma(1-\alpha)}\right)\,dv
  +
  \frac1{\Gamma(1-\alpha)}
  \sum_{j=1}^{m}\int_j^{j+1}(j^{-\alpha}-v^{-\alpha})\,dv
  \\
  &=
  1-\frac{1}{\Gamma(1-\alpha)(1-\alpha)}
  +
  \frac1{\Gamma(1-\alpha)}
  \left(
    \sum_{j=1}^{m}j^{-\alpha}
    -
    \frac{(m+1)^{1-\alpha}-1}{1-\alpha}
  \right)
  \\
  &=
  1+\frac1{\Gamma(1-\alpha)}
  \left(
    \sum_{j=1}^{m}j^{-\alpha}
    -
    \frac{(m+1)^{1-\alpha}}{1-\alpha}
  \right).
\end{align*}
Since
\[
  \sum_{j=1}^{m}j^{-\alpha}
  -\frac{(m+1)^{1-\alpha}}{1-\alpha}
  \to \zeta(\alpha)
  \qquad (m\to\infty),
\]
we obtain \eqref{eq:quadrature-c1}. Moreover, for every $n\ge1$,
\begin{align*}
  0
  \le
  \int_j^{j+1}(j^{-\alpha}-v^{-\alpha})\,dv
  &=
  \int_j^{j+1}\left(j^{-\alpha}-v^{-\alpha}\right)\,dv
  \\
  &\le
  C_\alpha\int_j^{j+1}j^{-\alpha-1}(v-j)\,dv
\le
  C_\alpha j^{-\alpha-1},
\end{align*}
where we used the mean value theorem and the monotonicity of
$x\mapsto x^{-\alpha}$. Summing over $j\ge n$ yields
\[
  \sum_{j=n}^{\infty}\int_j^{j+1}(j^{-\alpha}-v^{-\alpha})\,dv
  \le
  C_\alpha\sum_{j=n}^{\infty}j^{-\alpha-1}
  \le
  C_\alpha n^{-\alpha},
\]
which proves \eqref{eq:quadrature-c1-tail}.
\end{proof}

%%%%%%%%%%%%%%%%%%%%%%%%%%%%%%%%%%%%%%%%%%%%%%%%%%

\subsection{Lemma~\ref{lem:quadrature-defect-Gprime}}\label{app:proof-quadrature-defect-Gprime}

\begin{lemma}[Weakly singular quadrature defect for $\mathcal G'$]
\label{lem:quadrature-defect-Gprime}
Let
\[
  \mathcal G(t)=\frac1\gamma\left(1-E_\alpha(-\gamma t^\alpha)\right),
  \qquad t\ge0,
\]
and, for $1\le n\le N_\tau$, define
\begin{equation}\label{eq:Dn-def}
\begin{aligned}
  \mathcal D_n^\tau
  &:=
  \int_0^1\left(1-\frac{v^{-\alpha}}{\Gamma(1-\alpha)}\right)
  \frac1\tau \mathcal G'(t_n-v/\tau)\,dv
\\
  &\qquad
  +
  \frac1{\Gamma(1-\alpha)}
  \sum_{j=1}^{n-1}\int_j^{j+1}(j^{-\alpha}-v^{-\alpha})
  \frac1\tau \mathcal G'(t_n-v/\tau)\,dv .
\end{aligned}
\end{equation}
Then
\begin{equation}\label{eq:quadrature-defect-main}
  \mathcal D_n^\tau
  =
  c_1(\alpha)\frac1\tau \mathcal G'(t_n)+R_n^\tau,
\end{equation}
where
\begin{equation}\label{eq:Rn-bound}
  |R_n^\tau|\le C_T\tau^{-\alpha}n^{-\alpha},
  \qquad 1\le n\le N_\tau.
\end{equation}
\end{lemma}

\begin{proof}
Write
\[
  w_n(v):=\frac1\tau \mathcal G'(t_n-v/\tau),
  \qquad 0\le v\le n.
\]
Then \eqref{eq:Dn-def} becomes
\begin{align*}
  \mathcal D_n^\tau
  =
  \int_0^1\left(1-\frac{v^{-\alpha}}{\Gamma(1-\alpha)}\right)w_n(v)\,dv
  +
  \frac1{\Gamma(1-\alpha)}
  \sum_{j=1}^{n-1}\int_j^{j+1}(j^{-\alpha}-v^{-\alpha})w_n(v)\,dv.
\end{align*}
We split
\[
  w_n(v)=w_n(0)+\left(w_n(v)-w_n(0)\right).
\]

The contribution of the constant part $w_n(0)$ equals
\begin{align*}
  w_n(0)
  &\left[
    \int_0^1\left(1-\frac{v^{-\alpha}}{\Gamma(1-\alpha)}\right)\,dv
    +
    \frac1{\Gamma(1-\alpha)}
    \sum_{j=1}^{n-1}\int_j^{j+1}(j^{-\alpha}-v^{-\alpha})\,dv
  \right].
\end{align*}
By Lemma~\ref{lem:quadrature-c1},
\[
  \int_0^1\left(1-\frac{v^{-\alpha}}{\Gamma(1-\alpha)}\right)\,dv
  +
  \frac1{\Gamma(1-\alpha)}
  \sum_{j=1}^{n-1}\int_j^{j+1}(j^{-\alpha}-v^{-\alpha})\,dv
  =
  c_1(\alpha)+\mathcal O(n^{-\alpha}),
\]
and hence the constant part contributes
\[
  c_1(\alpha)\,w_n(0)+\mathcal O\!\left(n^{-\alpha}|w_n(0)|\right).
\]

Since
\[
  \mathcal G'(t)=\frac1\gamma f_{\alpha,\gamma}(t),
\]
and
\[
  |f_{\alpha,\gamma}(t)|\le C_T t^{\alpha-1},
  \qquad 0<t\le T,
\]
we obtain
\[
  |w_n(0)|
  =
  \frac1\tau |\mathcal G'(t_n)|
  \le
  C_T\tau^{-1}t_n^{\alpha-1}
  =
  C_T\tau^{-\alpha}n^{\alpha-1}.
\]
Therefore,
\[
  n^{-\alpha}|w_n(0)|
  \le
  C_T\tau^{-\alpha}n^{-1}
  \le
  C_T\tau^{-\alpha}n^{-\alpha},
\]
since $n\ge1$ and $\alpha\in(0,1)$.

It remains to estimate the oscillatory part
\begin{align*}
  \widetilde R_n^\tau
  &:=
  \int_0^1\left(1-\frac{v^{-\alpha}}{\Gamma(1-\alpha)}\right)
  \left(w_n(v)-w_n(0)\right)\,dv
  \\
  &\quad
  +
  \frac1{\Gamma(1-\alpha)}
  \sum_{j=1}^{n-1}\int_j^{j+1}(j^{-\alpha}-v^{-\alpha})
  \left(w_n(v)-w_n(0)\right)\,dv .
\end{align*}

We first record a derivative bound for $\mathcal G'$. Since
\[
  \mathcal G'(t)=\frac1\gamma f_{\alpha,\gamma}(t)
  =t^{\alpha-1}E_{\alpha,\alpha}(-\gamma t^\alpha),
\]
we may write
\[
  \mathcal G'(t)
  =
  \sum_{m=0}^\infty
  \frac{(-\gamma)^m}{\Gamma(\alpha m+\alpha)}\,t^{\alpha m+\alpha-1},
  \qquad t>0.
\]
For every $\delta>0$, this series and its termwise derivative converge uniformly
on $[\delta,T]$, so differentiation term by term is justified there and yields
\[
  (\mathcal G')'(t)
  =
  \sum_{m=0}^\infty
  \frac{(-\gamma)^m(\alpha m+\alpha-1)}{\Gamma(\alpha m+\alpha)}
  t^{\alpha m+\alpha-2}.
\]
The $m=0$ term is of order $t^{\alpha-2}$, while the sum over $m\ge1$ is bounded
by $C_T t^{2\alpha-2}$. Since $2\alpha-2>\alpha-2$ and $t\le T$, we obtain
\[
  |(\mathcal G')'(t)|
  \le C_T t^{\alpha-2},
  \qquad 0<t\le T.
\]

If $n=1$, the discrete sum is absent. Using the crude bound
\[
  |w_1(v)-w_1(0)|
  \le |w_1(v)|+|w_1(0)|
  \le C_T\tau^{-\alpha}\left((1-v)^{\alpha-1}+1\right),
  \qquad 0\le v\le1,
\]
we immediately get
\[
  |\widetilde R_1^\tau|
  \le C_T\tau^{-\alpha}
  = C_T\tau^{-\alpha}1^{-\alpha}.
\]
Hence it remains to consider $n\ge2$.

For the integral over $[0,1]$, the mean value theorem and the derivative bound above give
\[
  |w_n(v)-w_n(0)|
  \le
  C_T\tau^{-2}\,v\,(t_n-v/\tau)^{\alpha-2},
  \qquad 0\le v\le1.
\]
Since $t_n-v/\tau\ge (n-1)/\tau\asymp n/\tau$ for $n\ge2$, it follows that
\[
  |w_n(v)-w_n(0)|
  \le
  C_T\tau^{-\alpha}n^{\alpha-2}v.
\]
Because
\[
  1-\frac{v^{-\alpha}}{\Gamma(1-\alpha)}\in L^1(0,1),
\]
we obtain
\[
  \left|
    \int_0^1\left(1-\frac{v^{-\alpha}}{\Gamma(1-\alpha)}\right)
    \left(w_n(v)-w_n(0)\right)\,dv
  \right|
  \le
  C_T\tau^{-\alpha}n^{\alpha-2}
  \le
  C_T\tau^{-\alpha}n^{-\alpha},
\]
because $\alpha-2\le -\alpha$ for $\alpha\in(0,1)$.

For the discrete part, write
\[
  I_{j,n}
  :=
  \int_j^{j+1}(j^{-\alpha}-v^{-\alpha})\left(w_n(v)-w_n(0)\right)\,dv,
  \qquad 1\le j\le n-1.
\]
We split the sum into the regions $1\le j\le \lfloor n/2\rfloor$ and
$\lfloor n/2\rfloor+1\le j\le n-1$.

If $1\le j\le \lfloor n/2\rfloor$, then for $v\in[j,j+1]$ we have
$t_n-v/\tau\asymp n/\tau$ and $v\asymp j$. Hence
\[
  |w_n(v)-w_n(0)|
  \le
  C_T\tau^{-2}\,v\,(t_n-v/\tau)^{\alpha-2}
  \le
  C_T\tau^{-\alpha}j\,n^{\alpha-2}.
\]
Using
\[
  \int_j^{j+1}(j^{-\alpha}-v^{-\alpha})\,dv
  \le C_\alpha j^{-\alpha-1},
\]
we conclude that
\[
  |I_{j,n}|
  \le
  C_T\tau^{-\alpha}n^{\alpha-2}j^{-\alpha}.
\]
Therefore,
\[
  \sum_{j=1}^{\lfloor n/2\rfloor}|I_{j,n}|
  \le
  C_T\tau^{-\alpha}n^{\alpha-2}
  \sum_{j=1}^{\lfloor n/2\rfloor}j^{-\alpha}
  \le
  C_T\tau^{-\alpha}n^{-1}
  \le
  C_T\tau^{-\alpha}n^{-\alpha}.
\]

If $\lfloor n/2\rfloor+1\le j\le n-1$, then for $v\in[j,j+1]$ we use the crude bound
\[
  |w_n(v)-w_n(0)|
  \le
  |w_n(v)|+|w_n(0)|
  \le
  C_T\tau^{-\alpha}\left((n-v)^{\alpha-1}+n^{\alpha-1}\right).
\]
Moreover, $j\asymp n$ on this range, so
\[
  \int_j^{j+1}(j^{-\alpha}-v^{-\alpha})\,dv
  \le C_\alpha n^{-\alpha-1}.
\]
Hence
\begin{align*}
  \sum_{j=\lfloor n/2\rfloor+1}^{n-1}|I_{j,n}|
  &\le
  C_T\tau^{-\alpha}n^{-\alpha-1}
  \sum_{j=\lfloor n/2\rfloor+1}^{n-1}
  \int_j^{j+1}\left((n-v)^{\alpha-1}+n^{\alpha-1}\right)\,dv
  \\
  &\le
  C_T\tau^{-\alpha}n^{-\alpha-1}
  \left(
    \int_{n/2}^{n}(n-v)^{\alpha-1}\,dv
    +n^{\alpha-1}\sum_{j=\lfloor n/2\rfloor+1}^{n-1}1
  \right)
  \\
  &\le
  C_T\tau^{-\alpha}n^{-1}
  \le
  C_T\tau^{-\alpha}n^{-\alpha}.
\end{align*}

Combining the bounds for the integral part and the two discrete ranges, we obtain
\[
  |\widetilde R_n^\tau|
  \le
  C_T\tau^{-\alpha}n^{-\alpha},
  \qquad 1\le n\le N_\tau.
\]

Finally, combining the constant contribution and the oscillatory contribution,
we get
\[
  \mathcal D_n^\tau
  =
  c_1(\alpha)\,w_n(0)+R_n^\tau
  =
  c_1(\alpha)\frac1\tau \mathcal G'(t_n)+R_n^\tau,
\]
where
\[
  |R_n^\tau|
  \le
  C_T\tau^{-\alpha}n^{-\alpha}.
\]
This proves \eqref{eq:quadrature-defect-main}--\eqref{eq:Rn-bound}.
\end{proof}


%%%%%%%%%%%%%%%%%%%%%%%%%%%%%%%%%%%%%%%%%%%%%%%%%%

\subsection{Proof of Proposition~\ref{prop:cum-resolvent-uniform-comparison}}\label{app:proof-cum-resolvent-uniform-comparison}
We split the proof into five steps.

\medskip
\noindent
\textit{Step 1: error equation.}
By Lemma~\ref{lem:cum-resolvent-renewal}, the discrete cumulative resolvent
satisfies, for every $n\ge1$,
\begin{equation}\label{eq:Gtau-renewal-proof}
  \mathcal G_n^\tau
  =
  a_\tau\sum_{k=1}^n\varphi_k\,\mathcal G_{n-k}^\tau
  +\tau^{-\alpha}\Phi_n,
  \qquad
  \Phi_n:=\sum_{k=1}^n\varphi_k.
\end{equation}
On the continuous side, since
$\widehat{\mathcal G}(z)=\frac{1}{z(\gamma+z^\alpha)}$,
we have
$z^\alpha\widehat{\mathcal G}(z)=\frac1z-\gamma\widehat{\mathcal G}(z)$,
and hence $\mathcal G$ solves
\begin{equation}\label{eq:G-cont-fde-proof}
  D^\alpha\mathcal G(t)=1-\gamma\mathcal G(t),
  \qquad
  \mathcal G(0)=0.
\end{equation}
Define
\[
  e_n^\tau:=\mathcal G_n^\tau-\mathcal G(t_n),
  \qquad 0\le n\le N_\tau.
\]
Then
\begin{equation}\label{eq:error-renewal-proof}
  e_n^\tau
  =
  a_\tau\sum_{k=1}^n\varphi_k\,e_{n-k}^\tau
  -L_n^\tau,
\end{equation}
where
\begin{equation}\label{eq:local-defect-def-proof}
  L_n^\tau
  :=
  \mathcal G(t_n)
  -
  a_\tau\sum_{k=1}^n\varphi_k\,\mathcal G(t_{n-k})
  -
  \tau^{-\alpha}\Phi_n.
\end{equation}

\medskip
\noindent
\textit{Step 2: defect decomposition.}
Introduce the exact tail sequence
\[
  \bar\varphi_j:=\sum_{m=j+1}^\infty \varphi_m,
  \qquad j\ge0.
\]
For the kernel from Definition~\ref{def:params},
\begin{equation}\label{eq:exact-tail}
  \bar\varphi_0=1,
  \qquad
  \bar\varphi_j=\frac{j^{-\alpha}}{\Gamma(1-\alpha)},
  \qquad j\ge1.
\end{equation}
Since $\varphi_k=\bar\varphi_{k-1}-\bar\varphi_k$ and $\mathcal G(0)=0$,
summation by parts yields
\begin{equation}\label{eq:abel-G-proof}
  \mathcal G(t_n)-\sum_{k=1}^n\varphi_k\,\mathcal G(t_{n-k})
  =
  \sum_{j=0}^{n-1}\bar\varphi_j
  \left(\mathcal G(t_{n-j})-\mathcal G(t_{n-j-1})\right).
\end{equation}
By \eqref{eq:abel-G-proof} and \eqref{eq:exact-tail},
\begin{align}
  &\mathcal G(t_n)-\sum_{k=1}^n\varphi_k\,\mathcal G(t_{n-k})
  \nonumber\\
  &=
  \int_0^1 \frac1\tau \mathcal G'(t_n-v/\tau)\,dv
  +
  \frac1{\Gamma(1-\alpha)}
  \sum_{j=1}^{n-1}
  j^{-\alpha}\int_j^{j+1}\frac1\tau\mathcal G'(t_n-v/\tau)\,dv.
  \label{eq:discrete-caputo-proof}
\end{align}
On the other hand, from \eqref{eq:G-cont-fde-proof},
\begin{equation}\label{eq:continuous-caputo-proof}
  \tau^{-\alpha}\left(1-\gamma\mathcal G(t_n)\right)
  =
  \frac1{\Gamma(1-\alpha)}
  \int_0^n v^{-\alpha}\frac1\tau\mathcal G'(t_n-v/\tau)\,dv.
\end{equation}
Using $a_\tau=1-\gamma\tau^{-\alpha}$ and $\Phi_n=1-\bar\varphi_n$, we may
therefore rewrite \eqref{eq:local-defect-def-proof} as
\begin{equation}\label{eq:Ln-split-proof}
  L_n^\tau
  =
  \mathcal D_n^\tau
  +
  \gamma\tau^{-\alpha}
  \left(
    \sum_{k=1}^n\varphi_k\,\mathcal G(t_{n-k})-\mathcal G(t_n)
  \right)
  +\tau^{-\alpha}\bar\varphi_n,
\end{equation}
where $\mathcal D_n^\tau$ is the quadrature defect from
Lemma~\ref{lem:quadrature-defect-Gprime}.

\medskip
\noindent
\textit{Step 3: bound the main defect.}
By Lemma~\ref{lem:quadrature-defect-Gprime},
\[
  \mathcal D_n^\tau
  =
  c_1(\alpha)\,\frac1\tau \mathcal G'(t_n)
  +R_n^\tau,
\]
with
$|R_n^\tau|\le C_T\tau^{-\alpha}n^{-\alpha}$.
Since
$\mathcal G'(t)=\frac1\gamma f_{\alpha,\gamma}(t)\le C_T t^{\alpha-1}$,
we obtain
\begin{equation}\label{eq:Ln-main-bound}
  |\mathcal D_n^\tau|
  \le
  C_T|c_1(\alpha)|\,\tau^{-\alpha}n^{\alpha-1}
  +C_T\tau^{-\alpha}n^{-\alpha}.
\end{equation}

\medskip
\noindent
\textit{Step 4: bound the lower-order terms.}
By \eqref{eq:abel-G-proof}, \eqref{eq:exact-tail}, and the bound
\[
  |\mathcal G(t_{m})-\mathcal G(t_{m-1})|
  \le
  \int_{t_{m-1}}^{t_m}|\mathcal G'(s)|\,ds
  \le C_T\tau^{-\alpha}m^{\alpha-1},
  \qquad 1\le m\le N_\tau,
\]
we obtain
\begin{align*}
  \left|
    \mathcal G(t_n)-\sum_{k=1}^n\varphi_k\,\mathcal G(t_{n-k})
  \right|
  &\le
  C_T\tau^{-\alpha}
  \left(
    n^{\alpha-1}
    +
    \sum_{j=1}^{n-1}j^{-\alpha}(n-j)^{\alpha-1}
  \right).
\end{align*}
Splitting the discrete convolution at $j=\lfloor n/2\rfloor$ gives
\[
  \sum_{j=1}^{n-1}j^{-\alpha}(n-j)^{\alpha-1}
  \le
  Cn^{\alpha-1}\sum_{j\le n/2}j^{-\alpha}
  +
  Cn^{-\alpha}\sum_{r\le n/2}r^{\alpha-1}
  \le C,
\]
and hence,
\[
  \left|
    \mathcal G(t_n)-\sum_{k=1}^n\varphi_k\,\mathcal G(t_{n-k})
  \right|
  \le C_T\tau^{-\alpha}.
\]
Since $\bar\varphi_n=n^{-\alpha}/\Gamma(1-\alpha)$ for $n\ge1$, it follows that
\begin{equation}\label{eq:Ln-lower-order}
  \left|
    \gamma\tau^{-\alpha}
    \left(
      \sum_{k=1}^n\varphi_k\,\mathcal G(t_{n-k})-\mathcal G(t_n)
    \right)
    +\tau^{-\alpha}\bar\varphi_n
  \right|
  \le
  C_T\left(\tau^{-2\alpha}+\tau^{-\alpha}n^{-\alpha}\right).
\end{equation}
Combining \eqref{eq:Ln-split-proof}, \eqref{eq:Ln-main-bound}, and
\eqref{eq:Ln-lower-order}, we conclude that
\begin{equation}\label{eq:Ln-final-bound}
  |L_n^\tau|
  \le
  C_T\tau^{-\alpha}n^{-\alpha}
  +
  C_T|c_1(\alpha)|\,\tau^{-\alpha}n^{\alpha-1},
  \qquad 1\le n\le N_\tau.
\end{equation}

\medskip
\noindent
\textit{Step 5: renewal amplification.}
From \eqref{eq:error-renewal-proof}, inversion by the renewal kernel gives
\begin{equation}\label{eq:error-inversion-proof}
  e_n^\tau
  =
  -L_n^\tau-\sum_{j=1}^{n-1}\psi_{n-j}^\tau L_j^\tau.
\end{equation}
By Lemma~\ref{lem:psi-pointwise},
$\psi_m^\tau\le C_T m^{\alpha-1}$
for any $1\le m\le N_\tau$.
Hence, using \eqref{eq:Ln-final-bound},
\begin{align*}
  \sum_{j=1}^{n-1}\psi_{n-j}^\tau |L_j^\tau|
  \le
  C_T\tau^{-\alpha}
  \sum_{j=1}^{n-1}(n-j)^{\alpha-1}j^{-\alpha}
  +
  C_T|c_1(\alpha)|\,\tau^{-\alpha}
  \sum_{j=1}^{n-1}(n-j)^{\alpha-1}j^{\alpha-1}.
\end{align*}
The first sum is uniformly bounded in $n$, while the second satisfies
\[
  \sum_{j=1}^{n-1}(n-j)^{\alpha-1}j^{\alpha-1}
  \le C n^{2\alpha-1}.
\]
Since $n\le N_\tau\le T\tau$,
$\tau^{-\alpha}n^{2\alpha-1}\le C_T\tau^{-(1-\alpha)}$.
Therefore,
\[
  \sum_{j=1}^{n-1}\psi_{n-j}^\tau |L_j^\tau|
  \le
  C_T\tau^{-\alpha}
  +
  C_T|c_1(\alpha)|\,\tau^{-(1-\alpha)}.
\]
The same bound follows for $|L_n^\tau|$ itself from \eqref{eq:Ln-final-bound}. Combining this with \eqref{eq:error-inversion-proof}, we obtain
\[
  |e_n^\tau|
  \le
  C_T\tau^{-\alpha}
  +
  C_{\mathrm{res}}(\alpha)\tau^{-(1-\alpha)},
  \qquad 0\le n\le N_\tau,
\]
where
$C_{\mathrm{res}}(\alpha)\to0$
as $\alpha\uparrow1^-$.
Taking the supremum over $0\le n\le N_\tau$ proves
\eqref{eq:cum-resolvent-uniform-comparison}.

%%%%%%%%%%%%%%%%%%%%%%%%%%%%%%%%%%%%%%%%%%%%%%%%%%

\subsection{Proof of Proposition~\ref{prop:resolvent-transfer-from-cumulative}}\label{app:proof-resolvent-transfer-from-cumulative}

\begin{proof}
Define
\[
  (\mathcal Q_\tau g)_n
  :=
  \frac{1}{\gamma\tau}\sum_{j=1}^{n-1}\mathsf{S}_{n-j}^\tau g(t_j),
  \qquad
  (\mathcal Q g)(t_n)
  :=
  \frac1\gamma\int_0^{t_n}f_{\alpha,\gamma}(t_n-s)g(s)\,ds.
\]

Recall from the definition of $\mathcal G_n^\tau$ that
\[
  \mathcal G_n^\tau-\mathcal G_{n-1}^\tau
  =
  \frac{1}{\gamma\tau}\mathsf{S}_n^\tau,
  \qquad n\ge1,
\]
and
\[
  \mathcal G'(t)=\frac1\gamma f_{\alpha,\gamma}(t),\qquad t>0.
\]
In particular,
\[
  |\mathcal G'(t)|\le C_T t^{\alpha-1},
  \qquad 0<t\le T,
\]
and therefore,
\[
  |\mathcal G(t)-\mathcal G(s)|\le C_T|t-s|^\alpha,
  \qquad 0\le s,t\le T.
\]

Using summation by parts, for $1\le n\le N_\tau$,
\begin{equation}\label{eq:Qtau-abel}
  (\mathcal Q_\tau g)_n
  =
  g(0)\mathcal G_{n-1}^\tau
  +\sum_{j=0}^{n-2}\mathcal G_{n-j-1}^\tau
   \int_{t_j}^{t_{j+1}} g'(s)\,ds.
\end{equation}
On the continuous side, integration by parts yields
\begin{equation}\label{eq:Q-abel}
  (\mathcal Q g)(t_n)
  =
  g(0)\mathcal G(t_n)
  +\int_0^{t_n}\mathcal G(t_n-s)g'(s)\,ds.
\end{equation}
Subtracting \eqref{eq:Q-abel} from \eqref{eq:Qtau-abel}, we obtain
\begin{align*}
  &(\mathcal Q_\tau g)_n-(\mathcal Q g)(t_n)
    =
  g(0)\left(\mathcal G_{n-1}^\tau-\mathcal G(t_n)\right)
  \\
  &
  \qquad\qquad+
  \sum_{j=0}^{n-2}
  \int_{t_j}^{t_{j+1}}
  \left(\mathcal G_{n-j-1}^\tau-\mathcal G(t_n-s)\right)g'(s)\,ds
  -
  \int_{t_{n-1}}^{t_n}\mathcal G(t_n-s)g'(s)\,ds.
\end{align*}
Hence,
\begin{align}
  \left|(\mathcal Q_\tau g)_n-(\mathcal Q g)(t_n)\right|
  &\le
  |g(0)|\,|\mathcal G_{n-1}^\tau-\mathcal G(t_n)|
  \nonumber
  \\
  &\qquad
  +
  \sum_{j=0}^{n-2}
  \int_{t_j}^{t_{j+1}}
  \left|\mathcal G_{n-j-1}^\tau-\mathcal G(t_n-s)\right|
  |g'(s)|\,ds
  \nonumber
  \\
  &\qquad\qquad
  +
  \int_{t_{n-1}}^{t_n}|\mathcal G(t_n-s)|\,|g'(s)|\,ds.\label{3:terms:2}
\end{align}
For the first term on the right hand side in \eqref{3:terms:2}, we insert and subtract $\mathcal G(t_{n-1})$:
\[
  |\mathcal G_{n-1}^\tau-\mathcal G(t_n)|
  \le
  |\mathcal G_{n-1}^\tau-\mathcal G(t_{n-1})|
  +
  |\mathcal G(t_{n-1})-\mathcal G(t_n)|.
\]
By Proposition~\ref{prop:cum-resolvent-uniform-comparison} and the
$\alpha$-H\"older continuity of $\mathcal G$, this gives
\[
  |\mathcal G_{n-1}^\tau-\mathcal G(t_n)|
  \le
  C_T\left(\tau^{-\alpha}+{C_{\mathrm{res}}(\alpha)}\,\tau^{-(1-\alpha)}\right).
\]

We split
\[
  \left|\mathcal G_{n-j-1}^\tau-\mathcal G(t_n-s)\right|
  \le
  \left|\mathcal G_{n-j-1}^\tau-\mathcal G(t_{n-j-1})\right|
  +
  \left|\mathcal G(t_{n-j-1})-\mathcal G(t_n-s)\right|.
\]
By Proposition~\ref{prop:cum-resolvent-uniform-comparison}, the first term in \eqref{3:terms:2} is
bounded by
\[
  C_T\left(\tau^{-\alpha}+{C_{\mathrm{res}}(\alpha)}\,\tau^{-(1-\alpha)}\right).
\]
Since $\mathcal G$ is $\alpha$-H\"older continuous on $[0,T]$, the second term on the right hand side in \eqref{3:terms:2} is bounded
by $C_T\tau^{-\alpha}$ uniformly for $s\in[t_j,t_{j+1}]$. Therefore,
\[
  \left|\mathcal G_{n-j-1}^\tau-\mathcal G(t_n-s)\right|
  \le
  C_T\left(\tau^{-\alpha}+{C_{\mathrm{res}}(\alpha)}\,\tau^{-(1-\alpha)}\right).
\]
Also $\mathcal G$ is bounded on $[0,T]$, and thus the last integral in \eqref{3:terms:2} is bounded by
\[
  C_T\int_{t_{n-1}}^{t_n}|g'(s)|\,ds.
\]
Using \eqref{eq:g-assump-2}, we obtain
\[
  \int_{t_{n-1}}^{t_n}|g'(s)|\,ds
  \le
  C_g\int_{t_{n-1}}^{t_n}s^{\alpha-1}\,ds.
\]
If $n=1$, then
\[
  \int_{0}^{1/\tau}s^{\alpha-1}\,ds
  =\frac1\alpha\tau^{-\alpha}.
\]
If $n\ge2$, then $s\ge t_1=1/\tau$ on $[t_{n-1},t_n]$, so
\[
  \int_{t_{n-1}}^{t_n}s^{\alpha-1}\,ds
  \le
  \tau^{-1}(t_{n-1})^{\alpha-1}
  \le C_T\tau^{-1}
  \le C_T\tau^{-\alpha},
\]
since $\alpha\in(0,1)$. Therefore, uniformly for $1\le n\le N_\tau$,
\[
  \int_{t_{n-1}}^{t_n}|g'(s)|\,ds
  \le C_T C_g\,\tau^{-\alpha}.
\]
Combining the preceding bounds gives
\[
  \left|(\mathcal Q_\tau g)_n-(\mathcal Q g)(t_n)\right|
  \le
  C_T' \left(|g(0)|+C_g\right)\tau^{-\alpha}
  +
  C_{\mathrm{op}}(\alpha)\left(|g(0)|+C_g\right)\tau^{-(1-\alpha)},
\]
uniformly in $1\le n\le N_\tau$, where
$C_{\mathrm{op}}(\alpha)\to0$
as $\alpha\uparrow1^-$.
This proves \eqref{eq:resolvent-operator-compare-main}.
\end{proof}
%%%%%%%%%%%%%%%%%%%%%%%%%%%%%%%%%%%%%%%%%%%%%%%%%%

\subsection{Lemma~\ref{lem:resolvent-gz-compare}}\label{app:proof-resolvent-gz-compare}

\begin{lemma}[Resolvent comparison for the continuous Riccati nonlinearity]
\label{lem:resolvent-gz-compare}
For every fixed $\eta\in\R$ and $R>0$, there exists a coefficient
$C_g(\alpha)\ge0$ and a constant $C_{\eta,R,T}>0$ such that
\begin{equation}\label{eq:resolvent-gz-compare-main}
\begin{aligned}
  &\sup_{1\le n\le N_\tau}\sup_{z\in\Gamma_{\eta,R}}
  \left|
    \frac{1}{\gamma\tau}\sum_{j=1}^{n-1}\mathsf S_{n-j}^\tau F_z(h(t_j,z))
    -
    \frac1\gamma\int_0^{t_n}f_{\alpha,\gamma}(t_n-s)F_z(h(s,z))\,ds
  \right|
  \\
  &\le
  C_{\eta,R,T}\tau^{-\alpha}+C_g(\alpha)\tau^{-(1-\alpha)},
\end{aligned}
\end{equation}
with $C_g(\alpha)\to0$ as $\alpha\uparrow1^-$.
\end{lemma}

\begin{proof}
By Lemma~\ref{lem:h-time-regularity}, the functions
\[
  g_z(t):=F_z(h(t,z))
\]
satisfy uniformly for $z\in\Gamma_{\eta,R}$ the bounds
\[
  |g_z(t)-g_z(0)|\le C_{\eta,R,T}t^\alpha,
  \qquad
  |\partial_t g_z(t)|\le C_{\eta,R,T}t^{\alpha-1}.
\]
Moreover,
$g_z(0)=\frac12(z^2-z)$
is uniformly bounded on $\Gamma_{\eta,R}$. Applying
Proposition~\ref{prop:resolvent-transfer-from-cumulative} with $g=g_z$ yields
\begin{align*}
  &\sup_{1\le n\le N_\tau}\sup_{z\in\Gamma_{\eta,R}}
  \left|
    \frac{1}{\gamma\tau}\sum_{j=1}^{n-1}\mathsf{S}_{n-j}^\tau F_z(h(t_j,z))
    -
    \frac1\gamma\int_0^{t_n}f_{\alpha,\gamma}(t_n-s)F_z(h(s,z))\,ds
  \right|\\
  &\le
  C_{\eta,R,T}\tau^{-\alpha}+C_{g}(\alpha)\tau^{-(1-\alpha)},
\end{align*}
where
$C_g(\alpha)\to0$
as $\alpha\uparrow1^-$.
This proves \eqref{eq:resolvent-gz-compare-main}.
\end{proof}
%%%%%%%%%%%%%%%%%%%%%%%%%%%%%%%%%%%%%%%%%%%%%%%%%%

\subsection{Lemma~\ref{lem:discrete-weakly-singular-gronwall}}\label{app:proof-discrete-weakly-singular-gronwall}

\begin{lemma}[Discrete weakly singular Gr\"onwall on the $\tau$-grid]
\label{lem:discrete-weakly-singular-gronwall}
Fix $\alpha\in(0,1)$, $T>0$, and $B>0$. Suppose that a nonnegative sequence
$(a_n)_{n=0}^{N_\tau}$ satisfies
\begin{equation}\label{eq:discrete-weakly-singular-gronwall-assump}
  a_n\le A+B\tau^{-\alpha}\sum_{j=1}^{n-1}(n-j)^{\alpha-1}a_j,
  \qquad 1\le n\le N_\tau,
\end{equation}
for some $A\ge0$. Then
\begin{equation}\label{eq:discrete-weakly-singular-gronwall-concl}
  a_n\le A E_\alpha(B\Gamma(\alpha)t_n^\alpha),
  \qquad 1\le n\le N_\tau.
\end{equation}
In particular,
\[
  \sup_{0\le n\le N_\tau}a_n\le C_{\alpha,B,T}A.
\]
\end{lemma}

\begin{proof}
Define iteratively
\[
  b_n^{(0)}:=A,
  \qquad
  b_n^{(m+1)}
  :=
  A+B\tau^{-\alpha}\sum_{j=1}^{n-1}(n-j)^{\alpha-1}b_j^{(m)},
  \qquad m\ge0.
\]
Since the kernel is nonnegative,
\eqref{eq:discrete-weakly-singular-gronwall-assump} implies by induction on
$m$ that
\[
  a_n\le b_n^{(m)},
  \qquad 0\le n\le N_\tau,\quad m\ge0.
\]
We claim that, for every $m\ge0$,
\begin{equation}\label{eq:discrete-weakly-singular-gronwall-iterates}
  b_n^{(m)}
  \le
  A\sum_{r=0}^{m}
  \frac{\left(B\Gamma(\alpha)t_n^\alpha\right)^r}{\Gamma(r\alpha+1)},
  \qquad 0\le n\le N_\tau.
\end{equation}
The case $m=0$ is immediate. Assume the claim holds at rank $m$. Since
$\tau^{-\alpha}(n-j)^{\alpha-1}
  =
  \tau^{-1}(t_n-t_j)^{\alpha-1}$,
and since $\alpha-1<0$ while $r\alpha\ge0$, for every
$s\in[t_j,t_{j+1}]$, we have
\[
  (t_n-t_j)^{\alpha-1}t_j^{r\alpha}
  \le
  (t_n-s)^{\alpha-1}s^{r\alpha}.
\]
Since $\alpha-1<0$, the map $u\mapsto (t_n-u)^{\alpha-1}$ is increasing on
$[0,t_n)$, while $u\mapsto u^{r\alpha}$ is also increasing. Hence, for
$s\in[t_j,t_{j+1}]$,
\begin{align}
  (t_n-t_j)^{\alpha-1}\le (t_n-s)^{\alpha-1},
\qquad\text{and}\qquad
  t_j^{r\alpha}\le s^{r\alpha},
  \label{eq:gronwall-pointwise}
\end{align}
and multiplying the two inequalities in \eqref{eq:gronwall-pointwise} yields
\[
  (t_n-t_j)^{\alpha-1}t_j^{r\alpha}
  \le
  (t_n-s)^{\alpha-1}s^{r\alpha}.
\]
Therefore,
\begin{align*}
  \tau^{-\alpha}\sum_{j=1}^{n-1}(n-j)^{\alpha-1}t_j^{r\alpha}
  &=
  \tau^{-1}\sum_{j=1}^{n-1}(t_n-t_j)^{\alpha-1}t_j^{r\alpha}
  \\
  &\le
  \sum_{j=1}^{n-1}\int_{t_j}^{t_{j+1}}(t_n-s)^{\alpha-1}s^{r\alpha}\,ds
  \\
  &\le
  \int_0^{t_n}(t_n-s)^{\alpha-1}s^{r\alpha}\,ds
 =
  \frac{\Gamma(\alpha)\Gamma(r\alpha+1)}{\Gamma((r+1)\alpha+1)}
  t_n^{(r+1)\alpha}.
\end{align*}
Using the induction hypothesis in the definition of $b_n^{(m+1)}$ yields
\[
  b_n^{(m+1)}
  \le
  A\sum_{r=0}^{m+1}
  \frac{\left(B\Gamma(\alpha)t_n^\alpha\right)^r}{\Gamma(r\alpha+1)},
\]
which proves \eqref{eq:discrete-weakly-singular-gronwall-iterates}. For each
fixed $n$, the sequence $m\mapsto b_n^{(m)}$ is nondecreasing because the
kernel in \eqref{eq:discrete-weakly-singular-gronwall-assump} is nonnegative.
Moreover, \eqref{eq:discrete-weakly-singular-gronwall-iterates} provides a
uniform upper bound independent of $m$, so $b_n^{(m)}$ converges as
$m\to\infty$. Since $a_n\le b_n^{(m)}$ for every $m$, passing to the limit and
using the monotone convergence of the truncated Mittag--Leffler series gives
\eqref{eq:discrete-weakly-singular-gronwall-concl}. The final uniform bound
follows since $t_n\le T$.
\end{proof}

%%%%%%%%%%%%%%%%%%%%%%%%%%%%%%%%%%%%%%%%%%%%%%%%%%

\subsection{Proof of Proposition~\ref{prop:Volterra-rate-fixed}}\label{app:proof-Volterra-rate-fixed}
Recall from Lemma~\ref{lem:renewal-inversion} that
\begin{equation}\label{eq:Htau-resolvent-discrete-newproof}
\begin{aligned}
  H_n^\tau(z)
  &=
  \tau^{-\alpha}F_z(H_n^\tau(z))
  +\widehat E_n^\tau(z)
\\
  &\quad
  +
  \frac{a_\tau}{1-a_\tau}\tau^{-\alpha-1}
  \sum_{j=1}^{n-1}\mathsf{S}_{n-j}^\tau F_z(H_j^\tau(z))
  +
  \sum_{j=1}^{n-1}\psi_{n-j}^\tau\widehat E_j^\tau(z),
\end{aligned}
\end{equation}
where
\[
  \sup_{0\le n\le N_\tau}\sup_{z\in\Gamma_{\eta,R}}
  |\widehat E_n^\tau(z)|
  \le C_{\eta,R,T}\tau^{-2\alpha}.
\]

On the continuous side, Lemma~\ref{lem:continuous-resolvent} gives
\[
  h(t_n,z)
  =
  \frac1\gamma\int_0^{t_n}
    f_{\alpha,\gamma}(t_n-s)F_z(h(s,z))\,ds.
\]

Set
\[
  e_n^\tau(z):=H_n^\tau(z)-h(t_n,z).
\]
Subtracting the two equations in \eqref{eq:Htau-resolvent-discrete-newproof} and \eqref{eq:continuous-resolvent}, and inserting the intermediate discrete resolvent
operator applied to $F_z(h(\cdot,z))$, we obtain
\begin{align*}
  |e_n^\tau(z)|
  &\le
  \tau^{-\alpha}|F_z(H_n^\tau(z))|
  +|\widehat E_n^\tau(z)|
  +\sum_{j=1}^{n-1}\psi_{n-j}^\tau |\widehat E_j^\tau(z)|
  \\
  &\quad
  +
  \left|
    \frac{a_\tau}{1-a_\tau}\tau^{-\alpha-1}
    \sum_{j=1}^{n-1}\mathsf{S}_{n-j}^\tau F_z(H_j^\tau(z))
    -
    \frac{1}{\gamma\tau}\sum_{j=1}^{n-1}\mathsf{S}_{n-j}^\tau F_z(h(t_j,z))
  \right|
  \\
  &\quad\quad
  +
  \left|
    \frac{1}{\gamma\tau}\sum_{j=1}^{n-1}\mathsf{S}_{n-j}^\tau F_z(h(t_j,z))
    -
    \frac1\gamma\int_0^{t_n}
      f_{\alpha,\gamma}(t_n-s)F_z(h(s,z))\,ds
  \right|.
\end{align*}

The last term is bounded by
\[
  C_{\eta,R,T}
  \left(\tau^{-\alpha}+{C_g(\alpha)}\,\tau^{-(1-\alpha)}\right)
\]
by Lemma~\ref{lem:resolvent-gz-compare}. The explicit local term
$\tau^{-\alpha}|F_z(H_n^\tau(z))|$ is of order $\tau^{-\alpha}$ because
$H_n^\tau(z)$ remains uniformly bounded on bounded strips, as follows from
Proposition~\ref{prop:smallness-new33} and the rescaling
\eqref{eq:H-rescaled-def}. The same boundedness implies that $F_z$ is
uniformly Lipschitz on the relevant range:
\[
  |F_z(x)-F_z(y)|\le C_{\eta,R,T}|x-y|.
\]

Using
$\frac{a_\tau}{1-a_\tau}\tau^{-\alpha-1}
  =
  \frac{1}{\gamma\tau}
  +\mathcal O(\tau^{-1-\alpha})$,
together with
$\sum_{m\ge1}\mathsf{S}_m^\tau=\tau$,
we absorb the coefficient mismatch into another $C_{\eta,R,T}\tau^{-\alpha}$
term and obtain
\[
  |e_n^\tau(z)|
  \le
  C_{\eta,R,T}
  \left(\tau^{-\alpha}+{C_g(\alpha)}\,\tau^{-(1-\alpha)}\right)
  +
  \frac{C_{\eta,R,T}}{\tau}
  \sum_{j=1}^{n-1}\mathsf{S}_{n-j}^\tau |e_j^\tau(z)|
  +
  \sum_{j=1}^{n-1}\psi_{n-j}^\tau |\widehat E_j^\tau(z)|.
\]

Since
\[
  \sup_{1\leq j\leq n-1} |\widehat E_j^\tau(z)|\le C_{\eta,R,T}\tau^{-2\alpha},
  \qquad
  \sum_{m\ge1}\psi_m^\tau=\frac{a_\tau}{1-a_\tau}\lesssim \tau^\alpha,
\]
we have
\[
  \sum_{j=1}^{n-1}\psi_{n-j}^\tau |\widehat E_j^\tau(z)|
  \le
  C_{\eta,R,T}\tau^{-2\alpha}\sum_{m\ge1}\psi_m^\tau
  \le
  C_{\eta,R,T}\tau^{-\alpha}.
\]
Hence
\[
  |e_n^\tau(z)|
  \le
  C_{\eta,R,T}
  \left(\tau^{-\alpha}+{C_g(\alpha)}\,\tau^{-(1-\alpha)}\right)
  +
  \frac{C_{\eta,R,T}}{\tau}
  \sum_{j=1}^{n-1}\mathsf{S}_{n-j}^\tau |e_j^\tau(z)|.
\]

Finally, using the pointwise bound from Lemma~\ref{lem:psi-pointwise},
\[
  \frac{1}{\tau}\mathsf{S}_{n-j}^\tau
  =
  \frac{1-a_\tau}{a_\tau}\psi_{n-j}^\tau
  \le
  C_T\tau^{-\alpha}(n-j)^{\alpha-1}
  =
  \frac{C_T}{\tau^\alpha (n-j)^{1-\alpha}},
\]
and setting
\[
  a_n:=\sup_{0\le m\le n}\sup_{z\in\Gamma_{\eta,R}}|e_m^\tau(z)|,
\]
we obtain from the preceding inequality that
\[
  a_n
  \le
  C_{\eta,R,T}\tau^{-\alpha}
  +
  C_{\mathrm{Ric}}(\alpha)\tau^{-(1-\alpha)}
  +
  C_T\tau^{-\alpha}\sum_{j=1}^{n-1}(n-j)^{\alpha-1}a_j,
\]
where
$C_{\mathrm{Ric}}(\alpha)\to0$
as $\alpha\uparrow1^-$.
Applying Lemma~\ref{lem:discrete-weakly-singular-gronwall} with
\[
  A=C_{\eta,R,T}\tau^{-\alpha}
  +
  C_{\mathrm{Ric}}(\alpha)\tau^{-(1-\alpha)}
\]
and $B=C_T$, we obtain
\[
  \sup_{0\le n\le N_\tau}\sup_{z\in\Gamma_{\eta,R}}|e_n^\tau(z)|
  \le
  C_{\eta,R,T}\tau^{-\alpha}
  +
  C_{\mathrm{Ric}}(\alpha)\tau^{-(1-\alpha)}.
\]
This proves \eqref{eq:Volterra-rate-fixed}.

%%%%%%%%%%%%%%%%%%%%%%%%%%%%%%%%%%%%%%%%%%%%%%%%%%
\subsection{Proof of Lemma~\ref{lem:exponent-comparison}}\label{app:proof-exponent-comparison}

\begin{proof}
By Proposition~\ref{prop:Volterra-rate-fixed},
\[
  \sup_{0\le n\le N_\tau}\sup_{z\in\Gamma_{\eta,R}}
  \left|H_n^\tau(z)-h(t_n,z)\right|
  \le
  C_{\eta,R,T}\tau^{-\alpha}+{C_{\mathrm{Ric}}(\alpha)}\,\tau^{-(1-\alpha)},
\]
where $t_n:=n/\tau$.
Using
\[
  \log\Phi^\tau(z,T)
  =
  \frac{\theta}{\mu}(1-a_\tau)
  \sum_{m=1}^{N_\tau}\hat\mu_\tau(m)\,H_{N_\tau-m}^\tau(z),
\]
we write
\begin{align}
  &\left|
    \log\Phi^\tau(z,T)-\log\phi(z,T)
  \right|
  \nonumber
  \\
  &\le
  \frac{\theta}{\mu}(1-a_\tau)
  \sum_{m=1}^{N_\tau}\hat\mu_\tau(m)
  \left|
    H_{N_\tau-m}^\tau(z)-h\!\left(\frac{N_\tau-m}{\tau},z\right)
  \right|
  \nonumber
  \\
  &\qquad
  +
  \frac{\theta}{\mu}(1-a_\tau)
  \sum_{m=1}^{N_\tau}\hat\mu_\tau(m)
  \left|
    h\!\left(\frac{N_\tau-m}{\tau},z\right)-h\!\left(T-\frac{m}{\tau},z\right)
  \right|
  \nonumber
  \\
  &\qquad\qquad
  +
  \left|
    \frac{\theta}{\mu}(1-a_\tau)
    \sum_{m=1}^{N_\tau}\hat\mu_\tau(m)
    h\!\left(T-\frac{m}{\tau},z\right)
    -
    \log\phi(z,T)
  \right|.\label{3:terms}
\end{align}

For the first term in \eqref{3:terms}, Proposition~\ref{prop:Volterra-rate-fixed} gives
\[
  \left|
    H_{N_\tau-m}^\tau(z)-h\!\left(\frac{N_\tau-m}{\tau},z\right)
  \right|
  \le
  C_{\eta,R,T}\tau^{-\alpha}+{C_{\mathrm{Ric}}(\alpha)}\,\tau^{-(1-\alpha)}
\]
uniformly in $m$ and $z$. Moreover, Proposition~\ref{prop:baseline-riemann},
applied to the constant function $f\equiv1$, yields
\begin{equation}\label{use:again}
  \frac{1-a_\tau}{\mu}\sum_{m=1}^{N_\tau}\hat\mu_\tau(m)\le C_T.
\end{equation}
Hence the first term in \eqref{3:terms} is bounded by
\[
  C_{\eta,R,T}\tau^{-\alpha}+{C_{\mathrm{Ric}}(\alpha)}\,\tau^{-(1-\alpha)}.
\]

For the second term in \eqref{3:terms}, note that
\[
  \left|
    \frac{N_\tau-m}{\tau}-\left(T-\frac{m}{\tau}\right)
  \right|
  =
  \left|\frac{N_\tau}{\tau}-T\right|
  \le \frac1\tau.
\]
Since $h(\cdot,z)$ is uniformly $\alpha$-H\"older continuous on $[0,T]$ by
Proposition~\ref{prop:riccati-regularity}, we obtain
\[
  \left|
    h\!\left(\frac{N_\tau-m}{\tau},z\right)-h\!\left(T-\frac{m}{\tau},z\right)
  \right|
  \le C_{\eta,R,T}\tau^{-\alpha}
\]
uniformly in $m$ and $z$. Using again \eqref{use:again},
%\[
%  \frac{1-a_\tau}{\mu}\sum_{m=1}^{N_\tau}\hat\mu_\tau(m)\le C_T,
%\]
it follows that the second term in \eqref{3:terms} is bounded by
$C_{\eta,R,T}\tau^{-\alpha}$.

For the third term in \eqref{3:terms}, we apply Proposition~\ref{prop:baseline-riemann} with
$f(s)=h(s,z)$. By Proposition~\ref{prop:riccati-regularity}, the function $h(\cdot,z)$
is uniformly $\alpha$-H\"older continuous on $[0,T]$, and
\[
  |h(t,z)|\le C_{\eta,R,T}t^\alpha,
\]
uniformly for $z\in\Gamma_{\eta,R}$. Therefore,
\begin{align*}
  \left|
    \frac{\theta}{\mu}(1-a_\tau)
    \sum_{m=1}^{N_\tau}\hat\mu_\tau(m)
    h\!\left(T-\frac{m}{\tau},z\right)
    -
    \log\phi(z,T)
  \right|
  \le
  C_{\eta,R,T}\tau^{-\alpha}
  +
  C_{\mathrm{base}}(\alpha)\tau^{-(1-\alpha)},
\end{align*}
uniformly on $\Gamma_{\eta,R}$, where $C_{\mathrm{base}}(\alpha)\to0$ as $\alpha\uparrow1^-$.

Combining the bounds on the three terms in \eqref{3:terms}, and enlarging
$C_{\eta,R,T}$ if necessary, we obtain
\[
  \sup_{z\in\Gamma_{\eta,R}}
  \left|
    \log\Phi^\tau(z,T)-\log\phi(z,T)
  \right|
  \le
  C_{\eta,R,T}\tau^{-\alpha}
  +
  C_{\log}(\alpha)\tau^{-(1-\alpha)},
\]
where
$C_{\log}(\alpha):=C_{\mathrm{Ric}}(\alpha)+C_{\mathrm{base}}(\alpha)$,
so that
$C_{\log}(\alpha)\to0$
as $\alpha\uparrow1^-$.
This is exactly \eqref{eq:log-transform-bound}.
\end{proof}

%%%%%%%%%%%%%%%%%%%%%%%%%%%%%%%%%%%%%%%%%%%%%%%%%%

\subsection{Proof of Proposition~\ref{prop:transform-bound-fixed}}\label{app:proof-transform-bound-fixed}

\begin{proof}
For $z=\eta+\ii\xi\in\Gamma_{\eta,R}$, the triangle inequality gives
\[
  |\Phi^\tau(z,T)|
  =
  \left|\E\!\left[e^{zP_T^\tau}\right]\right|
  \le
  \E\!\left[e^{\eta P_T^\tau}\right]
  =
  \Phi^\tau(\eta,T),
\]
and similarly
\[
  |\phi(z,T)|
  =
  \left|\E\!\left[e^{zP_T}\right]\right|
  \le
  \E\!\left[e^{\eta P_T}\right]
  =
  \phi(\eta,T).
\]
Since $\eta\in\mathcal S_T$, both $\Phi^\tau(\eta,T)$ and $\phi(\eta,T)$ are finite,
and therefore the transforms are uniformly bounded on $\Gamma_{\eta,R}$ for all
sufficiently large $\tau$.

Now set
\[
  x:=\log\Phi^\tau(z,T),
  \qquad
  y:=\log\phi(z,T).
\]
By Lemma~\ref{lem:exponent-comparison},
\[
  |x-y|
  \le
  C_{\eta,R,T}\tau^{-\alpha}
  +
  C_{\log}(\alpha)\tau^{-(1-\alpha)},
\]
uniformly on $\Gamma_{\eta,R}$, with $C_{\log}(\alpha)\to0$ as
$\alpha\uparrow1^-$. Using
\[
  e^x-e^y=(x-y)\int_0^1 e^{tx+(1-t)y}\,dt,
\]
we obtain
\begin{align*}
  |\Phi^\tau(z,T)-\phi(z,T)|
  &=
  |e^x-e^y|
  \\
  &\le
  |x-y|
  \int_0^1 \left|e^{tx+(1-t)y}\right|\,dt.
\end{align*}
Moreover,
\[
  \left|e^{tx+(1-t)y}\right|
  \le
  e^{t\mathrm{Re}\ x+(1-t)\mathrm{Re}\ y}
  \le
  \max\{|\Phi^\tau(z,T)|,|\phi(z,T)|\}
  \le C_{\eta,R,T},
\]
uniformly for $z\in\Gamma_{\eta,R}$ and all sufficiently large $\tau$.
Therefore,
\[
  |\Phi^\tau(z,T)-\phi(z,T)|
  \le
  C_{\eta,R,T}\tau^{-\alpha}
  +
  C_{\Phi}(\alpha)\tau^{-(1-\alpha)},
\]
uniformly on $\Gamma_{\eta,R}$, where one may take
\[
  C_{\Phi}(\alpha):=C_{\eta,R,T}\,C_{\log}(\alpha),
\]
so that $C_{\Phi}(\alpha)\to0$ as $\alpha\uparrow1^-$.
This proves \eqref{eq:transform-bound-fixed}. The equivalent form
\eqref{eq:transform-bound-xi} is just the parametrization $z=\eta+\ii\xi$.
\end{proof}
%%%%%%%%%%%%%%%%%%%%%%%%%%%%%%%%%%%%%%%%%%%%%%%%%%

\subsection{Proof of Proposition~\ref{prop:pricing-local-trunc}}\label{app:proof-pricing-local-trunc}
Subtracting \eqref{eq:CRtau-def-trunc} and \eqref{eq:CR-def-trunc}, we obtain
\begin{align*}
  \frac{C_R^\tau(T,K)-C_R(T,K)}{S_0}
  =
  \frac{e^{-(\eta-1)k}}{2\pi}
  \int_{-R}^{R}
    e^{-\ii \xi k}
    \frac{
      \Phi^\tau(\eta+\ii \xi,T)-\phi(\eta+\ii \xi,T)
    }
    {(\eta+\ii \xi)(\eta+\ii \xi-1)}
  \,d\xi.
\end{align*}
Hence,
\begin{align*}
  \frac{|C_R^\tau(T,K)-C_R(T,K)|}{S_0}
  &\le
  \frac{e^{-(\eta-1)k}}{2\pi}
  \int_{-R}^{R}
    \frac{
      |\Phi^\tau(\eta+\ii \xi,T)-\phi(\eta+\ii \xi,T)|
    }
    {|(\eta+\ii \xi)(\eta+\ii \xi-1)|}
  \,d\xi.
\end{align*}
By Proposition~\ref{prop:transform-bound-fixed},
\[
  \sup_{|\xi|\le R}
  \left|
    \Phi^\tau(\eta+\ii \xi,T)-\phi(\eta+\ii \xi,T)
  \right|
  \le
  C_{\eta,R,T}\tau^{-\alpha}+C_2(\alpha)\tau^{-(1-\alpha)},
\]
where
$C_2(\alpha)\to0$
as $\alpha\uparrow1^-$.
Therefore,
\begin{align*}
  &|C_R^\tau(T,K)-C_R(T,K)|\\
  &\le
  \frac{S_0 e^{-(\eta-1)k}}{2\pi}
  \left(
    C_{\eta,R,T}\tau^{-\alpha}+C_2(\alpha)\tau^{-(1-\alpha)}
  \right)
  \int_{-R}^{R}\frac{d\xi}{|(\eta+\ii \xi)(\eta+\ii \xi-1)|}.
\end{align*}
Since $\eta>1$, the denominator does not vanish on the contour
$\{\eta+\ii\xi:\ |\xi|\le R\}$, and the integral is finite for each fixed
$R>0$. Absorbing the prefactor and the integral into the constants proves
Proposition~\ref{prop:pricing-local-trunc}. 

%%%%%%%%%%%%%%%%%%%%%%%%%%%%%%%%%%%%%%%%%%%%%%%%%%

\subsection{Corollary~\ref{cor:pricing-local-trunc-uniform}}\label{app:proof-pricing-local-trunc-uniform}

\begin{corollary}[Uniform truncated pricing error on compact moneyness sets]
\label{cor:pricing-local-trunc-uniform}
Fix $\eta>1$, $R>0$, $T>0$, and let $\mathcal K\subset\R$ be compact. Then,
for all sufficiently large $\tau$,
\begin{equation}\label{eq:pricing-local-trunc-uniform}
  \sup_{k\in\mathcal K}
  \left|C_R^\tau(T,S_0e^k)-C_R(T,S_0e^k)\right|
  \le
  C_1\tau^{-\alpha}+C_2(\alpha)\tau^{-(1-\alpha)},
\end{equation}
for some constants $C_1>0$ and $C_2(\alpha)\ge0$, independent of $\tau$, with
$C_2(\alpha)\to0$ as $\alpha\uparrow1^-$.
\end{corollary}

\begin{proof}
The estimate \eqref{eq:pricing-local-trunc-uniform} is uniform in $k$ except for the
prefactor $e^{-(\eta-1)k}$. Since $k\in\mathcal K$ and $\mathcal K$ is compact, 
\[
  \sup_{k\in\mathcal K}e^{-(\eta-1)k}
  =
  e^{-(\eta-1)\inf\mathcal K},
\]
where $\inf\mathcal K$ denotes the smallest element of the compact set $\mathcal K$, i.e.\ $\inf\mathcal K=\min_{k\in\mathcal K}k$.
Taking the supremum over $k\in\mathcal K$ in the estimate from
Proposition~\ref{prop:pricing-local-trunc} yields \eqref{eq:pricing-local-trunc-uniform}.
\end{proof}





\end{document}
